\documentclass[mat2, tisk]{fmfdelo}
% \documentclass[fin2, tisk]{fmfdelo}
% \documentclass[isrm2, tisk]{fmfdelo}
% \documentclass[ped, tisk]{fmfdelo}
% Če pobrišete možnost tisk, bodo povezave obarvane,
% na začetku pa ne bo praznih strani po naslovu, …

%%%%%%%%%%%%%%%%%%%%%%%%%%%%%%%%%%%%%%%%%%%%%%%%%%%%%%%%%%%%%%%%%%%%%%%%%%%%%%%
% METAPODATKI
%%%%%%%%%%%%%%%%%%%%%%%%%%%%%%%%%%%%%%%%%%%%%%%%%%%%%%%%%%%%%%%%%%%%%%%%%%%%%%%

% - vaše ime
\avtor{Gal Anton Gorše}

% - naslov dela v slovenščini
\naslov{Povsem pozitivne preslikave in robne upodobitve}

% - naslov dela v angleščini
\title{Completely positive maps and boundary representations}

% - ime mentorja/mentorice s polnim nazivom:
%   - doc.~dr.~Ime Priimek
%   - izr.~prof.~dr.~Ime Priimek
%   - prof.~dr.~Ime Priimek
%   za druge variante uporabite ustrezne ukaze
\mentor{prof.~dr.~Igor Klep}
% \somentor{...}
% \mentorica{...}
% \somentorica{...}
% \mentorja{...}{...}
% \somentorja{...}{...}
% \mentorici{...}{...}
% \somentorici{...}{...}

% - leto magisterija
\letnica{2026}

% - povzetek v slovenščini
%   V povzetku na kratko opišite vsebinske rezultate dela. Sem ne sodi razlaga
%   organizacije dela, torej v katerem razdelku je kaj, pač pa le opis vsebine.
\povzetek{V nekomutativni geometriji geometrijske pojme, ki jih opišemo z algebrami funkcij, pogosto posplošimo na nekomutativne $C^*$-algebre.
V tem magistrskem delu predstavimo Arvesonove robne upodobitve, ki so nekomutativen analog Choquetovega robu za funkcijske sisteme. 
Osrednja tema je uporaba povsem pozitivnih preslikav,
ki jih lahko razumemo kot posplošitev pozitivnih funkcionalov na $C^*$-algebrah. 
Spoznamo ključne izreke o povsem pozitivnih preslikavah,
kot na primer Stinespringov in Arvesonov izrek, ter s pomočjo
maksimalnih povsem pozitivnih preslikav dokažemo obstoj $C^*$-ogrinjače in ideala Šilova za poljuben operatorski sistem. 
Ta pristop vodi do izreka Davidsona in Kennedyja, po katerem robne upodobitve vsak operatorski sistem popolnoma normirajo, s čimer sta avtorja rešila 
46-letno domnevo v teoriji operatorskih algeber.
}

% - povzetek v angleščini
\abstract{In noncommutative geometry, geometric notions that can be described using algebras of functions are often generalized to noncommutative $C^*$-algebras.
In this master's thesis, we present Arveson's boundary representations as a noncommutative analogue of the Choquet boundary for function systems.
The central theme is the use of completely positive maps, which can be viewed as a generalization of positive functionals on $C^*$-algebras.
We discuss key theorems on completely positive maps, including the Stinespring and Arveson theorems, and use maximal completely positive maps to prove the existence of the $C^*$-envelope 
and the Šilov ideal for any operator system.
This approach leads to the Davidson-Kennedy theorem, which states that boundary representations completely norm every operator system, 
thereby resolving a 46-year-old conjecture in the theory of operator algebras.}

% - klasifikacijske oznake, ločene z vejicami
%   Oznake, ki opisujejo področje dela, so dostopne na strani https://www.ams.org/msc/
\klasifikacija{47L25, 46L52, 47A20}

% - ključne besede, ki nastopajo v delu, ločene s \sep
\kljucnebesede{\texorpdfstring{$C^*$}{C*}-algebra\sep povsem pozitivna preslikava\sep operatorski sistem\sep 
operatorski prostor\sep povsem omejena preslikava\sep dilatacija\sep abstrakten operatorski sistem\sep \texorpdfstring{$C^*$}{C*}-ogrinjača\sep rob Šilova\sep ideal Šilova\sep
Choquetov rob\sep robna upodobitev}

% - angleški prevod ključnih besed
\keywords{\texorpdfstring{$C^*$}{C*}-algebra\sep completely positive map\sep operator system\sep
operator space\sep completely bounded map\sep dilation\sep abstract operator system\sep \texorpdfstring{$C^*$}{C*}-envelope\sep Šilov boundary\sep Šilov ideal\sep
Choquet boundary\sep boundary representation}

% - neobvezna zahvala

% - ime datoteke z viri (vključno s končnico .bib), če uporabljate BibTeX
\literatura{literatura.bib}

%%%%%%%%%%%%%%%%%%%%%%%%%%%%%%%%%%%%%%%%%%%%%%%%%%%%%%%%%%%%%%%%%%%%%%%%%%%%%%%
% DODATNE DEFINICIJE
%%%%%%%%%%%%%%%%%%%%%%%%%%%%%%%%%%%%%%%%%%%%%%%%%%%%%%%%%%%%%%%%%%%%%%%%%%%%%%%
% naložite dodatne pakete, ki jih potrebujete
\usepackage{units}        % fizikalne enote kot \unit[12]{kg} s polovico nedeljivega presledka, glej primer v kodi
\usepackage{graphicx}     % za slike
\usepackage{tikz}
\usetikzlibrary{arrows.meta}
\usepackage{tikz-3dplot}
% VEČ ZANIMIVIH PAKETOV
% \usepackage{array}      % več možnosti za tabele
% \usepackage[list=true,listformat=simple]{subcaption}  % več kot ena slika na figure, omogoči slika 1a, slika 1b
% \usepackage[all]{xy}    % diagrami
% \usepackage{doi}        % za clickable DOI entrye v bibliografiji
% \usepackage{enumerate}     % več možnosti za sezname

% Za barvanje source kode
% \usepackage{minted}
% \renewcommand\listingscaption{Program}

% Za pisanje psevdokode
% \usepackage{algpseudocode}  % za psevdokodo
% \usepackage{algorithm}
% \floatname{algorithm}{Algoritem}
% \renewcommand{\listalgorithmname}{Kazalo algoritmov}

% deklarirajte vse matematične operatorje, da jih bo LaTeX pravilno stavil
\DeclareMathOperator{\bh}{\mathcal{B} (\mathcal{H})}
\DeclareMathOperator{\bk}{\mathcal{B} (\mathcal{K})}
\DeclareMathOperator{\imag}{Im}
\DeclareMathOperator{\real}{Re}
\DeclareMathOperator{\im}{im}
\newcommand{\quot}[2]{{\raisebox{0em}{$#1$}\left/\raisebox{0em}{$#2$}\right.}}
\DeclareMathOperator{\sa}{sa}
\DeclareMathOperator{\lin}{Lin}
\DeclareMathOperator{\PP}{PP}
\DeclareMathOperator{\po}{po}
\DeclareMathOperator{\PO}{PO}
\DeclareMathOperator{\id}{id}
\DeclareMathOperator{\ev}{ev}
\DeclareMathOperator{\co}{co}
\DeclareMathOperator{\supp}{supp}
\DeclareMathOperator{\tr}{Tr}
\DeclareMathOperator{\env}{env}
\DeclareMathOperator{\ext}{ext}

\DeclareMathOperator{\EPP}{EPP}
\DeclareMathOperator{\diag}{diag}



\newcommand{\mapsfrom}{\mathrel{\reflectbox{\ensuremath{\mapsto}}}}

% vstavite svoje definicije ...
\newcommand{\R}{\mathbb R}
\newcommand{\N}{\mathbb N}
\newcommand{\Z}{\mathbb Z}
% Lahko se zgodi, da je ukaz \C definiral že paket hyperref,
% zato dobite napako: Command \C already defined.
% V tem primeru namesto ukaza \newcommand uporabite \renewcommand
\newcommand{\C}{\mathbb C}
\newcommand{\Q}{\mathbb Q}
\newcommand{\T}{\mathbb T}
\newcommand{\F}{\mathbb F}
\newcommand{\E}{\mathbb E}
\newcommand{\A}{\mathbb A}
\newcommand{\D}{\mathbb D}
\newcommand{\oa}{\mathbf{OA}}
\newcommand{\op}{\mathbf{OP}}
\newcommand{\os}{\mathbf{OS}}
\hyphenation{al-geb-rah}





%%%%%%%%%%%%%%%%%%%%%%%%%%%%%%%%%%%%%%%%%%%%%%%%%%%%%%%%%%%%%%%%%%%%%%%%%%%%%%%
% ZAČETEK VSEBINE
%%%%%%%%%%%%%%%%%%%%%%%%%%%%%%%%%%%%%%%%%%%%%%%%%%%%%%%%%%%%%%%%%%%%%%%%%%%%%%%

\begin{document}


\section{Uvod}

Znameniti Gelfandov izrek pravi, da je vsaka komutativna $C^*$-algebra 
$*$-izomorfna $C^*$-algebri zveznih funkcij na kompaktnem Hausdorffovem prostoru $X$.
To vodi do fundamentalne ideje \emph{nekomutativne geometrije}: na splošno (ne nujno komutativno) $C^*$-algebro
lahko gledamo kot na algebro funkcij na nekomutativni posplošitvi pojma kompaktnega Hausdorffovega topološkega prostora.
Nekatere klasične geometrijske objekte, ki jih lahko formuliramo v jeziku komutativnih $C^*$-algeber $C(X)$, 
tako vpeljemo tudi v nekomutativnem kontekstu. 
Glavni cilj tega magistrskega dela je izpostaviti dva takšna geometrijska pojma, \emph{rob Šilova} in \emph{Choqueta},
in vpeljati njuna nekomutativna analoga.

Ključen pojem, ki ga bomo pri tem potrebovali, so \emph{povsem pozitivne preslikave}.
Eden izmed načinov, na katere lahko motiviramo povsem pozitivne preslikave motiviramo, pride iz nekomutativne verjetnosti.
Povsem pozitivne preslikave so namreč nekomutativna posplošitev pojma {Markovskega jedra} iz klasične verjetnosti.
Običajnemu verjetnostnemu prostoru $(X, \mathcal{F}, \mu)$ lahko priredimo $C^*$-algebro $\mathcal{B} (X)$,
ki jo sestavljajo vse omejene kompleksne naključne spremenljivke na $X$, in pričakovano vrednost $\E: \mathcal{B} (X) \to \C$.
To nam omogoča, da na verjetnostne prostore gledamo iz bolj abstraktnega vidika
kot par $(\mathfrak{A}, \E)$,
kjer je $\mathfrak{A}$ $C^*$-algebra in $\E: \mathfrak{A} \to \C$ pozitiven funkcional,
za katerega velja $\E (1) = 1$ (takšnim funkcionalom v žargonu pravimo \emph{stanja}).
Izkaže se, da lahko večino običajnih pojmov iz verjetnostne teorije, kot na primer \emph{dogodke}, \emph{slučajne spremenljivke}, ali pa \emph{pogojno pričakovano vrednost}, 
prevedemo tudi na nekomutativni kontekst
(za več o tej temi glej \cite{nicaspeicher2006}). Oglejmo si Markovsko jedro za merljiva prostora $(X, \mathcal{F})$ in $(Y, \mathcal{G})$,
ki je definirano kot funkcija $K: Y \times \mathcal{F} \to [0, 1]$,
tako da je za vsak $A \in \mathcal{F}$ preslikava $y \mapsto K(y, A)$ $\mathcal{G}$-merljiva in 
za vsak $y \in Y$ preslikava $B \mapsto K(y, B)$ verjetnostna mera na $(X, \mathcal{F})$.
To inducira preslikavo 
$$\varphi: \mathcal{B}(X) \to \mathcal{B}(Y),\quad (\varphi f)(y) = \int_X f(x)\, K(y, dx),$$
ki ima lastnost, da za poljubne funkcije $a_1, \dots, a_n \in \mathcal{B}(X)$ in $b_1, \dots, b_n \in \mathcal{B}(Y)$ velja 
\begin{equation}\label{eq:intro1}
    \sum_{i, j = 1} ^n b_i ^* \varphi(a_i ^* a_j) b_j \geq 0.
\end{equation}
Linearnim preslikavam $\varphi: \mathfrak{A} \to \mathfrak{B}$ med $C^*$-algebrami,
za katere velja lastnost \eqref{eq:intro1}, pravimo povsem pozitivne preslikave.
Povsem pozitivne preslikave imajo med drugim aplikacije v opisu \emph{odprtih kvantnih sistemov} s pomočjo
\emph{kvantnih dinamičnih polgrup} %, ki so nekomutativna posplošitev klasičnih markovskih polgrup 
(\cite{attaljoyepillet2006}),
ali pa v \emph{kvantni teoriji informacij}, kjer so znani kot \emph{kvantni kanali} (\cite{nielsenchuang2010}).

Pri organizaciji snovi v tem magistrskem delu sem želel okvirno slediti zgodovinskemu zaporedju teorije povsem pozitivnih preslikav.
Za njen začetek lahko vzamemo Stinespringov članek \cite{stinespring1955} iz leta 1955
in Arvesonov članek \cite{arveson1969} iz leta 1969 -- na slednjega se bomo sklicevali skozi celotno magistrsko delo.
Posebnega pomena sta Stinespringov izrek, ki je posplošitev Gelfand-Naimark-Segalovega izreka,
in Arvesonov razširitveni izrek, ki je posplošitev Hahn-Banachovega izreka za povsem pozitivne preslikave.
Posledica teh dveh izrekov je, da lahko povsem pozitivne preslikave $\varphi: \mathfrak{A} \to \mathfrak{B}$ smatramo za naravno posplošitev pozitivnih funkcionalov oz.~stanj
na $C^*$-algebrah, le da je kodomena $\mathfrak{B}$ sedaj poljubna $C^*$-algebra. Osnovne primere in lastnosti povsem pozitivnih preslikav 
obravnavamo v poglavju \ref{sec:2}.

Iz Arvesonovega dela je sledilo, da naravna domena za povsem pozitivne preslikave niso celotne $C^*$-algebre,
temveč njihovi enotski sebiadjungirani podprostori, ki jim v žargonu pravimo \emph{operatorski sistemi}.
Kot bomo videli, so povsem pozitivne preslikave morfizmi v kategoriji operatorskih sistemov.
Operatorski sistemi se sicer pojavijo tudi v nekomutativni geometriji%, ko obravnavamo \emph{spektralne trojice} -- to so trojice $(\mathfrak{A}, \mathcal{H}, D)$,
%kjer je $\mathfrak{A}$ enotska $*$-algebra omejenih operatorjev na Hilbertovem prostoru $\mathcal{H}$ in $D$ neomejen sebiadjungiran operator, 
%tako da je njegova resolventa $(D + i)^{-1}$ kompaktna in je $[D, a] = Da - aD$ omejen operator za vsak $a \in \mathfrak{A}$ (DOPOLNI).
%Če s $P$ označimo projekcijo na lastne podprostore $D$, potem želimo strukturo spektralne trojice $(\mathfrak{A}, \mathcal{H}, D)$
%izluščiti s pomočjo trojic $(P \mathfrak{A} P, P \mathcal{H}, PDP)$, toda $P\mathfrak{A}P$ v tem primeru ni več $*$-algebra,
%temveč operatorski sistem 
(za več o tem pristopu glej \cite{connes2021}), najdemo jih pa tudi v kvantni teoriji informacij kot \emph{kvantne grafe} (\cite{weaver2021}, \cite{duanseveriniwinter2013}).
Osnovne lastnosti operatorskih sistemov obravnavamo v poglavju \ref{sec:3}.

Operatorski sistemi so poseben primer operatorskih prostorov,
ki so enotski (ne nujno sebiadjungirani) podprostori $C^*$-algeber.
Vsakemu operatorskemu prostoru $\mathcal{M} \subseteq \mathfrak{A}$ 
pa lahko priredimo operatorski sistem $\mathcal{M} + \mathcal{M}^*$.
Morfizmi na operatorskih prostorih bodo povsem omejene preslikave.
Preplet povsem omejenih preslikav in povsem pozitivnih preslikav bomo obravnavali v poglavju \ref{sec:4}.
Glavna rezultata tega razdelka bosta, da so vse povsem pozitivne preslikave povsem omejene 
in da lahko vsako enotsko povsem skrčitev na operatorskem sistemu $\mathcal{M}$
razširimo do povsem pozitivne preslikave na operatorskem sistemu $\mathcal{M} + \mathcal{M}^*$.
Od 1980-ih let naprej so Ruan, Paulsen, Wittstock in ostali avtorji podrobneje razvijali
teorijo popolnoma omejenih preslikav in operatorskih prostorov, ki jo 
lahko razumemo kot tesen analog teorije popolnoma pozitivnih preslikav in operatorskih sistemov.
V tem magistrskem delu se s tem ne bomo ukvarjali; več podrobnosti lahko bralec najde v \cite{paulsen2002} ali
\cite{pisier2003}.

V krajšem razdelku \ref{sec:5} bomo nato dokazali Stinespringovo karakterizacijo povsem pozitivnih 
preslikav na komutativnih $C^*$-algebrah. Povsem pozitivne preslikave $\varphi: \mathfrak{A} \to \mathfrak{B}$,
kjer je vsaj ena izmed $C^*$-algeber $\mathfrak{A}, \mathfrak{B}$ komutativna, imajo posebno lep opis; 
sovpadajo namreč z navadnimi pozitivnimi preslikavami.

Sledi poglavje \ref{sec:6} o dilatacijah.
Teorija dilatacij, ki sta jo začela Sz.-Nagy in Foias v 1950-ih letih (\cite{nagyfoias2010}),
se okvirno ukvarja s tem, kdaj lahko 
operator $T$ na Hilbertovem prostoru $\mathcal{H}$ realiziramo kot $P_{\mathcal{H}} S |_{\mathcal{H}}$,
kjer je $S$ operator na večjem Hilbertovem prostoru $\mathcal{K} \subseteq \mathcal{H}$
in $P_{\mathcal{H}}$ ortogonalna projekcija s $\mathcal{K}$ na $\mathcal{H}$.
Ponavadi želimo, da ima operator $S$ (ki mu pravimo \emph{dilatacija} $T$)
še nekatere dodatne lastnosti: obravnavali bomo Sz.-Nagyjev zgled,
ki nam pove, da lahko vsak skrčitven operator dilatiramo do izometrije
in vsako izometrijo do unitarnega operatorja. 
Glavni izrek tega poglavja je že zgoraj omenjeni Stinespringov izrek, ki povezuje teorijo dilatacij s povsem pozitivnimi preslikavami.
Dilatacije srečamo med drugim v mehaniki,
kjer lahko evolucijsko polgrupo ireverzibilnega sistema pod določenimi pogoji dilatiramo do unitarne evolucijske grupe (večjega) reverzibilnega sistema 
(glej recimo \cite{MielkePeletierZimmer2025}).

Arvesonovem delu iz leta 1969 sta zgodovinsko sledila Krausov članek \cite{kraus1971} iz leta 1971 
ter Choijev članek \cite{choi1975} iz leta 1975, v katerem sta avtorja klasificirala povsem pozitivne preslikave $\varphi: M_n (\C) \to M_m(\C)$.
Njuno delo med drugim povzamemo v poglavju \ref{sec:7}, v katerem obravnavamo povsem pozitivne preslikave $\varphi: \mathfrak{A} \to \mathfrak{B}$,
kjer je vsaj ena izmed $C^*$-algeber $\mathfrak{A}, \mathfrak{B}$ oblike $M_n (\C)$. 
Posebnega pomena bo trditev, ki pravi, da lahko povsem pozitivno preslikavo $\varphi: \mathcal{S} \to M_n (\C)$
na poljubnem operatorskem sistemu $\mathcal{S} \subseteq \mathfrak{A}$ razširimo do povsem pozitivne preslikave na celotni $C^*$-algebri $\mathfrak{A}$.
V poglavju \ref{sec:8} bomo to trditev dokazali še za povsem pozitivne preslikave $\varphi: \mathcal{S} \to \bh$,
s čimer bomo dokazali Arvesonov razširitveni izrek.

Naslednji pomemben članek za teorijo povsem pozitivnih preslikav sta leta 1977 napisala Choi in Effros (\cite{choieffros1977}),
v katerem sta avtorja dokazala, da imajo vsak operatorski sistem abstraktno karakterizacijo, ki je torej neodvisna od konkretne $C^*$-algebre,
v kateri je operatorski sistem vložen. Choi-Effrosevo abstraktno karakterizacijo obravnavamo v razdelku \ref{sec:9}.
Zato se poraja vprašanje, ali obstaja $C^*$-algebra, v katero lahko vložimo abstrakten operatorski sistem, in ki je na nek način `minimalna'.
Ekvivalentno vprašanje lahko postavimo za konkreten operatorski sistem $\mathcal{S} \subseteq C^*(\mathcal{S})$: ali obstaja minimalna $C^*$-algebra 
$\mathfrak{A}$, tako da obstaja enotska povsem izometrija $i: \mathcal{S} \to \mathfrak{A}$ in je $\mathfrak{A} = C^*(i(\mathcal{S}))$?
Takšni $C^*$-algebri $\mathfrak{A}$ bomo rekli \emph{$C^*$-ogrinjača}.

Problem obstoja $C^*$-ogrinjače je obravnaval že Arveson v \cite{arveson1969}.
V ta namen je vpeljal nekomutativen analog geometrijskega pojma robu Šilova, ki smo ga omenili na začetku uvoda.
Če je $\mathcal{S} \subseteq C(X)$ operatorski sistem, ki loči točke na kompaktnem Hausdorffovem kompaktu $X$,
potem je njegov rob Šilova definiran kot najmanjša zaprta podmnožica $Y \subseteq X$,
tako da vsaka funkcija iz $\mathcal{S}$ na $Y$ doseže svojo supremum normo.
Rob Šilova izhaja iz preučevanja \emph{uniformiranih algeber}:
to so enotske podalgebre (ne nujno zaprte za $*$) $\mathcal{A} \subseteq C(X)$, ki {ločijo točke na $X$}
(za več o tem glej \cite{gamelin1969}). Da lahko vpeljemo nekomutativen rob Šilova,
moramo njegovo definicijo prevesti v besednjak $C^*$-algeber.
To, da operatorski sistem $\mathcal{S} \subseteq C(X)$ loči točke na $X$,
je po Stone-Weierstrassovem izreku ekvivalentno temu, da $\mathcal{S}$ generira $C^*$-algebro $C(X)$.
Ker so zaprte podmnožica $X$ v bijektivni korespondenci z zaprtimi ideali v $C(X)$, lahko vpeljemo 
nekomutativen rob Šilova za poljuben operatorski sistem $\mathcal{S} \subseteq C^*(\mathcal{S})$ kot zaprt ideal $I \subseteq C^*(\mathcal{S})$. 
Arveson je opazil, da če za operatorski sistem $\mathcal{S} \subseteq C^*(\mathcal{S})$ obstaja $C^*$-ogrinjača $C^* _{\env} (\mathcal{S})$,
potem obstaja tudi njegov ideal Šilova $I \subseteq C^*(\mathcal{S})$ in velja $\quot{C^*(\mathcal{S})}{I} \cong C_{\env} ^* (\mathcal{S})$.
V svojem delu je dokazal obstoj ideala Šilova le za določen tip operatorskih sistemov, vendar pa je vprašanje obstoja ideala Šilova in $C^*$-ogrinjače 
za splošen operatorski sistem ostalo nerešeno do
leta 1979, ko je nanj pritrdilno odgovoril Hamana v članku \cite{hamana1979} s svojo teorijo injektivnih ogrinjač.

Arveson je poleg robu Šilova vpeljal še nekomutativen analog tesno povezanega geometrijskega pojma: Choquetovega robu.
Choquetov rob je pojem iz Choquetove teorije, ki se okvirno ukvarja s problemi naslednjega tipa:
\emph{Naj bo $E$ lokalno konveksen prostor in $K \subseteq E$ njegova konveksna, kompaktna podmnožica.
Ali za vsako točko $x \in K$ obstaja verjetnostna mera $\mu$ na $K$, skoncentrirana na $\ext K$, tako da za vsak omejen linearen funkcional $f \in E^*$ velja $\int_X f\, d\mu = f(x)$?}
Pomemben rezultat v tej smeri je Choquetov izrek (\cite[str. 14]{phelps2001}), ki na to vprašanje odgovori pritrdilno, v kolikor je $E$ metrizabilen prostor (več o tej temi v knjigi \cite{phelps2001}).
Ponovno vzamemo operatorski prostor $\mathcal{S} \subseteq C^*(\mathcal{S})$, ki loči točke na $X$. 
Njegov Choquetov rob sestavljajo tiste točke $x \in X$, za katere velja:
če je $\mu$ regularna Borelova verjetnostna mera, tako da je $$f(x) = \int_X f\, d\mu,\quad \forall f \in \mathcal{S},$$
potem je $\mu = \delta_x$. Choquetov rob operatorskega sistema ima to lastnost, da je njegovo zaprtje enako robu Šilova.
Arveson je zato v svojem članku iz leta 1969 vpeljal \emph{robne upodobitve} operatorskega sistema, ki predstavljajo posplošitev točk v Choquetovem robu 
za operatorski sistem funkcij. V tem članku je postavil tudi domnevo, da za poljuben operatorski sistem $\mathcal{S}$ obstaja \emph{dovolj} robnih upodobitev:
to pomeni, da je presek vseh (do unitarne ekvivalence) robnih upodobitev $\mathcal{S}$ enak idealu Šilova.

Kljub temu, da je Hamana dokazal obstoj ideala Šilova, je ta domneva dlje časa ostala nerešena.
Pomemben preboj je nastal leta 2002, ko sta Dritschel in McCollough v članku \cite{dritschel2005}
konstruirala popolnoma nov dokaz obstoja $C^*$-ogrinjače s pomočjo teorije dilatacij. 
Njuno konstrukcijo za operatorske algebre je nato Arveson leta 2003 v neobjavljenem zapisu \cite{arveson2003} posplošil 
na operatorske sisteme. V poglavju \ref{sec:10} predstavimo njihov dokaz, ki mu sledi dokaz obstoja ideala Šilova v poglavju \ref{sec:11}.

Za razliko od Hamanine teorije injektivnih ogrinjač nam argumenti Arvesona, Dritschela in McCollougha omogočajo konstrukcijo robnih upodobitev.
Arveson je nato leta 2008 v članku \cite{arveson2008} dokazal obstoj dovolj robnih upodobitev za separabilne operatorske sisteme. 
Za splošne operatorske sisteme sta trditev dokazala Davidson in Kennedy leta 2015 v članku \cite{davidson2015},
s čimer sta rešila 46-letno domnevo na področju operatorskih algeber.
V zadnjem poglavju \ref{sec:12} predstavimo njuno rešitev.









\section{Povsem pozitivne preslikave}\label{sec:2}

V tem magistrskem delu bomo, če ni izrecno napisano drugače,
za vse $C^*$-algebre predpostavili, da so enotske in nad kompleksnimi skalarji.
Prav tako bomo vsi $*$-homomorfizmi $C^*$-algeber $\varphi: \mathfrak{A} \to \mathfrak{B}$ enotski,
torej bodo slikali $\varphi(1) = 1$. Vse vektorske prostore bomo obravnavali nad kompleksnimi skalarji. 

Naj bo $\mathfrak{A}$ $C^*$-algebra. Definiramo $M_n (\mathfrak{A})$ kot algebro $n \times n$ matrik z elementi v $\mathfrak{A}$, z običajnim matričnim seštevanjem in množenjem.
Na $M_n(\mathfrak{A})$ lahko definiramo matrično involucijo
$$[a_{ij}]_{i, j} ^* := [a_{ji} ^*]_{i, j},$$
s katero $M_n (\mathfrak{A})$ postane $*$-algebra.

\begin{definicija}
    Naj bosta $\mathfrak{A}, \mathfrak{B}$ $C^*$-algebri in $\phi: \mathfrak{A} \to \mathfrak{B}$ linearna preslikava.
    Za vsak $n \in \N$ definiramo 
    $$\phi^{(n)}: M_n (\mathfrak{A}) \to M_n (\mathfrak{B}),\quad [a_{ij}]_{i, j} \mapsto [\phi(a_{ij})]_{i, j}.$$
    Preslikavi $\phi^{(n)}$ pravimo $n$-ta \emph{ampliacija} $\phi$.
\end{definicija}

Če je $\mathfrak{A}$ $C^*$-algebra, potem po Gelfand--Naimark--Segalovem izreku (\cite[Theorem VIII.5.14., str. 250]{conway1990}, v nadaljevanju: izrek GNS)
obstaja Hilbertov prostor $\mathcal{H}$ in injektivna upodobitev $\pi: \mathfrak{A} \to \bh$, kar inducira injektivni $*$-homomorfizem 
$$\pi^{(n)}: M_n (\mathfrak{A}) \to M_n (\bh)$$
za vsak $n \in \N$. To pomeni, da lahko $M_n (\mathfrak{A})$ identificiramo z
$*$-podalgebro $C^*$-algebre $M_n (\bh)$.
Ker je slika $\pi^{(n)} (M_n (\mathfrak{A}))$ zaprta v $M_n (\bh)$,
je tudi $C^*$-podalgebra v $M_n (\bh)$.
To nam inducira normo na $M_n (\mathfrak{A})$, s katero $M_n (\mathfrak{A})$ postane $C^*$-algebra.
Spomnimo, da vsaka $*$-algebra premore največ eno normo, za katero postane $C^*$-algebra (glej recimo \cite[Corollary 2.1.2., str. 37]{murphy1990}),
kar pomeni, da je zgoraj opisana norma na $M_n (\mathfrak{A})$ kanonična oziroma neodvisna od izbora upodobitve $\pi$.
Kot bomo videli, nam kanonične norme na $M_n (\mathfrak{A})$ povedo dodatne informacije o začetni $C^*$-algebri $\mathfrak{A}$.

\begin{definicija}
    Element $a \in \mathfrak{A}$ je \emph{pozitiven}, če obstaja $b \in \mathfrak{A}$, tako da je $a = b^* b$.
    Množico vseh pozitivnih elementov v $\mathfrak{A}$ označimo z $\mathfrak{A}_+$.
\end{definicija}

\begin{opomba}
    Množica $\mathfrak{A}_+$ je \emph{konveksen stožec}: to pomeni, da je $\mathfrak{A}_+ + \mathfrak{A}_+ \subseteq \mathfrak{A}_+$,
    $c \mathfrak{A}_+ \subseteq \mathfrak{A}_+$ za $\alpha \geq 0$ ter $\mathfrak{A}_+ \cap - \mathfrak{A}_+ = \{0\}$.
\end{opomba}

\begin{definicija}
    Naj bosta $\mathfrak{A} , \mathfrak{B}$ $C^*$-algebri in $\varphi: \mathfrak{A} \to \mathfrak{B}$ linearna preslikava.
    \begin{enumerate}
        \item $\varphi$ je \emph{pozitivna}, če je $\varphi(\mathfrak{A}_+) \subseteq \mathfrak{B}_+$;
        \item $\varphi$ je \emph{$n$-pozitivna}, če je $\varphi^{(n)}$ pozitivna;
        \item $\varphi$ je \emph{povsem pozitivna}, če je $n$-pozitivna za vsak $n \in \N$.
    \end{enumerate}
    Množico vseh (enotskih) povsem pozitivnih preslikav označimo s $\PP (\mathfrak{A}, \mathfrak{B})$ (oziroma $\EPP (\mathfrak{A} , \mathfrak{B})$).
\end{definicija}

\begin{opomba}
    \begin{itemize}
        \item Če je $\varphi$ $n$-pozitivna, potem je tudi $k$-pozitivna za vse $k \leq n$.
        \item Kompozitum pozitivnih ($n$-pozitivnih, povsem pozitivnih) preslikav je prav tako pozitiven ($n$-pozitiven, povsem pozitiven).
        \item Če sta $\varphi, \psi: \mathfrak{A} \to \mathfrak{B}$ pozitivni ($n$-pozitivni, povsem pozitivni) preslikavi in $\alpha, \beta \geq 0$,
        potem je tudi $\alpha \varphi + \beta \psi$ pozitivna ($n$-pozitivna, povsem pozitivna).
    \end{itemize}
\end{opomba}

V nadaljevanju si bomo oglejmo nekaj primerov povsem pozitivnih preslikav
in njihovih osnovnih lastnosti. Model povsem pozitivnih preslikav so pozitivni funkcionali. 

\begin{zgled}\label{zgl:1.3}
    Dokažimo, da je vsak pozitivni linearni funkcional $\varphi$ na $C^*$-algebri $\mathfrak{A}$ povsem pozitiven.
    Dovolj je pokazati, da je $$\varphi^{(n)} : M_n (\mathfrak{A}) \to M_n (\C)$$
    pozitivna preslikava. Ključen je uvid, da za popolnoma poljuben funkcional $\varphi$, matriko $A = [a_{ij}]_{i, j} \in M_n (\mathfrak{A})$
    in vektorja $$\alpha = \begin{bmatrix}
        \alpha_1 \\ \vdots \\ \alpha_n
    \end{bmatrix},\quad \beta = \begin{bmatrix}
        \beta_1 \\ \vdots \\ \beta_n
    \end{bmatrix} \in \C^n$$ velja 
    \begin{align*}
        \beta^* \varphi^{(n)} (A) \alpha &= \begin{bmatrix}
            \beta_1  \\ \vdots \\ \beta_n 
        \end{bmatrix}^*
        \begin{bmatrix}
            \varphi (a_{11}) & \cdots & \varphi(a_{1n})\\
            \vdots & & \vdots\\
            \varphi (a_{n1}) & \cdots & \varphi(a_{nn})
        \end{bmatrix}
        \begin{bmatrix}
            \alpha_1 \\ \vdots \\ \alpha_n
        \end{bmatrix} \\
        &= \sum_{i, j} \overline{\beta_i} \varphi(a_{ij}) \alpha_j  = \varphi \left(\sum_{i, j} \overline{\beta_i} a_{ij} \alpha_j \right)\\
        &= \varphi(\beta^* A \alpha).
    \end{align*}
    Če je sedaj matrika $A \in M_n (\mathfrak{A})$ pozitivna, potem je za vsak $\alpha \in \C^n$
    tudi $\alpha^* A \alpha \in \mathfrak{A}$ pozitiven element in velja 
    $$\alpha^* \varphi^{(n)}(A) \alpha = \varphi(\alpha^* A \alpha) \geq 0.$$
    To pomeni, da je $\varphi^{(n)} (A) \in M_n (\C)$ pozitivna matrika, zato je $\varphi^{(n)}: M_n (\mathfrak{A}) \to M_n (\C)$ pozitivna preslikava.
\end{zgled}

To nam omogoča, da na povsem pozitivne preslikave gledamo kot na posplošitve stanj,
pri čemer je kodomena poljubna $C^*$-algebra. Izkazalo se bo, da lahko teorijo stanj posplošimo na 
povsem pozitivne preslikave.

\begin{zgled}
    Če je $\varphi: \mathfrak{A} \to \mathfrak{B}$ $*$-homomorfizem, je $\varphi$ pozitiven, saj 
    $\varphi(a^* a) = \varphi(a) ^* \varphi(a)$. Še več, v tem primeru je tudi $\varphi^{(n)}$ $*$-homomorfizem,
    zato je $\varphi^{(n)}$ pozitivna za vsak $n\in \N$. Od tod sledi,
    da je vsak $*$-homomorfizem med dvema $C^*$-algebrama povsem pozitiven.
\end{zgled}

\begin{zgled}\label{zgl:1.4}
    Obstajajo pozitivne preslikave, ki niso povsem pozitivne.
    Primer take preslikave je $$\varphi: M_2 (\C) \to M_2 (\C),\quad A \mapsto A^\top$$
    (tukaj operacija $A^\top$ pomeni realno transpozicijo matrike brez konjugiranja, saj mora biti $\varphi$ linearna preslikava).
    Očitno je $\varphi$ pozitivna preslikava, kar pa ne velja za njeno drugo ampliacijo
    $$\varphi^{(2)}:  M_4 (\C) \to M_4 (\C),\quad \begin{bmatrix}
        A_{11} & A_{12}\\
        A_{21} & A_{22}
    \end{bmatrix} \mapsto \begin{bmatrix}
        A_{11}^\top & A_{12}^\top\\
        A_{21}^\top & A_{22}^\top
    \end{bmatrix}.$$ Ta namreč slika pozitivno matriko
    $$\begin{bmatrix}
        1 & 0 & 0 & 1\\
        0 & 0 & 0 & 0\\
        0 & 0 & 0 & 0\\
        1 & 0 & 0 & 1
    \end{bmatrix} = \begin{bmatrix}
        1 \\ 0 \\ 0 \\ 0
    \end{bmatrix} \begin{bmatrix}
        1 \\ 0 \\ 0 \\ 0
    \end{bmatrix}^* \geq 0$$
    v
    $$\begin{bmatrix}
        1 & 0 & 0 & 0\\
        0 & 0 & 1 & 0\\
        0 & 1 & 0 & 0\\
        0 & 0 & 0 & 1
    \end{bmatrix},$$
    ki ima negativno lastno vrednost $-1$ (in zato seveda ne more biti pozitivna).
\end{zgled}

\begin{zgled}\label{zgl:1.5}
    Naj bosta $\mathcal{H}, \mathcal{K}$ Hilbertova prostora in $V \in \mathcal{B}(\mathcal{K}, \mathcal{H})$
    omejen operator. Dokazujemo, da je preslikava
    $$\varphi: \bh \to  \bk,\quad A \mapsto V^* A V$$
    povsem pozitivna. Vzemimo pozitivno matriko $[A_{ij}]_{i, j} \in M_n (\bh)_+$.
    Potem je 
    \begin{align*}
        \varphi^{(n)} ([A_{ij}]_{i, j}) &= \begin{bmatrix}
            \varphi(A_{11}) & \cdots & \varphi(A_{1n})\\
            \vdots & & \vdots\\
            \varphi(A_{n1}) & \cdots & \varphi(A_{nn})\\
        \end{bmatrix}\\
        &= \begin{bmatrix}
            V^* A_{11} V & \cdots & V^* A_{1n} V\\
            \vdots & & \vdots\\
            V^* A_{n1} V & \cdots & V^* A_{nn} V\\
        \end{bmatrix}\\
        &= \begin{bmatrix}
            V & &\\
             & \ddots & \\
             & & V
        \end{bmatrix}^* \begin{bmatrix}
            A_{11} & \cdots & A_{1n}\\
            \vdots & & \vdots\\
            A_{n1} & \cdots & A_{nn}\\
        \end{bmatrix} \begin{bmatrix}
            V & &\\
             & \ddots & \\
             & & V
        \end{bmatrix}\\
    \end{align*}
    pozitivna matrika, torej je $\varphi ^{(n)}$ pozitivna preslikava. 
    Posledično je $\varphi$ povsem pozitivna.

    Če je $\mathfrak{A}$ $C^*$-algebra in $\psi: \mathfrak{A} \to \bh$ povsem pozitivna preslikava,
    potem je tudi kompozicija 
    $$
        \mathfrak{A} \to  \bk,\quad A \mapsto V^* \psi (a) V 
    $$
    povsem pozitivna.
\end{zgled}

Razdelek zaključimo z uporabno karakterizacijo povsem pozitivnih preslikav \eqref{eq:intro1}
iz uvoda.


\begin{comment}
\begin{zgled}
    Naj bo $\psi: \mathfrak{A} \to \mathcal{B}$ povsem pozitivna in $b \in B$.
    Definirajmo $$\varphi :A \to B ,\quad a \mapsto b^* \psi(a) b.$$
    Dokazujemo, da je $\varphi$ povsem pozitivna. Za vsak $n \in \N$ in $[a_{ij}]_{i, j} \in M_n (A)_+$
    je namreč
    \begin{align*}
        \varphi^{(n)} ([a_{ij}]_{i, j}) &= [\varphi(a_{ij})]_{i, j} = \begin{bmatrix}
            b^* \psi (a_{11})  b & \cdots & b^* \psi (a_{1n}) b\\
            \vdots & \ddots & \vdots\\
            b^* \psi (a_{n1})  b & \cdots & b^* \psi (a_{nn}) b
        \end{bmatrix}\\
        &= \begin{bmatrix}
            b & & \\
             & \ddots & \\
              & & b
        \end{bmatrix}^* (\psi^{(n)} [(a_{ij})]_{i, j}) \begin{bmatrix}
            b & & \\
              & \ddots & \\
              & & b
        \end{bmatrix} \geq 0.
    \end{align*}
\end{zgled}
\end{comment}


\begin{lema}\label{lem:1.2}
    Če je $\mathfrak{A}$ $C^*$-algebra, potem je vsak pozitiven element $M_n (\mathfrak{A})$
    vsota $n$ pozitivnih matrik oblike $[a_i ^* a_j]_{i, j}$ za 
    neke $\{a_1, \dots, a_n\} \subseteq A$.
\end{lema}

\begin{dokaz}
    Najprej opazimo, da če imamo matriko $$R = \begin{bmatrix}
        0 & \cdots & 0\\
        \vdots & & \vdots\\
        a_1 & \cdots & a_n\\
        \vdots & & \vdots\\
        0 & \cdots & 0\\       
    \end{bmatrix} \in M_n (\mathfrak{A}),$$
    potem je $R^* R = [a_i^* a_j]_{i, j}$.
    Naj bo sedaj $P \in M_n (\mathfrak{A})_+$ poljubna pozitivna matrika, torej $P = B^* B$ za nek $B \in M_n (\mathfrak{A})$.
    Zapišemo $B = R_1 + \cdots + R_n$, kjer je $R_k$ matrika,
    ki ima $k$-to vrstico enako kot $B$, drugje pa je ničelna.
    Ker velja $R_k^* R_l = 0$ za poljubna različna indeksa $k \neq l$ iz $\{1, \dots, n\}$,
    sledi
    \begin{align*}
        P = B^* B &= (R_1 ^* + \cdots + R_n ^*) (R_1 + \cdots + R_n) \\
        &= R_1 ^* R_1 + \cdots + R_n ^* R_n. \qedhere
    \end{align*}
\end{dokaz}



\begin{lema}\label{lem:1.3}
    Naj bo $\mathfrak{B}$ $C^*$-algebra in $A = [a_{ij}]_{i, j} \in M_n(\mathfrak{B})$.
    Potem je $A$ pozitivna natanko tedaj, ko velja $$\begin{bmatrix}
        b_1 \\ \vdots \\ b_n
    \end{bmatrix}^* A \begin{bmatrix}
        b_1 \\ \vdots \\ b_n
    \end{bmatrix} = \sum_{i, j} ^n b_i ^* a_{ij} b_j \geq 0$$
    za poljubno $n$-terico $b_1, \dots, b_n \in \mathfrak{B}$.
\end{lema}

\begin{dokaz}
    Implikacija v levo je trivialna, zato dokazujemo implikacijo v desno.
    \begin{enumerate}
        \item  V prvem koraku dokažemo, da če je $A = [a_{ij}]_{i, j} \in M_n (\mathfrak{B})$ takšna matrika, da velja $$\begin{bmatrix}
        b_1 \\ \vdots \\ b_n
    \end{bmatrix}^* A \begin{bmatrix}
        b_1 \\ \vdots \\ b_n
    \end{bmatrix} = \sum_{i, j} ^n b_i ^* a_{ij} b_j = 0$$
    za poljubno $n$-terico $b_1, \dots, b_n \in \mathfrak{B}$, potem je $A = 0$.
    Dokazali bomo, da je vsaka komponenta $A$ ničelna. Če v zgornjo enačbo vstavimo $b_i = 1$ in $b_j = 0$ za vse $j \neq i$, potem dobimo
    $a_{ii} = 0$. Od tod sledi, da so vsi diagonalni elementi matrike $A$ ničelni. 
    Preostane še pokazati, da je $a_{ij} = 0$ za poljubna indeksa $i \neq j$.
    Če vstavimo $b_i := 1, b_j := c \in \mathfrak{B}$ in $b_k = 0$ za $k \neq i, j$, dobimo enakost
    $$a_{ii} + a_{ij} c + c^* a_{ji}  + c^* a_{jj} c = 0.$$
    Ker so diagonalni elementi ničelni, sledi
    $$a_{ij} c + c^* a_{ji} = 0.$$
    Če posebej vstavimo $c = 1$ in $c = i$, dobimo
    $$a_{ij} + a_{ji} = 0, \quad i a_{ij} - i a_{ji} = 0,$$
    od koder potem sledi $a_{ij} = a_{ji} = 0$. S tem smo dokazali, da je $A = 0$.
        \item V drugem koraku dokažemo, da če je $$\begin{bmatrix}
            b_1 \\ \vdots \\ b_n
        \end{bmatrix}^* A \begin{bmatrix}
            b_1 \\ \vdots \\ b_n
        \end{bmatrix} = \sum_{i, j} ^n b_i ^* a_{ij} b_j \in \mathfrak{B}$$
        sebiadjungiran element za vsako $n$-terico $b_1, \dots, b_n \in \mathfrak{B}$, potem je $A$ sebiadjungirana matrika.
        Opazimo namreč, da velja
        \begin{align*}
            \sum_{i, j} ^n b_i ^* a_{ij} b_i &= \left(\sum_{i, j} ^n b_i ^* a_{ij} b_i\right)^* \\
            &= \sum_{i, j} ^n b_j ^* a_{ij}^* b_i\\
            &=  \sum_{i, j} ^n b_i ^* a_{ji}^* b_j,
        \end{align*}
        od koder sledi
        $$\begin{bmatrix}
        b_1 \\ \vdots \\ b_n
    \end{bmatrix}^* (A - A^*) \begin{bmatrix}
        b_1 \\ \vdots \\ b_n
    \end{bmatrix} = \sum_{i, j} ^n b_i ^* (a_{ij} - a_{ji}^*) b_i = 0.$$ Po prvemu koraku tako dobimo $A - A^* = 0$.
        \item Nazadnje dokažemo, da če je $$\begin{bmatrix}
            b_1 \\ \vdots \\ b_n
        \end{bmatrix}^* A \begin{bmatrix}
            b_1 \\ \vdots \\ b_n
        \end{bmatrix} = \sum_{i, j} ^n b_i ^* a_{ij} b_j$$
        pozitiven za poljubno $n$-terico $b_1, \dots, b_n \in \mathfrak{B}$, potem je $A$ pozitivna matrika.
        Ker je $A$ po prejšnjem koraku sebiadjungirana, imamo po zveznem funkcijskem računu (glej recimo \cite[Section VIII.2, str. 236-240]{conway1990}) 
        dekompozicijo $A = A_+ - A_-$,
        kjer sta $A_+, A_-$ pozitivni matriki, tako da velja $A_+ A_- = A_- A_+ = 0$.
        Po predpostavki imamo
        $$\begin{bmatrix}
            b_1 \\ \vdots \\ b_n
        \end{bmatrix}^* A_- A A_- \begin{bmatrix}
            b_1 \\ \vdots \\ b_n
        \end{bmatrix}\geq 0,$$
        hkrati pa je izraz na levi strani enak
        $$- \begin{bmatrix}
            b_1 \\ \vdots \\ b_n
        \end{bmatrix}^* A_- ^3 \begin{bmatrix}
            b_1 \\ \vdots \\ b_n
        \end{bmatrix}\leq 0.$$ Od tod sledi
        $$- \begin{bmatrix}
            b_1 \\ \vdots \\ b_n
        \end{bmatrix}^* A_- ^3 \begin{bmatrix}
            b_1 \\ \vdots \\ b_n
        \end{bmatrix} = 0$$ za poljubno $n$-terico $b_1, \dots, b_n \in \mathfrak{B}$.
        Po prvem koraku nato dobimo $A_- ^3 = 0$, od koder po zveznem funkcijskem računu sledi $A_- = 0$.
        To pa pomeni, da je $A = A_+$ pozitivna matrika. \qedhere
    \end{enumerate}
\end{dokaz}

Iz lem \ref{lem:1.2} in \ref{lem:1.3} sledi naslednja trditev.

\begin{trditev}
    Preslikava $\varphi: \mathfrak{A} \to \mathfrak{B}$ je povsem pozitivna natanko tedaj,
    ko za vsak $n \in \N$ in $n$-terici $a_1, \dots, a_n \in \mathfrak{A}$ ter $b_1, \dots, b_n \in \mathfrak{B}$
    velja $$\sum_{i, j} ^n b_i ^* \varphi (a_i ^* a_j) b_j \geq 0.$$
\end{trditev}

\section{Operatorski sistemi}\label{sec:3}

V zgornjih zgledih smo obravnavali povsem pozitivne preslikave na $C^*$-algebrah.
V praksi pa imamo pogosto opravka s povsem pozitivnimi preslikavami, ki so definirane zgolj na 
podprostorih $C^*$-algebre, ki jim rečemo operatorski sistemi. V tem poglavju obravnavamo nekatere lastnosti 
pozitivnih in povsem pozitivnih preslikav na operatorskih sistemih. 

\begin{definicija}
    Naj bo $\mathfrak{A}$ enotska $C^*$-algebra. Njenemu podprostoru $\mathcal{M} \subseteq \mathfrak{A}$,
    ki vsebuje enico, pravimo \emph{operatorski prostor}. Operatorskemu prostoru $\mathcal{S} \subseteq \mathfrak{A}$,
    za katerega velja $\mathcal{S} = \mathcal{S}^*$, pravimo \emph{operatorski sistem}.
\end{definicija}

\begin{opomba}
    V splošni literaturi (glej na primer \cite{paulsen2002}) so operatorski sistemi običajno definirani zgolj kot podprostori $C^*$-algebre,
    torej brez zahteve, da vsebujejo enico. V tem magistrskem delu bomo obravnavali zgolj enotske operatorske prostore,
    zato dodamo to kot del definicije.
\end{opomba}

\begin{zgled}
    Če je $\mathcal{M}$ operatorski prostor, potem je $\mathcal{S} = \mathcal{M} + \mathcal{M}^*$ operatorski sistem.
\end{zgled}

Če je $\mathcal{M}$ (oziroma $\mathcal{S}$) operatorski prostor (sistem) v $\mathfrak{A}$,
potem je $M_n (\mathcal{M})$ operatorski prostor (sistem) v $M_n (\mathfrak{A})$ za vse $n \in \N$.
Pozitivni elementi v operatorskem prostoru so $\mathcal{M}_+ = \mathcal{M} \cap \mathfrak{A}_+$.
Na enak način kot za $C^*$-algebre podamo tudi definicijo pozitivne, $n$-pozitivne ter povsem pozitivne preslikave
$\varphi: \mathcal{M} \to \mathfrak{B}$ na operatorskem prostoru $\mathcal{M}$.

\begin{zgled}
    Če je $\varphi: \mathfrak{A} \to \mathfrak{B}$ povsem pozitivna preslikava
    in $\mathcal{M} \subseteq \mathfrak{A}$ poljuben operatorski prostor,
    potem je tudi $\varphi|_\mathcal{S}: \mathcal{S} \to \mathfrak{B}$ povsem pozitivna preslikava.
\end{zgled}

\begin{zgled}
    Če je $\mathcal{S} \subseteq \mathfrak{A}$ operatorski sistem in $\varphi: \mathcal{S} \to \C$
    pozitiven funkcional, potem je $\varphi$ povsem pozitivna preslikava. Dokaz poteka popolnoma enako kot v zgledu
    \ref{zgl:1.3}.
\end{zgled}

Operatorski sistemi imajo pomembno vlogo v teoriji se v teoriji povsem pozitivnih preslikav, saj je vsak njihov
element linearna kombinacija pozitivnih elementov.
Res, vsak element v operatorskemu sistemu $a \in \mathcal{S}$ je linearna kombinacija dveh sebiadjungiranih elementov 
$$a = \underbrace{\frac{a + a^*}{2}}_{\real a} + i \cdot \underbrace{\frac{a - a^*}{2i}}_{\imag a},$$
ki prav tako ležita v $\mathcal{S}$.
Poleg tega lahko vsak sebiadjungiran element zapišemo kot razliko dveh pozitivnih elementov 
$$a = \frac{1}{2} (\| a\| + a) - \frac{1}{2} (\| a\| - a),$$
ki ležita v $\mathcal{S}$. 

Operatorski sistemi so zato `naraven' objekt, na katerem so morfizmi povsem pozitivne preslikave.
Z oznako $\osys$ označimo kategorijo operatorskih sistemov, v kateri so morfizmi povsem pozitivne preslikave,
v $\osys_1$ pa enotske povsem pozitivne preslikave.

\begin{zgled}
    Za razliko od operatorskih sistemov lahko operatorski prostori vsebujejo le trivialne sebiadjungirane elemente, torej $\R \cdot 1$.
    Oglejmo si operatorski prostor $\mathcal{M} := \C [z] \subseteq C(\T)$.
    Najprej dokažimo, da je $\{z^k\}_{k \in \Z} \subseteq C(\T)$ linearno neodvisna množica.
    Denimo, da obstajajo cela števila $k_1 \leq k_2 \leq \cdots \leq k_l$ in skalarji $\alpha_1, \alpha_2, \dots, \alpha_l \in \C$,
    da velja $\sum_{j = 1} ^l \alpha_j z^{k_j} = 0$ na $\T$. Potem je tudi polinom 
    $\sum_{j = 1} ^l \alpha_j z^{k_j - k_1}$ ničelen na $\T$. Po osnovnem izreku algebre 
    ima vsak neničelen kompleksen polinom končno število ničel, zato mora biti $$\alpha_1 = \alpha_2 = \cdots = \alpha_l = 0.$$
    Sedaj lahko dokažemo, da so edini sebiadjungirani elementi v $\C[z]$ trivialni.
    Vzemimo poljuben sebiadjungiran polinom $p = \sum_{j = 0} ^l \alpha_j z^j$. Potem je
    \begin{align*}
        \sum_{j = 0} ^l \alpha_j z^j = \overline{\sum_{j = 0} ^l \alpha_j z^j} = \sum_{j = 0} ^l \overline{\alpha_j} \overline{z}^j = \sum_{j = 0} ^l \overline{\alpha_j} z^{-j}.
    \end{align*} 
    Od tod sledi $\alpha_0 \in \R$ in $\alpha_j = 0$ za $j \geq 1$.
\end{zgled}

Kot nam pove zgornji zgled, imajo operatorski sistemi lahko premalo 
pozitivnih elementov, da bi povsem pozitivne preslikave na njih bile zanimive.
Kljub temu pa obstaja vrsta preslikav
na operatorskem prostoru $\mathcal{M}$, ki jih lahko razširimo do povsem pozitivnih preslikav na operatorskem sistemu $\mathcal{M} + \mathcal{M}^*$.
Takšne preslikave bodo potem naravni morfizem na operatorskih prostorih.
V nadaljevanju tega razdelka bomo opisali to konstrukcijo.

\begin{lema}
    Naj bo $\mathcal{S}$ operatorski sistem in $\mathfrak{B}$ $C^*$-algebra.
    Če je $\varphi: \mathfrak{S} \to \mathfrak{B}$ pozitivna preslikava, potem je \emph{$*$-linearna}, torej
    $$\varphi(a^*) = \varphi(a) ^*,\quad \forall a \in \mathcal{S}.$$
\end{lema}

\begin{dokaz}
    Ta trditev očitno velja za $a \in \mathcal{S}_+$.
    Sedaj predpostavimo, da je $a \in \mathcal{S}$ sebiadjungiran (torej $a = a^*$). 
    Vemo, da lahko vsak tak $a$ zapišemo kot $a = \frac{1}{2} (\| a\| + a) - \frac{1}{2}(\| a\| - a)$.
    Sedaj dobimo
    \begin{align*}
        \varphi (a)^* &= \varphi \left(\frac{1}{2} \left(\| a \| + a \right) - \frac{1}{2} \left(\| a \| - a \right)\right)^*\\
        &= \varphi \left(\frac{1}{2} (\| a \| + a)\right)^* - \varphi \left(\frac{1}{2} \| a\| - a\right)^*\\
        &= \varphi \left(\frac{1}{2} (\| a \| + a)\right) - \varphi \left(\frac{1}{2} \| a\| - a \right)\\
        &= \varphi \left(\frac{1}{2} \left(\| a \| + a \right) - \frac{1}{2} \left(\| a \| - a \right)\right)\\
        &= \varphi(a) = \varphi(a^*).
    \end{align*}
    Tako za povsem splošen $a \in \mathcal{S}$ dobimo
    \begin{align*}
        \varphi(a)^* &= \varphi(\real a + i \cdot \imag a)^*\\
        &= \varphi(\real a)^* - i \varphi(\imag a)^*\\
        &= \varphi(\real a) - i \varphi (\imag a)\\
        &= \varphi(\real a - i \cdot \imag a) = \varphi(a^*). \qedhere
    \end{align*}
\end{dokaz}

\begin{lema}\label{lem:1.1}
    Naj bo $\mathcal{S}$ operatorski sistem in $\varphi: \mathcal{S} \to \C$ funkcional.
    Potem je $\varphi$ pozitiven natanko tedaj, ko je $\| \varphi\| = \varphi(1)$.
\end{lema}

\begin{dokaz}
    Ker za vsak $a \in \mathcal{S}_+$ velja
    $$-\| a\| \leq a\leq \| a\|,$$
    sledi 
    $$\varphi(a) \leq \varphi(\|a\|),\quad \forall a \in \mathcal{S}.$$
    Za $a, b \in \mathcal{S}$ definiramo seskvilinearno formo $\langle a, b\rangle := \varphi(b^* a)$,
    od koder po Cauchy-Schwarzevi neenakosti sledi
    \begin{equation}
        |\varphi(b^* a)|^2 \leq \varphi (b^* b) \cdot \varphi(a^* a). \label{eq:1.2}
    \end{equation}
    Če vstavimo $b = 1$, dobimo 
    \begin{align*}
        |\varphi(a)|^2 &\leq \varphi(a^* a) \cdot \varphi(1) \leq \varphi(\| a^* a\|) \cdot \varphi(1)\\
        & = \| a^* a\| \cdot \varphi(1)^2 = \| a\|^2 \cdot \varphi(1)^2.
    \end{align*}
    Če enačbo korenimo, dobimo $\| \varphi\| \leq \varphi(1)$.
    Ker je $\varphi(1) \leq \| \varphi\|$, smo implikacijo v desno dokazali.
    Za implikacijo v levo vzamemo
    $a \in \mathcal{S}_+$ in $\varphi (a) := \alpha + i \beta.$
    Za vsak $t \in \R$ je
  \begin{align*}
    \alpha^2 + (\beta + t \|\varphi\|)^2 &= |\alpha + i (\beta + t \varphi(1))|^2\\
    &= |\varphi (a + it)|^2\\
    &\leq \| a + it\|^2 \cdot \|\varphi\|^2 \\
    &= \left( \|a\|^2 + t^2 \right) \|\varphi\|^2. 
  \end{align*}
    Od tod sledi $2 \beta t \|\varphi\| \leq \|a\|^2 \cdot \|\varphi\|^2$.
    Ker je bil $t \in \R$ poljuben, imamo $\beta = 0$ in $\varphi(a) = \alpha \in \R$.
    Tako dobimo
    \begin{align*}
        \| a\| \cdot \|\varphi\| - \varphi(a) &= \varphi(\|a\| - a)\\
        &\leq \|\|a\| - a\| \cdot \|\varphi\|\\
        &\leq \| a\| \cdot \| \varphi \|,
    \end{align*}
    torej $\varphi(a) \geq 0$.
\end{dokaz}

Za poljubno pozitivno preslikavo $\varphi: \mathcal{S} \to \mathfrak{B}$ imamo na voljo slabšo oceno.

\begin{trditev}\label{trd:1.1}
    Naj bo $\mathcal{S}$ operatorski sistem in $\mathfrak{B}$ $C^*$-algebra.
    Če je $\varphi: \mathcal{S} \to \mathfrak{B}$ pozitivna preslikava, potem je omejena in velja 
    $\| \varphi\| \leq 2 \cdot \varphi(1)$.
\end{trditev}

\begin{dokaz}
    Za sebiadjungiran $a \in \mathcal{S}$ velja
    $$-\| a\| \leq a \leq \| a\|$$
    in posledično $$-\varphi(\| a\|) \leq \varphi(a) \leq \varphi (\| a\|).$$
    Od tod sledi $| \varphi (a)| \leq \| a \| \cdot \varphi (1)$.
    Sedaj za poljuben $a \in \mathcal{S}$ dobimo 
    \begin{align*}
        \| \varphi(a)\| &= \| \varphi \left(\real a + i \cdot \imag a\right) \| \\
        &\leq \| \varphi (\real a)\| + \| \varphi (\imag a)\| \\
        &\leq (\| \real a\| + \| \imag a \|)\cdot \|\varphi (1)\| \\
        &\leq 2 \| a \| \cdot \| \varphi(1)\|. \qedhere
    \end{align*}
\end{dokaz}

Naslednji zgled (\cite[Appendix A.2, str. 221]{arveson1969}) pokaže, da je konstanta $2$ v oceni iz trditve \ref{trd:1.1} tudi dosežena.
\begin{zgled}[Arveson]\label{zgl:1.6}
    Naj bo $\mathcal{S} \subseteq C(\mathbb{T})$ operatorski sistem, ki ga razpenjajo funkcije $1, z, \overline{z}$.
    Definirajmo preslikavo
    $$\varphi: \mathcal{S} \to M_2 (\C), \quad \varphi (a + bz + c \overline{z}) = \begin{bmatrix}
        a & 2b\\
        2c & a
    \end{bmatrix}.$$
    Dokazujemo, da je $\varphi$ pozitivna preslikava.
    
    Trdimo, da velja 
    $$a + bz + c \overline{z} \geq 0 \quad \Leftrightarrow \quad \text{$b = \overline{c}$ in $a \geq 2|b|$}.$$
    Najprej implikacija v levo: če velja $b = \overline{c}$ in $a \geq 2|b|$, potem je 
    $$a + bz + c \overline{z} = a + 2\real{bz} \geq 0,$$
    saj $a \geq 2|b| \geq 2 \real{bz}$ za vse $z \in \mathbb{T}$.
    Nato še implikacija v desno: če je 
    $$\overline{a + bz + c \overline{z}} = a + bz + c \overline{z},$$
    sledi $a = \overline{a}$ in $c = \overline{b}$. Od tod dobimo
    $a + 2 \real (bz) \geq 0$ za vse $z \in \mathbb{T}$. Na krožnici obstaja tak $z \in \mathbb{T}$,
    da velja $bz = - |b|$, zato mora veljati tudi $a \geq 2 |b|$. 
    
    Če je $a + bz + \overline{bz} \geq 0$, potem se s 
    $\varphi$ preslika v matriko
    $$\begin{bmatrix}
        a & 2b\\
        2\overline{b} & a
    \end{bmatrix}.$$
    Slednja ima nenegativne diagonalne elemente in determinanto
    $$\det \begin{bmatrix}
        a & 2b\\
        2\overline{b} & a
    \end{bmatrix} = a^2 - 4 |b|^2 \geq 0,$$
    zato je pozitivna in $\varphi$ je pozitivna preslikava.
    
    Nazadnje opazimo, da velja
    $$\| \varphi \| \leq 2 \cdot \| \varphi(1) \| = 2 = \| \varphi(z)\| \leq \| \varphi\|,$$
    torej $\| \varphi \| = 2 \cdot \| \varphi(1)\|$.
\end{zgled}

\begin{opomba}
    V zgornjem zgledu sta dve težavi: kodomena $M_2 (\C)$ je nekomutativna (glej lemo \ref{lem:1.1})
    in operatorski sistem $\mathcal{S}$ `ne premore' dovolj pozitivnih elementov (glej trditev \ref{pos:russo-dye}).
    Zgled \ref{zgl:1.6} nam pove, da za posplošitev stanj na preslikave $\varphi: \mathfrak{A} \to \mathfrak{B}$
    ni dovolj predpostaviti zgolj pozitivnost, temveč potrebujemo povsem pozitivnost. 
\end{opomba}

Poglavje zaključimo s trditvijo, da je enotsko skrčitev na operatorskem prostoru $\mathcal{M}$
moč razširiti do pozitivne preslikave na operatorskem sistemu $\mathcal{M} + \mathcal{M}^*$.

\begin{trditev}
    Naj bo $\mathcal{S}$ operatorski sistem in $\mathfrak{B}$ $C^*$-algebra.
    Potem je vsaka enotska skrčitev $\varphi: \mathcal{S} \to \mathfrak{B}$ pozitivna.
\end{trditev}

\begin{dokaz}
    Po izreku GNS lahko (kot bomo v nadaljevanju pogosto storili) 
    brez škode za splošnost predpostavili, da je $\mathfrak{B} = \bh$ 
    za nek Hilbertov prostor $\mathcal{H}$.
    Fiksirajmo nek $x \in \mathcal{H}$, tako da velja $\| x\| = 1$.
    Oglejmo si funkcional
    $$\rho_x: \mathcal{S} \to \C,\quad \rho_x (a) = \langle \varphi(a) x, x \rangle.$$
    Po Cauchy-Schwarzevi neenakosti sledi
    $$|\rho_x (a)| \leq \| \varphi(a) x \| \cdot \| x\| \leq \| \varphi(a)\| \cdot \| x\|^2 = \|\varphi(a)\| \leq \| a\|,$$
    torej $\| \rho_x \| \leq 1$. Ker pa velja $\rho_x (1) = 1$, smo dokazali $\| \rho_x \| = 1 = \rho_x(1)$.
    Lema \ref{lem:1.1} nam zato pove, da je $\rho_x$ pozitiven funkcional.
    Sedaj vzemimo poljuben $a \in \mathcal{S}_+$. 
    Ker je
    $$\langle \varphi(a) x, x \rangle = \rho_x (a) \geq 0$$
    za vsak enotski vektor $x \in \mathcal{H}$, 
    mora $\varphi (a)$ biti pozitiven.
\end{dokaz}

\begin{trditev}
    Naj bo $\mathcal{M} \subseteq \mathfrak{A}$ operatorski prostor in $\mathfrak{B}$ $C^*$-algebra. 
    Potem lahko vsako enotsko skrčitev $\varphi: \mathcal{M} \to \mathfrak{B}$
    enolično razširimo do pozitivne preslikave na operatorskem sistemu $\mathcal{M} + \mathcal{M}^*$ s predpisom 
    $$\tilde{\varphi} : \mathcal{M} + \mathcal{M}^* \to \mathfrak{B},\quad \tilde{\varphi} (a + b^*) = \varphi(a) + \varphi (b)^*.$$
\end{trditev}

\begin{dokaz}
    Najprej dokažimo enoličnost. Če je $\psi: \mathcal{M} + \mathcal{M}^* \to \mathfrak{B}$ pozitivna razširitev
    $\varphi$, mora biti $*$-linearna. To pomeni, da za vsaka $a, b \in \mathcal{M}$ velja
    \begin{align*}
        \psi (a + b^*) &= \psi(a) + \psi(b^*) \\
        &= \psi(a) + \psi(b)^* \\
        &= \varphi(a) + \varphi(b)^* \\
        &= \tilde{\varphi} (a + b^*).
    \end{align*} 
    
    Nato dokažemo, da je $\tilde{\varphi}$ sploh dobro definirana preslikava.
    Vzemimo operatorski sistem 
    $$\mathcal{S} := \{a\ |\ \text{$a$ in $a^*$ sta elementa $\mathcal{M}$}\}.$$
    ($\mathcal{S}$ vsebuje enico, zato je neprazen).
    Ker je $\varphi: \mathcal{S} \to \mathfrak{B}$ enotska skrčitev, je pozitivna in v posebnem $*$-linearna.
    To pomeni, da za vsak $a \in \mathcal{S}$ velja
    $$\varphi (a^*) = \varphi(a)^*.$$
    Sedaj predpostavimo, da velja $a + b^* = c + d^* \in \mathcal{M} + \mathcal{M}^*$.
    Potem je $$a - c = (-b + d)^* \in \mathcal{S},$$
    torej 
    \begin{align*}
        \varphi(a) - \varphi(c) &= \varphi(a - c)\\
        &=\varphi(-b + d)^* \\
        &=-\varphi(b)^* + \varphi(d)^*.
    \end{align*}
    Od tod potem sledi
    \begin{align*}
        \tilde{\varphi}(a + b^*) &= \varphi(a) + \varphi(b)^*\\
        &= \varphi(c)  + \varphi(d)^*\\
        &= \tilde{\varphi}(c + d^*),
    \end{align*}
    torej je $\tilde{\varphi}$ dobro definirana.

    Nazadnje dokažimo, da je $\tilde{\varphi}$ pozitivna.
    Ponovno lahko predpostavimo, da je $\mathfrak{B} = \bh$ za nek Hilbertov prostor $\mathcal{H}$.
    Vzemimo nek $x \in \mathcal{H}$, za katerega velja $\| x\| = 1$. 
    Kot v prejšnji trditvi definiramo $$\rho_x : \mathcal{M} \to \C,\quad \rho_x (a) = \langle {\varphi} (a) x, x\rangle$$
    in 
    $$\tilde{\rho}_x : \mathcal{M} + \mathcal{M}^* \to \C,\quad \tilde{\rho}_x (a) = \langle \tilde{\varphi} (a) x, x\rangle.$$
    Dovolj je dokazati, da je $\tilde{\rho}_x$ pozitiven funkcional.
    Ker velja $ 1 = \rho_x (1) = \| \rho_x\|$, je $\rho_x$ po lemi \ref{lem:1.1} pozitiven funkcional. 
    Po Hahn-Banachovem izreku (\cite[Corollary III.6.5., str. 79]{conway1990}) ga lahko
    razširimo do funkcionala $\rho_x ': \mathcal{M} + \mathcal{M}^* \to \C$,
    za katerega prav tako velja $\| \rho_x '\| = 1 = \rho_x '(1)$.
    Slednji je zato pozitiven in posledično $*$-linearen.
    Od tod sledi, da za vsaka $a, b \in \mathcal{M}$ velja
    \begin{align*}
        \rho_x ' (a + b^*) = \rho_x (a) + \rho_x (b)^* = \tilde{\rho}_x (a + b^*).
    \end{align*}
    Torej je $\tilde{\rho}_x = \rho_x '$ pozitiven.
\end{dokaz}


\section{Povsem omejene preslikave}\label{sec:4}

Sedaj lahko postavimo zadostni pogoj za preslikavo $\varphi: \mathcal{M} \to \mathfrak{B}$,
da jo lahko razširimo do povsem pozitivne preslikave $\tilde{\varphi}: \mathcal{M} + \mathcal{M}^* \to \mathfrak{B}$.
Takšnim preslikavam pravimo povsem omejene preslikave. Ta razdelek je namenjen osnovnim lastnostim in zgledom povsem omejenih preslikav.

\begin{definicija}
    Naj bosta $\mathfrak{A}, \mathfrak{B}$ $C^*$-algebri in $\mathcal{M} \subseteq \mathfrak{A}$ operatorski prostor.
    Potem je preslikava $\varphi: \mathcal{M} \to \mathfrak{B}$:
    \begin{enumerate}
        \item \emph{povsem omejena}, če je $\| \varphi\|_{\po} := \sup_{n \in \N} \| \varphi^{(n)}\| < \infty$;
        \item \emph{povsem skrčitev}, če je $\| \varphi\|_{\po} \leq 1$;
        \item \emph{povsem izometrija}, če je $\varphi^{(n)}$ izometrija za vsak $n \in \N$. 
    \end{enumerate}
    Množico vseh povsem omejenih preslikav označujemo s $\PO (\mathcal{M}, \mathfrak{B})$.
\end{definicija}

\begin{opomba}
    Če je $n \geq m$, potem je $\| \varphi^{(n)} \| \leq \| \varphi^{(m)}\|$.
\end{opomba}
 
\begin{trditev}
    Množica $\PO (\mathcal{M}, \mathfrak{B})$ je Banachov prostor.
\end{trditev}

\begin{dokaz}
    Enostavno je videti, da je $\PO (\mathcal{M}, \mathfrak{B})$ vektorski prostor in $\| \cdot \|_{\po}$ res ustreza aksiomom za normo.
    Da dokažemo polnost, vzemimo poljubno Cauchyjevo zaporedje $\{\varphi_k\}_{k \in \N} \subseteq \PO (\mathcal{M}, \mathfrak{B})$.
    Ker je $\| \cdot \| \leq \| \cdot\|_{\po}$, je $\{\varphi_k\}$ Cauchyjevo tudi v navadni normi $\| \cdot\|$, torej $\varphi_k \to \psi$
    v navadni normi. Od tod sledi $\varphi_k ^{(n)} \to \psi^{(n)}$ za vsak $n \in \N$.
    Za poljuben $\varepsilon > 0$ obstaja tak $M \in \N$, da za vsaka $k, l \geq M$
    velja $\| \varphi_k - \varphi_l\|_{\po} < \varepsilon$, torej $\| \varphi_k ^{(n)} - \varphi_l ^{(n)}\|_{\po} < \varepsilon$ za vsak $n \in \N$.
    Če pošljemo $l \to \infty$, dobimo $\| \varphi_k ^{(n)} - \psi^{(n)} \| \leq \varepsilon$ za vsak $n \in \N$, od koder sledi 
    $\| \varphi_k - \psi\|_{\po} \leq \varepsilon$ za vse $k \geq M$.
    To pomeni, da $\varphi_k \to \psi$ tudi v normi $\| \cdot \|_{\po}$.    
\end{dokaz}

\begin{trditev}
    Naj bo $\mathcal{S}$ operatorski sistem in $\mathfrak{B}$ $C^*$-algebra.
    Potem je vsaka enotska povsem skrčitev $\varphi: \mathcal{S} \to \mathfrak{B}$ povsem pozitivna.
\end{trditev}

\begin{dokaz}
    Za vsak $n \in \N$ je $\varphi^{(n)}: M_n (\mathcal{S}) \to M_n (\mathfrak{B})$ enotska skrčitev,
    zato je pozitivna.
\end{dokaz}

\begin{posledica}\label{pos:1.2}
    Vsaka enotska povsem skrčitev $\varphi: \mathcal{M} \to \mathfrak{B}$ premore enolično določeno povsem pozitivno razširitev 
    $\tilde{\varphi}: \mathcal{M} + \mathcal{M}^* \to \mathfrak{B}$.
\end{posledica}

\begin{dokaz}
    Za poljuben $n \in \N$ je $\varphi^{(n)}: M_n (\mathcal{M}) \to M_n (\mathfrak{B})$ enotska skrčitev, 
    torej je $\widetilde{(\varphi^{(n)})}: M_n (\mathcal{M}) + M_n (\mathcal{M})^* \to M_n (\mathfrak{B})$ pozitivna.
    Ker je $M_n (\mathcal{M}) + M_n (\mathcal{M})^* = M_n (\mathcal{M} + \mathcal{M}^*)$ in preslikava
    $\widetilde{(\varphi^{(n)})}$ sovpada s $(\tilde{\varphi})^{(n)}: M_n (\mathcal{M} + \mathcal{M}^*) \to M_n (\mathfrak{B})$,
    je slednja preslikava prav tako pozitivna.
    \begin{comment}
    Vzemimo matriki $[a_{ij}]_{i, j} \in M_n (\mathcal{M}), [b_{ij} ^*]_{i, j} \in M_n (\mathcal{M})$. Potem je 
    \begin{align*}
        \widetilde{(\varphi^{(n)})} ([a_{ij}]_{i, j} + [b_{ij} ^*]_{i, j}) &= \varphi^{(n)} ([a_{ij}]_{i, j}) + \varphi^{(n)} ([b_{ji}]_{i, j})^*\\
        &= [\varphi(a_{ij})]_{i, j} + [\varphi(b_{ji})]_{i, j}^*\\
        &= [\varphi(a_{ij})]_{i, j} + [\varphi(b_{ij})^*]_{i, j}\\
        &= [\varphi(a_{ij})]_{i, j} + [\varphi(b_{ij}^*)]_{i, j}\\
        &= [\varphi(a_{ij} + b_{ij}^*)]_{i, j}\\
        &= [\tilde{\varphi} (a_{ij} + b_{ij}^*)]_{i,j}\\
        &= \tilde{\varphi}^{(n)} ([ a_{ij} + b_{ij}^*]_{i,j})\\
        &= \tilde{\varphi}^{(n)} ([ a_{ij}]_{i, j} + [b_{ij}^*]_{i,j}),
    \end{align*}
    torej je $\tilde{\varphi}^{(n)}$ pozitivna za vsak $n \in \N$.
    \end{comment}
\end{dokaz}

\begin{opomba}
    Z oznako $\os$ bomo označevali kategorijo operatorskih prostorov, v kateri so morfizmi povsem omejene preslikave, 
    v $\os_1$ pa enotske povsem omejene preslikave.
\end{opomba}

Oglejmo si nekatere zglede povsem omejenih preslikav.

\begin{zgled}
    Dokažimo, da je vsak omejen funkcional $\varphi: \mathcal{M} \to \C$ povsem omejen.
    Za poljubna enotska vektorja $\alpha, \beta \in \C^n$ in $A \in M_m (\mathcal{M})$ imamo
    \begin{align*}
        \left| \beta^* \varphi^{(n)} (A) \alpha \right| &= \left| \varphi(\beta^* A \alpha)\right|\\
        &\leq \| \varphi\| \cdot \| \beta^* A \alpha\|\\
        &\leq \| \varphi\| \cdot \| A\| \cdot \| \alpha \| \cdot \| \beta\|\\
        &= \| \varphi\| \cdot \| A\|.
    \end{align*}
    Če vzamemo supremum obeh strani po $\alpha, \beta$, dobimo 
    $$\| \varphi^{(n)} (A)\| \leq \| \varphi\| \cdot \| A\|$$
    in posledično $\| \varphi^{(n)} \| = \| \varphi\|= \| \varphi \|_{\po}$.
\end{zgled}

\begin{comment}
\begin{zgled}
    Iz zgleda \ref{zgl:1.3} že vemo, da je vsak pozitivni funkcional $\varphi: \mathcal{S} \to \C$ povsem pozitiven.
    Dokazujemo, da je povsem omejena, in sicer da velja $\| \varphi^{(n)}\| = \| \varphi\|$ za vsak $n \in \N$.
    Brez škode za splošnost predpostavimo, da je $\| \varphi \| = \varphi(1) = 1$.
    Dovolj je dokazati Cauchy-Schwarzevo neenakost
    \begin{equation}
        \varphi^{(n)} (A) ^* \varphi^{(n)}(A) \leq \varphi^{(n)} (A^* A) \label{eq:1.3}
    \end{equation}
    za poljubno matriko $A \in M_n (\mathcal{S})$. Od tod po pozitivnosti $\varphi^{(n)}$ sledi 
    $$\|\varphi^{(n)} (A)\|^2 = \| \varphi^{(n)} (A)^* \varphi^{(n)} (A) \| \leq \| \varphi^{(n)} (A^* A)\| \leq \| \varphi^{(n)} (1)\| \cdot \| A^* A\| = \| A\|^2,$$
    torej $\| \varphi^{(n)}\| = \| \varphi\| = 1$ za vsak $n \in \N$. Preostane le še, da dokažemo neenakost \eqref{eq:1.3}.
    Vzemimo poljuben vektor $\alpha = (\alpha_1, \dots, \alpha_n) \in \C^n$. Potem je po uporabi neenakosti \eqref{eq:1.2}
    \begin{align*}
        \langle \varphi^{(n)} (A)^* \varphi^{(n)} (A) \alpha, \alpha \rangle &= \langle \varphi^{(n)} (A) \alpha, \varphi^{(n)} (A) \alpha \rangle = \sum_{i = 1} ^n \left| \sum_{k = 1} ^n \varphi(a_{ik}) \alpha_k \right|^2\\
        &= \sum_{i = 1} ^n \left| \varphi\left(\sum_{k = 1} ^n a_{ik} \alpha_k \right) \right|^2 \leq \sum_{i = 1} ^n \varphi\left(\left(\sum_{k = 1} ^n a_{ik} \alpha_k \right)^* \left(\sum_{j = 1} ^n a_{ij} \alpha_j \right)\right) \\
        &= \sum_{i = 1} ^n \varphi\left(\sum_{k = 1} ^n a_{ik}^* \overline{\alpha_k} \sum_{j = 1} ^n a_{ij} \alpha_j \right) = \sum_{i = 1} ^n \sum_{k = 1} ^n \sum_{j = 1} ^n \alpha_j \varphi\left( a_{ik}^* a_{ij} \right)  \overline{\alpha_k} \\
        &= \sum_{k = 1} ^n \sum_{j = 1} ^n \sum_{i = 1} ^n  \alpha_j \varphi\left( a_{ik}^* a_{ij} \right)  \overline{\alpha_k} = \sum_{k = 1} ^n \sum_{j = 1} ^n  \alpha_j \varphi\left(\sum_{i = 1} ^n  a_{ik}^* a_{ij} \right)  \overline{\alpha_k} \\
        &= \langle \varphi^{(n)} (A^* A) \alpha, \alpha \rangle, 
    \end{align*}
    torej $\varphi^{(n)} (A)^* \varphi^{(n)} (A) \leq \varphi^{(n)} (A^* A)$.
\end{zgled}
\end{comment}

\begin{zgled}
    Vemo že, da je vsak enotski $*$-homomorfizem $\varphi: \mathfrak{A} \to \mathfrak{B}$ povsem pozitiven.
    Velja pa tudi, da je $\varphi$ skrčitev. Ker so preslikave $\varphi^{(n)}$ prav tako enotski $*$-homomorfizmi,
    so tudi skrčitve in je $\varphi$ povsem skrčitev. Če je $\varphi$ enotski $*$-monomorfizem, je izometrija in enako velja za vse njegove ampliacije $\varphi^{(n)}$.
    Torej je vsak enotski $*$-monomorfizem med $C^*$-algebrama povsem izometrija.
\end{zgled}

Opazimo, da so vse povsem pozitivne preslikave, ki jih do sedaj poznamo (pozitivni funkcionali, $*$-homomorfizmi itd.), tudi povsem omejene.
V preostanku tega poglavja dokažemo, da to velja tudi v splošnem. Argumentom, ki sledijo, pravimo tudi \emph{dilatacijski triki}.

\begin{lema}\label{lem:random1}
    Naj bo $\mathcal{H}$ Hilbertov prostor in $P, Q, A \in \bh$.
    Potem je $$\begin{bmatrix}
            P & A\\
             A^* & Q
        \end{bmatrix} \geq 0$$
        natanko tedaj, ko sta $P, Q$ pozitivna operatorja in 
        $|\langle Ax, y\rangle|^2 \leq \langle Py, y\rangle \cdot \langle Qx, x \rangle$
        za vse $x, y \in \mathcal{H}$.
        V posebnem velja
        $$\begin{bmatrix}
            1 & A\\
             A^* & Q
        \end{bmatrix} \geq 0$$
        natanko tedaj, ko $A^* A \leq Q$.
\end{lema}

\begin{dokaz}
    Predpostavimo, da velja $|\langle Ax, y\rangle|^2 \leq \langle Py, y\rangle \cdot \langle Qx, x \rangle$ za vse $x, y \in \mathcal{H}$.
        Potem sledi
        \begin{align*}
            \left\langle \begin{bmatrix}
                P & A\\
                 A^* & Q
            \end{bmatrix} \begin{bmatrix}
                y\\ x
            \end{bmatrix}, \begin{bmatrix}
                y \\ x
            \end{bmatrix} \right\rangle &= \langle Py, y \rangle + 2\real \langle Ax, y \rangle + \langle Qx, x\rangle\\
            &\geq \langle Py, y \rangle - 2 |\langle Ax, y \rangle| + \langle Qx, x\rangle\\
            &\geq \langle Py, y \rangle - 2 \langle P y, y \rangle^{\frac{1}{2}} \langle Q x, x \rangle^{\frac{1}{2}}  + \langle Qx, x\rangle\\
            &= (\langle Py, y \rangle^{\frac{1}{2}} - \langle Qx, x \rangle^{\frac{1}{2}})^2 \geq 0.
        \end{align*}
        Predpostavimo še obratno, torej da obstajata $x, y$, tako da
        $|\langle Ax, y\rangle|^2 > \langle Py, y\rangle \cdot \langle Qx, x \rangle$.
        Brez škode za splošnost smemo predpostaviti, da je $\langle P y, y \rangle = \langle Qx, x \rangle = 1$, 
        saj bi v nasprotnem primeru lahko pomnožili $x, y$ z ustreznima faktorjema.
        Prav tako lahko predpostavimo $\langle Ax, y \rangle \leq 0$,
        saj lahko v nasprotnem primeru preprosto zavrtimo $x$ za ustrezen kot $e^{i \theta}$. 
        Od tod sledi $$\langle Ax, y \rangle = - |\langle Ax, y \rangle | < -1$$
        in 
        \begin{align*}
            \left\langle \begin{bmatrix}
                P & A\\
                 A^* & Q
            \end{bmatrix} \begin{bmatrix}
                y\\ x
            \end{bmatrix}, \begin{bmatrix}
                y \\ x
            \end{bmatrix} \right\rangle &= \langle Py, y \rangle + 2\real \langle Ax, y \rangle + \langle Qx, x\rangle\\
            &< 1 - 2 + 1 = 0.
        \end{align*}
        Za dokaz druge trditve je dovolj pokazati,
        da velja $|\langle Ax, y\rangle|^2 \leq \langle y, y\rangle \cdot \langle Qx, x \rangle$ za vse $x, y$.
        Če vstavimo $y = Ax$, dobimo
        $$\langle Ax, Ax \rangle \leq \langle Qx, x \rangle$$ za vse $x$.
        Ker je leva stran enačbe enaka $\langle A^* A x, x \rangle$, dobimo $A^* A \leq Q$.
        Obratno, če velja $A^* A \leq Q$, potem po Cauchy-Schwarzove neenakosti sledi
        \begin{align*}
            |\langle Ax, y \rangle|^2 &\leq \| Ax \|^2 \cdot \| y\|^2 \\
            &\leq \langle A^* A x, x \rangle \cdot \langle y, y \rangle \\
            &\leq \langle Qx, x \rangle \cdot \langle y, y \rangle. \qedhere
        \end{align*}
\end{dokaz}

\begin{posledica}
    Če je $$\begin{bmatrix}
            P & A\\
             A^* & Q
        \end{bmatrix} \geq 0,$$
    potem velja $A^* A \leq \| P\| Q$ in $A^* A \leq \| Q\| P$.
\end{posledica}

\begin{dokaz}
    Ker je $0 \leq P \leq \| P\|$, imamo
        $$\begin{bmatrix}
            \| P\| & A\\
             A^* & Q
        \end{bmatrix} \geq \begin{bmatrix}
            P & A\\
             A^* & Q
        \end{bmatrix} \geq 0.$$
        Od tod sledi 
        $$\begin{bmatrix}
            1 & \frac{A}{\| P\|}\\
            \frac{A^*}{\| P\|} &\frac{Q}{\| P\|}
        \end{bmatrix} \geq 0$$ in 
        $$\frac{A^*}{\| P\|} \frac{A}{\| P\|} \leq \frac{Q}{\| P\|}$$
        po prejšnji lemi. Od tod sledi $A^* A \leq \| P\| Q$.
        Po enakem argumentu dokažemo tudi drugo neenakost.
\end{dokaz}

\begin{posledica}\label{pos:1.3}
    Naj bo $\mathcal{S}$ operatorski sistem in $\varphi: \mathcal{S} \to \mathfrak{B}$ $2$-pozitivna preslikava.
    Potem je $\| \varphi \| = \| \varphi(1)\|$.
\end{posledica}

\begin{dokaz}
    Ker je $\varphi$ $1$-pozitivna, je tudi $*$-linearna.
    Vzemimo $a \in \mathcal{S}$, tako da je $\| a\| \leq 1$.
    Potem je $$\begin{bmatrix}
        1 & a\\
        a^* & 1
    \end{bmatrix} \geq 0,$$
    torej 
    $$\begin{bmatrix}
        \varphi(1) & \varphi(a)\\
        \varphi(a)^* & \varphi(1)
    \end{bmatrix} \geq 0.$$
    Po prejšnji posledici dobimo $$\varphi (a) ^* \varphi(a) \leq \| \varphi(1)\| \varphi(1)$$
    in zato $\| \varphi(a)\|^2 \leq \| \varphi(1)\|^2$.
    S tem smo dokazali $\| \varphi \| \leq \| \varphi(1)\|$.
\end{dokaz}

\begin{posledica}[Cauchy-Schwarz za $2$-pozitivne preslikave]
    Naj bosta $\mathfrak{A}, \mathfrak{B}$ $C^*$-algebri in $\varphi: \mathfrak{A} \to \mathfrak{B}$ 
    $2$-pozitivna preslikava. Potem velja $$\varphi (a^* b) \varphi (b^* a) \leq  \| \varphi (a^* a ) \| \cdot \varphi (b^* b)$$
    in v posebnem
    $$\varphi (a)^* \varphi (a) \leq \| \varphi(1)\| \varphi(a^* a).$$
\end{posledica}

\begin{dokaz}
    Ker je
    $$\begin{bmatrix}
        a & b\\
        0 & 0
    \end{bmatrix}^* \begin{bmatrix}
        a & b\\
        0 & 0
    \end{bmatrix} = \begin{bmatrix}
        a^* a & a^* b\\
        b^* a & b^* b
    \end{bmatrix} \geq 0,$$
    dobimo
    $$\begin{bmatrix}
        \varphi(a^* a) & \varphi(a^* b)\\
        \varphi(a^* b)^* & \varphi(b^* b)
    \end{bmatrix} \geq 0$$
    in posledično $\varphi (a^* b) \varphi(b a^*) \leq \| \varphi(a^* a) \| \varphi(b^* b)$.
\end{dokaz}

Sedaj lahko zapišemo glavno trditev tega poglavja.

\begin{trditev}
    Naj bo $\mathcal{S} \subseteq \mathfrak{A}$ operatorski sistem in $\mathfrak{B}$ $C^*$-algebra.
    Če je $\varphi: \mathcal{S} \to \mathfrak{B}$ povsem pozitivna, potem je tudi povsem omejena in velja
    $\| \varphi \| = \| \varphi \|_{\po} = \| \varphi (1)\|$.
\end{trditev}

\begin{dokaz}
    Ker je $\| \varphi (1) \| \leq \| \varphi\| \leq \| \varphi\|_{\po}$, zadošča pokazati
    $\| \varphi\|_{\po} \leq \| \varphi(1)\|$. 
    Vzemimo $n \in \N$ in $a \in M_n (\mathcal{S})$,
    tako da velja $\| a\| = 1$. Ker je $\varphi$ povsem pozitivna preslikava, je
    $$ \varphi^{(2n)} : M_2 (M_n (\mathcal{S})) \to M_2 (M_n (\mathcal{S}))$$
    pozitivna. 
    Ker pa velja $ \varphi^{(2n)} = (\varphi^{(n)})^{(2)}$, je $\varphi^{(n)}$ $2$-pozitivna preslikava.
    Po posledici \ref{pos:1.3} nato sledi $\| \varphi^{(n)}\| = \| \varphi^{(n)} (I_n)\| = \| \varphi (1)\|$.
\end{dokaz}

\begin{opomba}
    V prejšnjem dokazu smo identificirali $M_{2n} \equiv M_2 (M_n(\mathcal{S}))$.
    To smemo narediti, saj imamo za poljubna $m, n \in \N$ $*$-izomorfizem
    $$\Phi_{m, n} : M_m (M_n (\mathfrak{A})) \to M_{mn} (\mathfrak{A}),\quad \begin{bmatrix}
        [A_{11}] & \cdots & [A_{1m}]\\
        \vdots & & \vdots\\
        [A_{m1}] & \cdots & [A_{mm}]
    \end{bmatrix} \mapsto \begin{bmatrix}
        A_{11} & \cdots & A_{1m}\\
        \vdots & & \vdots\\
        A_{m1} & \cdots & A_{mm}
    \end{bmatrix},$$
    ki ohranja pozitivne elemente. 
\end{opomba}

\begin{posledica}
    Če je $\mathcal{M} \subseteq \mathfrak{A}$ operatorski prostor in $\varphi: \mathcal{M} \to \mathfrak{B}$
    enotska povsem izometrija, potem je tudi $\tilde{\varphi}: \mathcal{M} + \mathcal{M}^* \to \mathfrak{B}$
    povsem izometrija.
\end{posledica}

\begin{dokaz}
    Ker je $\varphi: \mathcal{M} \to \mathfrak{B}$ enotska skrčitev, je tudi 
    $$\tilde{\varphi}: \mathcal{M} + \mathcal{M}^* \to \varphi(\mathcal{M}) + \varphi(\mathcal{M})^*$$
    povsem skrčitev. Po drugi strani pa je tudi 
    $\varphi^{-1}: \varphi(\mathcal{M}) \to \mathfrak{A}$ enotska skrčitev, zato je
    $$\widetilde{\varphi^{-1}}: \varphi(\mathcal{M}) + \varphi(\mathcal{M})^* \to \mathcal{M} + \mathcal{M}^*$$
    povsem skrčitev. Ker sta preslikavi $\tilde{\varphi}$ in $\widetilde{\varphi^{-1}}$
    druga drugi inverzni in obe povsem skrčitvi, sta obe povsem izometrični.
\end{dokaz}

\begin{posledica}
    Množica $\PP (\mathcal{S}, \mathfrak{B})$ je konveksen stožec v $\PO (\mathcal{S}, \mathfrak{B})$.
\end{posledica}

Naslednja trditev je posplošitev leme \ref{lem:1.1} za povsem pozitivne preslikave.

\begin{posledica}
    Naj bo $\varphi: \mathcal{S} \to \mathfrak{B}$ enotska preslikava.
    Potem je $\varphi$ povsem pozitivna natanko tedaj, ko je $\| \varphi \|_{\po} = \| \varphi\| = 1$.
\end{posledica}

\begin{dokaz}
    Sledi direktno iz posledice \ref{pos:1.2} in prejšnje trditve.
\end{dokaz}

\begin{comment}

\begin{lema}
    Če je $\mathcal{M}$ operatorski prostor in $\varphi: \mathcal{M} \to M_n (\C)$ omejena preslikava,
    potem je povsem omejena in velja $\| \varphi\|_{\po} \leq n \| \varphi\|$.
\end{lema}

\begin{dokaz}
    Najprej dokažimo trditev za primer $n = 1$. Tedaj imamo za poljubna vektorja $\alpha, \beta \in \C^n$ in $A \in M_m (\mathcal{M})$
    \begin{align*}
        \left| \langle \varphi^{(m)} (A) \alpha, \beta \rangle \right| &= \left| \varphi(\beta^\top A \alpha)\right|\\
        &\leq \| \varphi\| \cdot \| \beta^\top A \alpha\|\\
        &\leq \| \varphi\| \cdot \| A\| \cdot \| \alpha \| \cdot \| \beta\|.
    \end{align*}
    Če vzamemo supremum obeh strani neenačbe po enotskih vektorjih $\| \alpha\| = \| \beta \| = 1$, dobimo 
    $$\| \varphi^{(m)} (A)\| \leq \| \varphi\| \cdot \| A\|$$
    in posledično $\| \varphi \|_{\po} = \| \varphi\|$.
    Sedaj dokažimo trditev za splošen $n \in \N$.
    Definirajmo preslikave $$\varphi_{ij}: \mathcal{M} \to \C,\quad \varphi_{ij} (a) = \langle \varphi(a) e_j, e_i \rangle.$$
    Ker je $\varphi (a) = \sum_{i, j} ^n \varphi_{ij} (a) E_{ij}$, lahko zapišemo njene ampliacije kot  
    $$\varphi^{(m)} (A) = \sum_{i, j} ^n \varphi_{ij}^{(m)} (A) \otimes E_{ij} \in M_m (M_n (\C)).$$
    Opazimo tudi, da je $$\Phi: M_n (\C) \otimes M_m (\C) \to M_m (\C) \otimes M_n (\C),\quad A \otimes B \mapsto B \otimes A$$
    $*$-izomorfizem $C^*$-algeber, zato ohranja normo.
    Od tod sledi 
    \begin{align*}
        \| \varphi^{(m)} (A)\| &= \left\| \sum_{i, j} ^n \varphi_{ij}^{(m)} (A) \otimes E_{ij} \right\|\\
        &=  \left\|\sum_{i, j} ^n E_{ij} \otimes \varphi_{ij}^{(m)} (A) \right\|\\
        &= \left\| \begin{bmatrix}
            \varphi_{11} ^{(m)} (A) & \cdots & \varphi_{1n} ^{(m)} (A)\\
            \vdots & & \vdots\\
            \varphi_{n1} ^{(m)} (A) & \cdots & \varphi_{nn} ^{(m)} (A)\\
        \end{bmatrix} \right\|\\
        &= \left\| \begin{bmatrix}
            \left\| \varphi_{11} ^{(m)} (A) \right\| & \cdots & \left\| \varphi_{1n} ^{(m)} (A) \right\|\\
            \vdots & & \vdots\\
            \left\| \varphi_{n1} ^{(m)} (A) \right\| & \cdots & \left\| \varphi_{nn} ^{(m)} (A) \right\|\\
        \end{bmatrix} \right\|\\
        &\leq \left\| \begin{bmatrix}
            \left\| \varphi_{11} ^{(m)} (A) \right\| & \cdots & \left\| \varphi_{1n} ^{(m)} (A) \right\|\\
            \vdots & & \vdots\\
            \left\| \varphi_{n1} ^{(m)} (A) \right\| & \cdots & \left\| \varphi_{nn} ^{(m)} (A) \right\|\\
        \end{bmatrix} \right\|_{F}\\
        &\leq n \max_{ij} \| \varphi_{ij} ^{(m)} (A)\|\\
        &\leq n \cdot \| \varphi \| \cdot \| A\|,
    \end{align*}
    pri čemer $\| \cdot \|_F$ označuje Frobeniusovo matrično normo.
\end{dokaz}



\begin{zgled}
    Preslikava $$\varphi: M_2 (\C) \to M_2 (\C),\quad A \mapsto A^\top$$
    iz zgleda \ref{zgl:1.4} je po prejšnji lemi povsem omejena, toda ni povsem pozitivna.
\end{zgled}
\end{comment}

\section{Povsem pozitivne preslikave na $C(X)$}\label{sec:5}


V prejšnjih poglavjih smo že srečali primere povsem pozitivnih preslikav 
$\varphi: \mathfrak{A} \to \mathfrak{B}$, kot na primer pozitivne funkcionale, $*$-homomorfizme
in enotske povsem skrčitve. V tem poglavju se omejimo na primer, ko je vsaj ena izmer $C^*$-algeber $\mathfrak{A}, \mathfrak{B}$
komutativna. Po Gelfandovi transformaciji (\cite[Theorem VIII.2.1., str. 236]{conway1990}) je vsaka komutativna $C^*$-algebra
$*$-izomorfna $C(X)$, kjer je $X$ nek kompakten Hausdorffov topološki prostor (če ni izrecno napisano drugače,
bomo v tem magistrskem delu za $C(X)$ vedno implicitno predpostavili, da je $X$ kompakten ter Hausdorfov).
Stinespring (\cite{stinespring1955}) je dokazal, da za povsem pozitivne preslikave
$\varphi: C(X) \to \mathfrak{B}$ in $\varphi: \mathfrak{A} \to C(X)$ obstaja preprosta karakterizacija.

\begin{trditev}\label{trd:1.2}
    Preslikava $\varphi: C(X) \to \mathfrak{B}$ je povsem pozitivna natanko tedaj, ko je pozitivna. 
\end{trditev}

\begin{dokaz}
    Brez škode za splošnost predpostavimo, da je $\mathfrak{B} = \bh$ za nek Hilbertov prostor $\mathcal{H}$.
            Fiksirajmo neka vektorja $u, v\in \mathcal{H}$. Po Rieszovem reprezentacijskem izreku \cite[Theorem 6.18, str. 129]{rudin1987} obstaja (enolična) 
            regularna kompleksna Borelova mera $\mu_{u, v}$,
            tako da velja
            $$\langle \varphi(a) v, u \rangle = \int_X a \, d\mu_{u, v},\quad \forall a \in \mathfrak{A}.$$
            Vzemimo poljubnih $n$ vektorjev $u_1, \dots, u_n \in \mathcal{H}$ in definiramo pozitivno mero
            $\mu := \sum_{i, j} ^n |\mu_{u_i, u_j}|$. Označimo s
            $h_{u_i, u_j} = \frac{d \mu_{u_i, u_j}}{d \mu}$ Radon--Nikodymove odvode. 
            Dokazujemo, da je matrika $$H(x) = [h_{u_i, u_j} (x)]_{i, j}$$
            pozitivna $\mu$-skoraj povsod.    
            Vzemimo poljubne skalarje $\lambda_1, \dots, \lambda_n \in \C$ in definiramo $u := \sum_{i = 1} ^n \lambda_i u_i$.
            Tedaj je
            \begin{align*}
                \langle \varphi (a) u, u \rangle &= \left\langle \varphi (a) \left( \sum_{j = 1} ^n \lambda_j u_j \right), \sum_{i = 1} ^n \lambda_i u_i \right\rangle\\
                &= \sum_{i, j} ^n \overline{\lambda_i} \lambda_j \langle \varphi (a) u_j, u_i \rangle\\
                &= \sum_{i, j} ^n \overline{\lambda_i} \lambda_j \int_X a\, d\mu_{u_i, u_j}.
            \end{align*}
            Ker je $\varphi$ pozitivna preslikava, je
            $$\mu_{u, u} = \sum_{i, j} ^n \overline{\lambda_i} \lambda_j \mu_{u_i, u_j}$$
            pozitivna mera.
            Opazimo, da velja $\mu_{u, u} \ll \mu$, torej je njen Radon--Nikodymov odvod
            $$\frac{ d\mu_{u, u}}{d \mu} = \sum_{i, j} ^n \overline{\lambda_i} \lambda_j h_{u_i, u_j} \geq 0$$
            $\mu$-skoraj povsod.
            Dokazali smo, da je za vsak $\lambda \in \C^n$ funkcija $\langle H(x) \lambda, \lambda \rangle$ nenegativna $\mu$-skoraj povsod.
            Vzemimo nek gost števen podprostor $D \subseteq \C^n$, recimo $D = (\Q + i \Q)^n$.
            Za vsak $\lambda \in D$ označimo $$N_{\lambda} := \{x \in X\ |\ \langle H(x) \lambda, \lambda \rangle < 0\}$$
            ter $N := \bigcup_{\lambda \in D} N_\lambda$.
            Ker je $\mu(N_\lambda) = 0$ za vsak $\lambda \in D$ in je $D$ števna množica, je tudi $\mu(N) = 0$.
            Vzemimo poljuben $x \notin N$.
            Potem je zvezna funkcija $\lambda \mapsto \langle H(x) \lambda, \lambda \rangle$ nenegativna na $D \subseteq \C^n$, torej je nenegativna na celotnem
            prostoru $\C^n$. To pomeni, da je $H(x)$ nenegativna matrika povsod razen na $\mu$-ničelni množici $N$.
            Sedaj vzemimo poljubne $n \in \N$, $a_i, \dots, a_n \in \mathfrak{A}$ in $u_1, \dots, u_n \in \mathcal{H}$.
            Potem velja
            \begin{align*}
                \left\langle \varphi^{(n)} \left(\begin{bmatrix}
                    a_1^* a_1  & \cdots & a_1 ^* a_n\\
                    \vdots & & \vdots \\
                    a_n ^* a_1  & \cdots & a_n ^* a_n\\
                \end{bmatrix}\right) \begin{bmatrix}
                    u_1 \\ \vdots \\ u_n
                \end{bmatrix}, 
                \begin{bmatrix}
                    u_1 \\ \vdots \\ u_n
                \end{bmatrix}
                \right\rangle &= \sum_{i, j} ^n \langle \varphi (a_i ^* a_j) u_j, u_i \rangle \\
                &= \sum_{i, j} ^n \int_X \overline{a_i} (x) a_j (x)\, d\mu_{u_i, u_j}\\
                &= \sum_{i, j} ^n \int_X \overline{a_i} (x) a_j (x) h_{u_i, u_j}(x)\, d\mu\\
                &= \int_X \sum_{i, j} ^n \overline{a_i} (x) a_j (x) h_{u_i, u_j}(x)\, d\mu \geq 0\\ 
                &= \int_X \left\langle H(x) \begin{bmatrix}
                    a_1 (x) \\ \vdots \\ a_n (x)
                \end{bmatrix}, \begin{bmatrix}
                    a_1 (x) \\ \vdots \\ a_n (x)
                \end{bmatrix} \right\rangle\, d\mu\\ 
                &\geq 0.  \qedhere
            \end{align*}
\end{dokaz}

\begin{opomba}
    Če je $\mathcal{S} \subseteq C(X)$ operatorski sistem, potem pozitivna preslikava $\varphi: \mathcal{S} \to \mathfrak{B}$
    še ni nujno povsem pozitivna.
    Protiprimer smo že srečali v zgledu \ref{zgl:1.6}, kjer smo definirali 
    pozitivno preslikavo $\varphi: \mathcal{S} \to M_2 (\C)$ na operatorskem sistemu $\mathcal{S} = \lin\{1, z, \overline{z}\} \subseteq C(\T)$,
    tako da $\varphi$ ni povsem pozitivna.
    Težava je v tem, da razcep iz leme \ref{lem:1.2} ne velja za poljuben operatorski sistem, temveč zgolj za celotno $C^*$-algebro $\mathfrak{A}$.
\end{opomba}

Sedaj dokažimo podobno trditev še za primer, ko je kodomena $C(X)$.

\begin{trditev}
    Naj bo $\mathcal{S} \subseteq \mathfrak{A}$ poljuben operatorski sistem.
    Potem je preslikava $\varphi: \mathcal{S} \to C(X)$ povsem pozitivna natanko tedaj, ko je pozitivna.
\end{trditev}


\begin{dokaz}
    \begin{enumerate}
        Vzemimo poljubno pozitivno matriko $A = [a_{ij}] \in M_n(\mathcal{S})$ in $b_1, \dots, b_n \in C(X)$. Dokazati želimo, da velja
        $\sum_{i, j} ^n b_j^* \varphi(a_{ij}) b_i \geq 0$. Za poljuben $x \in X$ imamo
        \begin{align*}
            \left(\sum_{i, j} ^n b_j^* \varphi(a_{ij}) b_i \right) (x) &= \sum_{i, j} ^n \overline{b_j} (x) \cdot \varphi(a_{ij}) (x) \cdot b_i (x)\\
            &= \sum_{i, j} ^n  \varphi(\overline{b_j (x)} \cdot a_{ij} \cdot b_i (x)) (x) \\
            &=  \varphi \left(\sum_{i, j} ^n \overline{b_j (x)} \cdot a_{ij} \cdot b_i (x) \right) (x) \\
            &= \varphi \left( \begin{bmatrix}
                b_1 (x) \\ \vdots \\ b_n (x)
            \end{bmatrix}^* A \begin{bmatrix}
                b_1 (x) \\ \vdots \\ b_n (x)
            \end{bmatrix} \right) (x) \geq 0. \qedhere
        \end{align*}
    \end{enumerate}
\end{dokaz}

\begin{posledica}
    Če je $\varphi: C(X) \to \mathfrak{B}$ enotska skrčitev, potem je povsem skrčitev.
    Enako velja tudi za enotsko skrčitev $\varphi: \mathcal{M} \to C(X)$,
    kjer je $\mathcal{M} \subseteq \mathfrak{A}$ poljuben operatorski prostor.
\end{posledica}

\begin{dokaz}
    Ker je $\varphi: C(X) \to \mathfrak{B}$ enotska skrčitev, je pozitivna.
    Po trditvi \ref{trd:1.2} je tudi povsem pozitivna in zato povsem skrčitev. 
    Enak argument velja tudi za enotsko skrčitev $\varphi: \mathcal{M} \to C(X)$.
\end{dokaz}

Naslednjo posledico bomo večkrat uporabili kasneje.

\begin{posledica}\label{pos:aux}
    Naj bo $\mathcal{S} \subseteq C(X)$ operatorski sistem.
    Če je $\varphi: \mathcal{S} \to C(Y)$ enotska izometrija, potem je povsem izometrija.
\end{posledica}

\begin{dokaz}
    Ker je $\varphi: \mathcal{S} \to C(X)$ enotska skrčitev, je po prejšnji posledici povsem skrčitev.
    Po drugi strani pa je njen inverz $\varphi^{-1}: \varphi(\mathcal{S}) \to C(X)$ prav tako enotska skrčitev,
    zato je povsem skrčitev. Ker sta tako $\varphi$ in $\varphi^{-1}$ povsem skrčitvi, je $\varphi$ povsem izometrija.
\end{dokaz}

\section{Dilatacije}\label{sec:6}

Dilatacije so orodje, s katerimi operator na Hilbertovem prostoru razširimo do `pohlevnejšega' operatorja na večjem Hilbertovem prostoru.
To nam omogoča, da lahko nekatere trditve za splošne operatorje reduciramo na bolj poseben razred operatorjev. 
V tem razdelku dokažemo Stinespringov izrek (\cite{stinespring1955}), ki dilatacije poveže s povsem pozitivnimi preslikavami.


\begin{definicija}
    \emph{Dilatacija} operatorja $T \in \bh$ je operator $S \in \bk$, pri čemer 
    je $\mathcal{H}$ podprostor Hilbertovega prostora $\mathcal{K}$, $P_{\mathcal{H}}$ ortogonalna projekcija na $\mathcal{H}$
    in $T = P_{\mathcal{H}} S\big|_{\mathcal{H}}$.
    Pravimo tudi, da je $T$ \emph{kompresija} $S$.
    Ker je $\mathcal{K} = \mathcal{H} \oplus \mathcal{H}^\perp$, lahko operator $S$ zapišemo v matrični obliki
    $$S = \begin{bmatrix}
        T & *\\
        * & *
    \end{bmatrix}.$$    
\end{definicija}


Pogosto lahko za operator $T$ najdemo dilatacijo, ki ima določene dodatne lastnosti.
To ilustriramo v naslednjih dveh zgledih.


\begin{zgled}\label{zgl:1.1}
    Dokažimo, da ima vsaka izometrija unitarno dilatacijo.
    Naj bo $V \in \bh$ izometrija in $P := I - V V^*$
    projekcija na $(\im V)^\perp$. Zapišimo $\mathcal{K} := \mathcal{H} \oplus \mathcal{H}$ ter 
    definirajmo $U \in \bk$ v matrični obliki kot 
    $$U := \begin{bmatrix}
        V & P\\
        0 & V^*
    \end{bmatrix}.$$
    Očitno je $U$ dilatacija $V$. Dokazati moramo še, da je tudi unitarna. Ker velja
    \begin{align*}
        U^* U &= \begin{bmatrix}
            V^* & 0\\
            P & V
        \end{bmatrix} \begin{bmatrix}
            V & P\\
            0 & V^*
        \end{bmatrix}\\
        &= \begin{bmatrix}
            V^* V &  V^* P\\
            PV & P^2 + V V^*
        \end{bmatrix}\\
        &= \begin{bmatrix}
            I & 0\\
            0 & I
        \end{bmatrix}
    \end{align*}
    in 
    \begin{align*}
         U U^* &= \begin{bmatrix}
            V & P\\
            0 & V^*
        \end{bmatrix} \begin{bmatrix}
            V^* & 0\\
            P & V
        \end{bmatrix} \\
        &= \begin{bmatrix}
             P^2 + V V^* &  P V\\
            V^* P &V^* V 
        \end{bmatrix}\\
        &= \begin{bmatrix}
            I & 0\\
            0 & I
        \end{bmatrix},
    \end{align*}
    je $U$ res unitarna dilatacija $V$. Še več, v tem primeru je 
    $U$ \emph{potenčna dilatacija} $V$. To pomeni, da je za vsak $n \in \N$ operator
    $$U^n = \begin{bmatrix}
        V^n & *\\
        0 & (V^*) ^n
    \end{bmatrix}$$
    dilatacija za $V^n$.
\end{zgled}

\begin{zgled}\label{zgl:1.2}
    Dokažemo lahko tudi, da ima vsaka skrčitev izometrično dilatacijo.
    Vzemimo $T \in \bh$, tako da je $\| T\| \leq 1$. Definirajmo 
    $$D_T := (I - T^* T)^{\frac{1}{2}} \in \bh,$$tako da za vsak $x \in \mathcal{H}$ velja
    \begin{align*}
        \| T x \|^2 + \| D_T x\|^2 &= \langle Tx, Tx \rangle + \langle D_T x, D_T x \rangle\\
        &= \langle T^* T x, x\rangle + \langle D_T ^2 x, x\rangle\\
        &= \langle T^* T x, x\rangle + \langle (I - T^* T) x, x\rangle\\
        &= \langle x, x \rangle = \| x\|^2.
    \end{align*}
    Nato definirajmo Hilbertov prostor $\mathcal{K} := \bigoplus_{n \in \N} \mathcal{H}$, ki je prostor zaporedij
    $$\{(x_1, x_2, \dots)\ |\ x_n \in \mathcal{H},\ \sum_{n = 1} ^\infty \| x_n\|^2 < \infty\}$$
    s skalarnim produktom
    $$\langle (x_1, x_2, \dots), (y_1, y_2, \dots) \rangle := \sum_{n = 1} ^\infty \langle x_i, y_i \rangle.$$
    Nazadnje definiramo še operator 
    $$V: \mathcal{K} \to \mathcal{K},\quad (x_1, x_2,x_3, \dots) \mapsto (T x_1, D_T x_2, x_3, \dots),$$
    ki je izometrija, saj velja
    \begin{align*}
        \| V(x_1, x_2, x_3, \dots)\|^2 &= \| T x_1\|^2 + \| D_T x_2\|^2 + \sum_{n = 3} ^\infty \| x_n\|^2 \\
        &= \sum_{n = 1} ^\infty \| x_n\|^2 \\
        &= \| (x_1, x_2, x_3, \dots )\|^2.
    \end{align*}
    Če identificiramo $\mathcal{H}$ s $\mathcal{H} \oplus 0 \oplus 0 \oplus \cdots \subseteq \mathcal{K}$, potem velja 
    $$T^n = P_{\mathcal{H}} V^n \big|_{\mathcal{H}}$$
    za vse $n \in \N$.
\end{zgled}

Kot posledico zgledov \ref{zgl:1.1} in \ref{zgl:1.2} dobimo naslednji izrek.

\begin{izrek}[Sz.-Nagy]
    Naj bo $T \in \bh$ skrčitev. Potem obstaja Hilbertov prostor $\mathcal{K} \supseteq \mathcal{H}$ 
    in unitaren operator $U \in \bk$, tako da velja
    $$T^n = P_{\mathcal{H}} U^n \big|_{\mathcal{H}}$$
    za vse $n \in \N$.
\end{izrek}

Sz.-Nagyjev izrek nam omogoča, da nekatere trditve o skrčitvah reduciramo na trditve o unitarnih operatorjih.
Za ilustracijo tega argumenta dokažimo naslednjo trditev.

\begin{posledica}[von Neumannova neenakost]\label{pos:1.1}
    Naj bo $T \in \bh$ skrčitev $p \in \C [z]$ kompleksen polinom. Potem imamo
    $$\| p(T)\| \leq \sup \{|p(z)|\ |\ z \in \mathbb{T}\}.$$
\end{posledica}

\begin{dokaz}
    Naj bo $U$ potenčna dilatacija $T$, torej
    $T^n = P_{\mathcal{H}} U^n \big|_{\mathcal{H}}$
    za vse $n \in \N \cup\{0\}$. Potem velja $p(T) = P_{\mathcal{H}} p(U)\big|_{\mathcal{H}}$
    in posledično $\| p(T)\| \leq \| p(U)\|$. Ker je $U$ normalna, imamo po spektralnem izreku
    $$\| p(U) \| = \sup \{|p(\lambda)|\ |\ \lambda \in \sigma(U)\}.$$
    Nato upoštevamo še, da je $U$ unitarna in zanjo velja $\sigma(U) \subseteq \mathbb{T}$, torej
    $$\sup \{|p(\lambda)|\ |\ \lambda \in \sigma(U)\} \leq \sup \{|p(\lambda)|\ |\ \lambda \in \mathbb{T}\}.$$
    Vse to skupaj implicira
    \begin{equation*}
        \| p(T)\| \leq \| p(U)\| = \sup \{|p(\lambda)|\ |\ \lambda \in \sigma(U)\} \leq \sup \{|p(\lambda)|\ |\ \lambda \in \mathbb{T}\}. \qedhere
    \end{equation*} 
\end{dokaz}

\begin{zgled}\label{zgl:1.7}
    Naj bo $T \in \bh$ skrčitev in
    $$\mathcal{M} := \C [z] \subseteq \C (T)$$
    operatorski sistem. Potem je preslikava 
    $$\varphi: \mathcal{M} \to \bh,\quad p \mapsto p(T)$$
    enotska skrčitev po von Neumannovi neenakosti. To pomeni, da 
    jo lahko razširimo do pozitivne preslikave $$\tilde{\varphi}: \mathcal{M} + \mathcal{M}^* \to \bh,\quad p + q^* \mapsto p(T) + q(T)^*.$$
    Po Stone-Weierstrassovem izreku je
    \begin{equation*}
        \mathcal{M} + \mathcal{M}^* = \{p (z) + \overline{q ({z})}\ |\ \text{$p, q$ polinoma}\}. 
    \end{equation*}
    gosta podmnožica v $C(\T)$, zato lahko $\tilde{\varphi}$ razširimo do pozitivne preslikave $\psi: C(\T) \to \bh$.
    Po trditvi \ref{trd:1.2} sledi, da je tudi povsem pozitivna.
\end{zgled}

V prejšnjem zgledu smo dokazali, da za vsako skrčitev $T \in \bh$
obstaja enolična povsem pozitivna preslikava $\psi: C(\T) \to \bh$,
tako da velja $\psi (p) = p(T)$ za poljuben polinom $p \in C(\T)$.
Takoj lahko navedemo dve posledici tega dejstva.

\begin{posledica}[Russo-Dye]\label{pos:russo-dye}
    Če sta $\mathfrak{A}, \mathfrak{B}$ $C^*$-algebri in 
    $\varphi: \mathfrak{A} \to \mathfrak{B}$ pozitivna preslikava, 
    potem je $\| \varphi \| = \| \varphi(1)\|$.
\end{posledica}

\begin{dokaz}
    Brez škode za splošnost predpostavimo, da je $\mathfrak{A} \subseteq \bh$.
    Vzemimo poljuben $a \in \mathfrak{A}$ z $\| a \|\leq 1$. 
    Po prejšnji opombi obstaja povsem pozitivna preslikava $\psi: C(\T) \to \bh$,
    tako da je $\psi(z) = a$. Po konstrukciji je 
    $$\psi (C(\T)) = \overline{\{p(a) |\ p \in \C[z]\}} \subseteq \mathfrak{A}.$$ 
    Posledično je preslikava
    $$\varphi \circ \psi: C(\T) \to \mathfrak{B}$$ pozitivna, 
    zato je po trditvi \ref{trd:1.2} povsem pozitivna 
    Od tod sledi 
    $$\| \varphi(a)\| = \| (\varphi \circ \psi) (z)\| \leq \| (\varphi \circ \psi) (1) \| = \| \varphi (1)\|$$
    in posledično $\| \varphi \| = \| \varphi(1)\|$.
\end{dokaz}

\begin{posledica}[Matrična von Neumannova neenakost]
    Naj bo $T \in \bh$ skrčitev in $p = [p_{ij}]_{i, j} \in M_n(\C [z])$ kompleksna matrični polinom. 
    Potem velja
    $$\| p(T)\| = \| [p_{ij} (T)]_{i, j}\| \leq \sup \{\|[p_{ij}(z)]_{i, j}\|\ |\ z \in \mathbb{T}\}.$$
\end{posledica}

\begin{proof}
    Ker je $\psi: C(\T) \to \bh$ enotska povsem pozitivna preslikava, velja $\| \psi \|_{\po} = 1$.
    V posebnem velja $\| \psi^{(n)} \| =1$ za vsak $n \in \N$. Potem za $[p_{ij}]_{i, j} \in M_n (\C [z])$ velja
    \begin{equation*}
        \| [p_{ij} (T)]_{i, j}\| = \| \psi^{(n)} ([p_{ij} ]_{i, j}) \| \leq \| [p_{ij} ]_{i, j} \|. \qedhere
    \end{equation*}
\end{proof}

Vrnimo se nazaj k dilatacijam.
Na tej točki še ni popolnoma očitno, kaj imajo te skupnega s povsem pozitivnimi preslikavami, na kar smo namignili 
v začetku tega poglavja. Vemo pa, da je za povsem pozitivno preslikavo $\psi: \mathfrak{A} \to \bh$ in $V \in \mathcal{B}(\mathcal{K}, \mathcal{H})$ tudi 
\begin{equation}
    \mathfrak{A} \to  \bk,\quad A \mapsto V^* \psi (a) V \label{eq:1.1}
\end{equation}
povsem pozitivna preslikava. To pomeni, da sta dilatacija in kompresija povsem pozitivni preslikavi na $\bh$.

\begin{zgled}
    Denimo, da sta $\mathcal{H} \subseteq \mathcal{K}$ Hilbertova prostora.
    Potem sta po \eqref{eq:1.1} tudi preslikavi
    $$
        \bk \to \bh,
        \quad
        A \mapsto P_{\mathcal{H}}A|_{\mathcal{H}},
    $$
    in
    $$
        \bh \to \bk,
        \quad
        A \mapsto
        \begin{bmatrix}
            A & 0\\
            0 & 0
        \end{bmatrix}
    $$
    povsem pozitivni. 
    
    Ker je kompozitum povsem pozitivnih preslikav
    ponovno povsem pozitiven, je za vsako povsem pozitivno preslikavo
    $\varphi : \mathfrak{A} \to \bh$
    tudi preslikava
    $$
        \psi : \mathfrak{A} \to \bk,
        \quad
        a \mapsto
        \begin{bmatrix}
            \varphi(a) & 0\\
            0 & 0
        \end{bmatrix}
    $$
    povsem pozitivna.

    Prav tako je za vsako povsem pozitivno preslikavo
    $\psi : \mathfrak{A} \to \bk$
    tudi preslikava
    $$
        \varphi : \mathfrak{A} \to \bh,
        \quad
        \varphi(a)=P_{\mathcal{H}}\psi(a)|_{\mathcal{H}}
    $$
    povsem pozitivna. 
\end{zgled}

Stinespringov izrek nam pove, da je 
vsaka povsem pozitivna preslikava 
$\varphi : \mathcal{A} \to \bh$
oblike \eqref{eq:1.1}. 
Pravzaprav velja še več.

\begin{izrek}[Stinespring]\label{izr:1.1}
    Naj bo $\mathfrak{A}$ $C^*$-algebra in $\varphi: \mathfrak{A} \to \bh$ povsem pozitivna preslikava.
    Potem obstaja Hilbertov prostor $\mathcal{K}$, omejen operator $V \in \mathcal{B} (\mathcal{H}, \mathcal{K})$
    in $*$-homomorfizem 
    $\pi: \mathfrak{A} \to \bk$, tako da je
    $$\varphi(a) = V^* \pi(a) V,\quad \forall a \in \mathfrak{A}.$$
    Če je $\varphi$ enotska preslikava, potem je $V$ izometrija.
\end{izrek}


Za trenutek se ustavimo na primeru, ko je $\varphi: \mathfrak{A} \to \bh$ enotska povsem pozitivna preslikava.
Tedaj je $V$ izometrija, zato lahko
Hilbertov prostor $\mathcal{H}$ identificiramo s podprostorom
$V\mathcal{H} \leq \mathcal{K}$.
S to identifikacijo operator $V^*$ sovpada s projekcijo
$P_{\mathcal{H}}$ iz $\mathcal{K}$ na $\mathcal{H}$, zato velja
$$
    \varphi(a)=P_{\mathcal{H}}\pi(a)|_{\mathcal{H}},\quad \forall a \in \mathfrak{A}.
$$
Stinespringov izrek nam torej pove, da lahko vsako enotsko povsem pozitivno preslikavo $\varphi: \mathfrak{A} \to \bh$ dilatiramo do
$*$-homomorfizma.

\begin{dokaz}[Dokaz izreka \ref{izr:1.1}]
    Oglejmo si algebraični tenzorski produkt vektorskih prostorov $\mathfrak{A} \otimes \mathcal{H}$.
    Na njem definiramo seskvilinearno formo 
    $$\langle a \otimes x, b \otimes y \rangle = \langle \varphi(b^* a) x, y\rangle,$$
    ki jo razširimo linearno. Za vsak $u = \sum_{j = 1} ^n a_j \otimes x_j \in \mathfrak{A} \otimes \mathcal{H}$ imamo
    \begin{align*}
        \langle u, u \rangle &= \sum_{i, j} ^n \langle \varphi(a_i ^* a_j)x_j,x_i \rangle\\
        &= \left\langle \varphi^{(n)} \left(\begin{bmatrix}
            a_1 \\ \vdots \\ a_n
        \end{bmatrix}^* \begin{bmatrix}
            a_1 \\ \vdots \\ a_n
        \end{bmatrix}\right) \begin{bmatrix}
           x_1 \\ \vdots \\x_n
        \end{bmatrix}, \begin{bmatrix}
           x_1 \\ \vdots \\x_n
        \end{bmatrix} \right\rangle \geq 0,
    \end{align*}
    kar pomeni, da je $\langle \cdot, \cdot \rangle$ polskalarni produkt na $\mathfrak{A} \otimes \mathcal{H}$. 

    Vsakemu elementu $a \in \mathfrak{A}$ lahko priredimo preslikavo
    $$\pi_0 (a): \mathfrak{A} \otimes \mathcal{H} \to \mathfrak{A} \otimes \mathcal{H},\quad \sum_{j = 1} ^n a_j \otimes x_j \mapsto \sum_{j = 1} ^n a a_j \otimes x_j.$$
    Sedaj imamo za $u = \sum_{j = 1} ^n a_j \otimes x_j$ in $v = \sum_{i = 1} ^m b_i \otimes y_i$
    \begin{align}
        \langle u, \pi_0 (a) v \rangle &= \langle \sum_{j = 1} ^n a_j \otimes x_j, \sum_{i = 1} ^m a b_i \otimes y_i \rangle \nonumber\\ 
        &= \sum_{i, j} \langle \varphi((a b_i)^* a_j)x_j, y_i \rangle \nonumber\\
        &= \sum_{i, j} \langle \varphi (b_i ^* a^* a_j ) x_j, y_i \rangle \nonumber\\
        &= \langle \pi_0 (a^*) u, v \rangle. \label{eq:1.4}
    \end{align}

    Definirajmo $$\mathcal{N} := \{u \in \mathfrak{A} \otimes \mathcal{H}\ |\ \langle u, u \rangle = 0\}.$$
    Po Cauchy-Schwarzevi neenakosti je $\mathcal{N}$ podprostor v $\mathfrak{A} \otimes \mathcal{H}$ in $\langle \cdot, \cdot \rangle$
    inducira skalarni produkt na $\quot{\mathfrak{A} \otimes \mathcal{H}}{\mathcal{N}}$. 
    Slednji prostor nato dopolnimo do Hilbertovega prostora $\mathcal{K}$.

    Vzemimo poljuben $u \in \mathfrak{A} \otimes \mathcal{H}$ in definirajmo $$f(a) := \langle \pi_0 (a) u, u \rangle.$$ 
    Potem je po \eqref{eq:1.4} $f$ pozitiven linearen funkcional na $\mathfrak{A}$, saj je
    \begin{align*}
        f(a^* a) &= \langle \pi_0 (a^* a) u, u \rangle\\
        &= \langle \pi_0 (a^*) \pi_0 (a) u, u\rangle \\
        &= \langle \pi_0 (a) u, \pi_0 (a) u\rangle \geq 0.
    \end{align*}
    Od tod sledi
    $$\langle \pi_0 (a) u, \pi_0 (a) u\rangle = f(a^* a) \leq \| a^* a\| \cdot f(1) = \| a\|^2 \cdot \langle u, u \rangle.$$
    Posledično je $\pi_0 (a) (\mathcal{N}) \subseteq \mathcal{N}$, torej je $\mathcal{N}$ invarianten podprostor za $\pi_0 (a)$ v $\mathfrak{A} \otimes \mathcal{H}$, zaradi česar 
    $\pi_0 (a)$ inducira omejen operator na $\quot{\mathfrak{A} \otimes \mathcal{H}}{\mathcal{N}}$,
    ki ga nato razširimo do omejenega operatorja na $\mathcal{K}$. 
    Zato $\pi_0$ inducira preslikavo (ki je hkrati $*$-homomorfizem) 
    $$\pi:\mathfrak{A} \to \bk,\quad \pi (a) (u + \mathcal{N}) = \pi_0 (a)u + \mathcal{N}.$$

    Nazadnje definiramo še $$V: \mathcal{H} \to \mathcal{K},\quad x \mapsto 1 \otimes x + \mathcal{N}.$$
    Ker je $$\| V x \|^2 = \langle 1 \otimes x, 1 \otimes x \rangle_{\mathcal{K}} = \langle \varphi (1) x, x \rangle_{\mathcal{H}} \leq \| \varphi (1)\| \cdot \| x \|^2,$$
    je $V$ omejen operator. Za poljubna $x, y \in \mathcal{H}$ in $a \in \mathfrak{A}$ imamo
    \begin{align*}
        \langle V x, a \otimes y + \mathcal{N} \rangle &= \langle 1 \otimes x + \mathcal{N}, a \otimes y + \mathcal{N}\rangle\\
        &= \langle \varphi (a^*) x, y \rangle\\
        &= \langle x, \varphi (a^*)^* y \rangle\\
        &= \langle x, \varphi(a) y \rangle,
    \end{align*}
    torej je $V^* (a \otimes y + \mathcal{N}) = \varphi(a) y$. Sedaj je
    $$V^* \pi(a) V y = V^* \pi (a) (1 \otimes y + \mathcal{N}) = V^* (a \otimes y + \mathcal{N}) = \varphi(a) y.$$

    Če je $\varphi$ enotska preslikava, potem za poljuben vektor $y \in \mathcal{H}$ velja
    $$V^* V y = V^* \pi(1) V y = \varphi(1) y = y,$$
    torej je $V^* V = I$ in je $V$ izometrija.
\end{dokaz}

\begin{opomba}
    Stinespringov izrek je posplošitev izreka GNS o upodobitvah. 
    Če namreč vzamemo $\mathcal{H} = \C$, potem je $\mathcal{B}(\C) = \C$ in $\mathfrak{A} \otimes \C = \mathfrak{A}$, tako da je izometrija 
    $V: \C \to \mathcal{K}$ določena z $V(1) = x$. Tako imamo 
    $$\varphi(a) = \varphi(a) (1) \cdot 1 = V^* \pi(a) V(1) \cdot 1 = \langle \pi (a) V(1), V(1) \rangle_{\mathcal{K}}$$
\end{opomba}

Trojici $(\pi, V, \mathcal{K})$ iz Stinespringovega izreka pravimo \emph{Stinespringova upodobitev}
za $\varphi$. Naravno se je vprašati, ali je Stinespringova upodobitev enolična. 
Če že imamo eno Stinespringovo upodobitev $(\pi, V, \mathcal{K})$ za $\varphi$, lahko definiramo nov prostor $\mathcal{K}_1 := \overline{\pi(\mathfrak{A}) V \mathcal{H}} \leq \mathcal{K}$, ki reducira $\pi (\mathfrak{A})$.
To nam omogoča, da definiramo novo upodobitev $\pi_1: \mathfrak{A} \to \mathcal{B}(\mathcal{K}_1)$, za katero prav tako velja $$\varphi(a) = V^* \pi_1 (a) V,\quad \forall a \in \mathfrak{A}$$
in $(\pi_1, V, \mathcal{K}_1)$ je še ena Stinespringova upodobitev za $\varphi$. To je motivacija za naslednjo definicijo.

\begin{definicija}
    Stinespringova upodobitev $(\pi, V, \mathcal{K})$ povsem pozitivne preslikave $\varphi: \mathfrak{A} \to \bh$ je \emph{minimalna}, če je $\overline{\pi(\mathfrak{A}) V \mathcal{H}} = \mathcal{K}$.
\end{definicija}

\begin{trditev}
    Naj bo $\mathfrak{A}$ $C^*$-algebra in $\varphi: \mathfrak{A} \to \bh$ povsem pozitivna preslikava. 
    Če sta $(\pi_1, V_1, \mathcal{K}_1), (\pi_2, V_2, \mathcal{K}_2)$ 
    minimalni Stinespringovi upodobitvi, potem obstaja unitaren operator $U: \mathcal{K}_1 \to \mathcal{K}_2$, tako da velja
    $U V_1 = V_2$ in $U \pi_1 U^* = \pi_2$.
\end{trditev}

\begin{dokaz}
    Definiramo preslikavo
    $$\widetilde{U} : \pi_1 (\mathfrak{A}) V_1 \mathcal{H} \to \pi_2 (\mathfrak{A}) V_2 \mathcal{H},\quad \sum_i \pi_1 (a_i) V_1 x_i \mapsto \sum_i \pi_2 (a_i) V_2 x_i.$$
    Takoj lahko opazimo, da je $\widetilde{U}$ surjekcija.

    Nadalje dokažemo, da je $\widetilde{U}$ izometrija:
    \begin{align*}
        \left\| \sum_i \pi_1 (a_i) V_1 x_i \right\|^2 &= \left\langle \sum_i \pi_1 (a_i) V_1 x_i, \sum_j \pi_1 (a_j) V_1 x_j \right\rangle\\
        &= \sum_i \sum_j \langle \pi_1 (a_i) V_1 x_i, \pi_1 (a_j) V_1 x_j\rangle\\
        &= \sum_i \sum_j \langle V_1^* \pi_1 ^* (a_j) \pi_1 (a_i) V_1 x_i,  x_j\rangle\\
        &= \sum_i \sum_j \langle V_1^* \pi_1 (a_j^* a_i) V_1 x_i,  x_j\rangle\\
        &= \sum_i \sum_j \langle \varphi(a_j^* a_i) x_i,  x_j\rangle\\
        &= \sum_i \sum_j \langle V_2^* \pi_2 (a_j^* a_i) V_2 x_i,  x_j\rangle\\
        &= \left\langle \sum_i \pi_2 (a_i) V_2 x_i, \sum_j \pi_2 (a_j) V_2 x_j \right\rangle\\
        &= \left\| \sum_i \pi_2 (a_i) V_2 x_i \right\|^2.
    \end{align*}
    Torej lahko razširimo $\widetilde{U}$ do izometrije $U: \mathcal{K}_1 \to \mathcal{K_2}$
    Ker je $U(\mathcal{K}_1)$ gosta v $\mathcal{K}_2$, je $U$ surjektivna.

    Enakost $U V_1 = V_2$ sledi direktno iz definicije $U$.
    Do dokažemo $U \pi_1 = \pi_2 U$, opazimo, da za vsak $a \in \mathfrak{A}$ velja 
    \begin{align*}
        \pi_2 (a) U \big|_{\pi_1 (\mathfrak{A}) V_1 \mathcal{H}}
        &=  \pi_2 (a) \widetilde{U} \big|_{\pi_1 (\mathfrak{A}) V_1 \mathcal{H}}\\
        &= \widetilde{U} \pi_1 (a) \big|_{\pi_1 (\mathfrak{A}) V_1 \mathcal{H}}\\
        &= U \pi_1 (a) \big|_{\pi_1 (\mathfrak{A}) V_1 \mathcal{H}},
    \end{align*}
    torej $\pi_2 (a) U = U \pi_1 (a)$. 
\end{dokaz}

\begin{opomba}
    V dokazu izreka \ref{izr:1.1} je množica
    $$\pi(\mathfrak{A}) V \mathcal{H} = \quot{\mathfrak{A} \otimes \mathcal{H}}{\mathcal{N}}$$
    gosta v svoji napolnitvi $\mathcal{K}$, torej v dokazu konstruiramo minimalno Stinespringovo upodobitev.
\end{opomba}


\begin{posledica}[Sz.-Nagyjeva minimalna unitarna dilatacija] \label{pos:1.4}
    Naj bo $T \in \bh$ skrčitev. Potem obstaja Hilbertov prostor $\mathcal{K} \supseteq \mathcal{H}$
    in unitaren operator $U \in \bk$, tako da je $\mathcal{K}$ najmanjši reducirajoč prostor za $U$, da velja 
    $$T^n = P_{\mathcal{H}} U^n \big|_{\mathcal{H}},\quad \forall n \in \N.$$

    Če je $(U', \mathcal{K}')$ še en takšen par,
    potem obstaja unitaren operator $V: \mathcal{K} \to \mathcal{K}'$, tako da je $V\big|_{\mathcal{H}} = \id_{\mathcal{H}}$
    in $V U V^* = U'$. 
\end{posledica}

\begin{dokaz}
    Stinespringov izrek uporabimo na pozitivni preslikavi $\psi: C(\T) \to \bh$ iz zgleda \ref{zgl:1.7},
    ki preslika poljuben polinom $p \in \C[z]$ v $p(T) \in \bh$. 
    Potem obstaja minimalna Stinespringova upodobitev $(\pi, V, \mathcal{K})$,
    pri čemer je $V$ izometrija, saj je $\psi$ enotska preslikava. 
    Posledično lahko identificiramo $\mathcal{H} \equiv V \mathcal{H}$.

    Sedaj za poljuben $n \in \N$ velja
    $$T^n = \psi(z^n) = P_{\mathcal{H}} \pi(z) ^n \big|_{\mathcal{H}} = P_{\mathcal{H}} U ^n \big|_{\mathcal{H}},$$
    kjer je $U := \pi(z) \in \bk$ unitaren operator.
    Po zveznem funkcijskem računu je $\pi(f) = f(U)$ za poljubno funkcijo $f \in C(\T)$.
    Ker je $(\pi, V, \mathcal{K})$ minimalna Stinespringova upodobitev, je $\overline{\pi(C(\T)) \mathcal{H}} = \mathcal{K}$.
    Ker pa je $\pi(C(\T)) = C^*(U)$, je to ekvivalentno $$\overline{\{U^n h\ |\ n \in \Z,\ h \in \mathcal{H}\}} = \mathcal{K}.$$
    Slednje pa je ekvivalentno temu, da je $\mathcal{K}$ najmanjši reducirajoč prostor za $U$.
    Zadnja trditev v izreku nato sledi direktno iz unitarne ekvivalence minimalnih Stinespringovih upodobitev.
\end{dokaz}


Razdelek zaključimo s posledico Stinespringovega izreka, ki jo bomo potrebovali kasneje.
Gre za `nekomutativno posplošitev' Radon-Nikodymovega izreka.
Najprej na množici $\PP(\mathfrak{A}, \mathfrak{B})$ definiramo delno urejenost  
$$\varphi \leq \psi \quad \Leftrightarrow \quad \psi - \varphi \in \PP(\mathfrak{A}, \mathfrak{B}).$$
Nadalje za poljubno podmnožico $C^*$-algebre $\mathcal{S} \subseteq \mathfrak{A}$ definiramo njen \emph{komutant} kot
$$\mathcal{S}' := \{a \in A\ |\ as = sa,\ \forall s \in \mathcal{S}\}.$$
Sedaj velja naslednje.

\begin{trditev} \label{trd:1.4}
    Naj bosta $\varphi, \psi \in \PP(\mathfrak{A},\bh)$ in
    $(\pi, V, \mathcal{K})$ minimalna Stinespringova upodobitev $\psi$.
    Potem je $\psi \geq \varphi$ natanko tedaj, ko obstaja enoličen $T \in \pi(\mathfrak{A})'$,
    tako da je
    $$\varphi(a) = V^* T \pi(a) V$$ in $0 \leq T \leq 1$.
\end{trditev}

\begin{dokaz}
    Denimo, da je $\varphi = V^* T \pi(a) V$, pri čemer velja $0 \leq T \leq 1$ in $T \in \pi(\mathfrak{A})'$.
    Ker je $1 - T \geq 0$, lahko definiramo $S := (1 - T)^{\frac{1}{2}}$. 
    Opazimo, da velja $1 - T \in \pi(\mathfrak{A})'$, zato imamo po zveznem funkcijskem računu $S \in \pi(\mathfrak{A})'$.
    Od tod pa sledi, da je $$\psi - \varphi = V^* (1 - T) \pi(a) V = V^* S^* \pi(a) S V$$
    povsem pozitivna.
    
    Za obratno implikacijo predpostavimo, da je $\psi - \varphi$ povsem pozitivna.
    Ker so vse minimalne Stinespringove dilatacije $\psi$ unitarno ekvivalentne, lahko brez škode za splošnost vzamemo trojico $(\pi, V, \mathcal{K})$,
    ki smo jo konstruirali v dokazu izreka \ref{izr:1.1}.
    V istem dokazu smo vpeljali polskalarna produkta 
    $$\langle a \otimes x, b \otimes y \rangle_{\varphi} := \langle \varphi(b^* a) x, y\rangle,\quad \langle a \otimes x, b \otimes y \rangle_{\psi} := \langle \psi(b^* a) x, y\rangle$$
    na prostoru $\mathfrak{A} \otimes \mathcal{H}$.
    Ker je $\psi - \varphi$ povsem pozitivna, sledi $$\langle u, u \rangle_{\varphi} \leq \langle u, u \rangle_{\psi}$$
    za poljuben $u \in \mathfrak{A} \otimes \mathcal{H}$. Po Cauchy-Schwarzevi neenakosti imamo za poljubna $u, v \in \mathcal{H}$
    \begin{equation}
        \langle u, v\rangle_{\varphi} \leq \langle u, u\rangle_{\varphi} \langle v, v \rangle_{\varphi} \leq \langle u, u\rangle_{\psi} \langle v, v \rangle_{\psi}. \label{eq:1.5}
    \end{equation} 
    Označimo $\mathcal{N} = \{u \in \mathfrak{A} \otimes \mathcal{H}\ |\ \langle u, u \rangle_{\psi}\}$, tako da $\langle \cdot, \cdot \rangle_{\psi}$ tvori skalarni produkt na prostoru $\quot{\mathfrak{A} \otimes \mathcal{H}}{\mathcal{N}}.$
    Slednjega nato napolnimo do Hilbertovega prostora $\mathcal{K}$.

    Po neenakosti \eqref{eq:1.5} $\langle \cdot, \cdot \rangle_\varphi$ inducira omejen polskalarni produkt na $\mathcal{K}$.
    Po Rieszovem izreku obstaja operator $T \in \bk$, tako da je
    $$\langle u, v \rangle_{\varphi} = \langle T u, v \rangle_{\psi},\quad \forall u, v \in \mathcal{K}.$$
    Potem za vsak $a \in \mathfrak{A}$ in $x, y \in \mathcal{H}$ velja
    \begin{align*}
        \langle (V^* T \pi(a) V) x, y\rangle &= \langle T \pi(a) V x, V y\rangle_{\psi}\\
        &= \langle T \pi(a) (1 \otimes x + \mathcal{N}), (1 \otimes y + \mathcal{N}) \rangle _{\psi}\\
        &= \langle T (a \otimes x + \mathcal{N}), (1 \otimes y + \mathcal{N}) \rangle _{\psi}\\
        &= \langle (a \otimes x + \mathcal{N}), (1 \otimes y + \mathcal{N}) \rangle _{\varphi}\\
        &= \langle \varphi (a) x, y \rangle,
    \end{align*}
    od koder sledi $\varphi(a) = V^* T \pi(a) V$. Prav tako za vsak $u \in \mathcal{K}$ velja
    $$0 \leq \langle T u, u \rangle_{\psi} = \langle u, u \rangle_{\varphi} \leq \langle u, u \rangle_{\psi},$$
    torej $0 \leq T \leq 1$. 
    
    Da dokažemo $T \in \pi(\mathfrak{A})'$, vzamemo poljubne $a, b, c \in \mathfrak{A}$ ter $x, y \in \mathcal{H}$.
    Potem \begin{align*}
        \langle T \pi(a) (b \otimes x + \mathcal{N}), (c \otimes y + \mathcal{N}) \rangle_\psi &= \langle T (ab \otimes x + \mathcal{N}), (c \otimes y + \mathcal{N}) \rangle_\psi\\
        &= \langle (ab \otimes x + \mathcal{N}), (c \otimes y + \mathcal{N}) \rangle_\varphi\\
        &= \langle \varphi (c^* ab)x, y\rangle\\
        &= \langle \varphi ((a^* c)^* b)x, y \rangle\\
        &= \langle (b \otimes x + \mathcal{N}), (a^* c \otimes y + \mathcal{N}) \rangle_\varphi\\
        &= \langle T (b \otimes x + \mathcal{N}), \pi(a^*) (c \otimes y + \mathcal{N}) \rangle_{\psi}\\
        &= \langle\pi(a) T  (b \otimes x + \mathcal{N}), (c \otimes y + \mathcal{N}) \rangle_\psi     
    \end{align*}
    torej $\langle T \pi(a) u, v \rangle_{\psi} = \langle \pi(a) T u, v \rangle_{\psi}$ za vse $u, v \in \mathcal{K}$. Od tod dobimo $T \pi(a) = \pi(a) T$.

    Za dokaz enoličnosti vzemimo nek $T' \in \pi(\mathfrak{A})'$, tako da je $0 \leq T' \leq 1$ in $\varphi = V^* T' \pi V$.
    Ponovno vzemimo poljubna $a, b \in \mathfrak{A}$ in $x, y \in \mathcal{H}$. Potem je 
    \begin{align*}
        \langle T (a \otimes x + \mathcal{N}), (b \otimes y + \mathcal{N}) \rangle _{\psi} &= \langle (V^* T \pi(b^* a) V) x, y\rangle\\
        &= \langle (V^* T' \pi(b^* a) V) x, y\rangle\\
        &= \langle T' (a \otimes x + \mathcal{N}), (b \otimes y + \mathcal{N}) \rangle _{\psi},
    \end{align*}
    zato je $\langle Tu, v \rangle_{\psi} = \langle T' u, v\rangle_{\psi}$ za poljubna $u, v \in \mathcal{K}$.
    Od tod dobimo $T' = T$.
\end{dokaz}



\begin{comment}


Nadaljujemo s posplošitvijo Radon-Nikodymovega izreka za povsem pozitivne preslikave.

\begin{definicija}
    Na množici $\PP(\mathfrak{A}, \mathfrak{B})$ lahko definiramo delno urejenost  
    $$\varphi \leq \psi \quad \Leftrightarrow \quad \psi - \varphi \in \PP(\mathfrak{A}, \mathfrak{B}).$$
\end{definicija}


Za zaključek podrazdelka dokažemo uporabno posledico Stinespringovega izreka, ki klasificira vse povsem pozitivne preslikave med 
prostori končno razsežnih matrik.

\begin{izrek}[Choi]\label{izr:1.2}
    Naj bo $\varphi: M_n (\C) \to M_m (\C)$ povsem pozitivna.
    Potem obstaja $r \leq m \cdot n$, tako da imamo $r$ kompleksnih matrik $V_1, \dots, V_r$ dimenzije $n \times m$,
    ki zadoščajo 
    \begin{equation}
        \varphi(A) = \sum_{k = 1} ^r V_k ^* A V_k,\quad \forall A \in M_n (\C). \label{eq:1.6}
    \end{equation}
\end{izrek}

Za dokaz potrebujemo naslednjo lemo.

\begin{lema}
    Naj bo $\mathfrak{A}$ $C^*$-algebra in $\pi: M_n (\mathfrak{A}) \to \bk$ upodobitev.
    Potem obstaja upodobitev $\rho: \mathfrak{A} \to \bh$, tako da je  
    $\pi$ unitarno ekvivalentna $\rho^{(n)}: M_n (\mathfrak{A}) \to \mathcal{B}(\mathcal{H}^n)$.
\end{lema} 

\begin{proof}
    Za poljuben par indeksov $i, j$ definirajmo $P_{ij} := \pi (E_{ij}) \in \bk$.
    Očitno velja $P_{ij} ^* = P_{ji}$ in 
    \begin{equation}\label{eq:1.7}
       P_{ij} P_{kl} = \delta_{jk} P_{il}. 
    \end{equation}
    V posebnem je 
    $$P_{ii} ^2 = P_{ii} = P_{ii}^*,$$
    torej je $P_{ii}$ ortogonalna projekcija na zaprt (po izreku o zaprtem grafu)
    podprostor $\mathcal{H}_i \subseteq \mathcal{K}$.
    Ker je $$P_{11} + \cdots + P_{nn} = I,\quad P_{ii} P_{jj} = 0,$$
    mora veljati $$\mathcal{H}_1 + \cdots + \mathcal{H}_n = \mathcal{K},\quad \mathcal{H}_i \perp \mathcal{H}_j$$
    in posledično $$\mathcal{H}_1 \oplus \cdots \oplus \mathcal{H}_n = \mathcal{K}.$$
    Za vsak $a \in \mathfrak{A}$ ter $P_{ii}$ velja $P_{ii} \pi(aI) = \pi(aI) P_{ii}$, torej je vsak $\mathcal{H}_i$
    reducirajoč prostor za $\pi(aI)$ in je 
    $$\pi_{ii}: \mathfrak{A} \to \mathcal{B}(\mathcal{H}_i),\quad \pi_{ii} (a) = \pi(aI) \big|_{\mathcal{H}_i}$$
    upodobitev.

    Po enačbi \eqref{eq:1.7} imamo $P_{ij} (\mathcal{H}_k) = \delta_{jk} \mathcal{H}_{i}$, 
    kar pomeni, da lahko $P_{ij}$ zapišemo v matrični obliki na $\mathcal{H}_1 \oplus \cdots \oplus \mathcal{H}_n$ kot
    $$P_{ij} = \begin{bmatrix}
    0 & 0 & 0 \\
    0  & 0 & \underbrace{P_{ij}\big|_{\mathcal{H}_j}}_{(i, j)} \\
     0 & 0 & 0 
 \end{bmatrix}.$$
    Iz \eqref{eq:1.7} prav tako sledi $$P_{ij} ^* P_{ij} = P_{ji} P_{ij} = P_{jj},$$
    od koder dobimo
    $$(P_{ij}\big|_{\mathcal{H}_j})^* (P_{ij}\big|_{\mathcal{H}_j}) = (P_{ji}\big|_{\mathcal{H}_i}) (P_{ij}\big|_{\mathcal{H}_j}) = P_{jj} \big|_{\mathcal{H}_j} = \id_{\mathcal{H}_j}.$$
    To implicira, da je $P_{ij}\big|_{\mathcal{H}_j}: \mathcal{H}_j \to \mathcal{H}_i$ unitaren operator
    in $$P_{ij} \pi_{jj}(a) P_{ji} = \pi_{ii} (a),\quad \forall a \in \mathfrak{A}.$$

    Sedaj definiramo $\mathcal{H} := \mathcal{H}_1$, $\rho := \pi_{11}$ in $U_i := P_{i1} \big|_{\mathcal{H}_1} : \mathcal{H}_1 \to\mathcal{H}_i$, zato je 
    $$U = \oplus_i ^n U_i : \mathcal{H}^n \to \oplus_{i = 1} ^n \mathcal{H}_i = \mathcal{K}$$
    unitaren operator. Ker je
    \begin{align*}
        U^* \pi(a E_{ij}) U &= \begin{bmatrix}
            U_1 ^* & & \\
            & \ddots & \\
            & & U_n ^* 
        \end{bmatrix} \pi_{ii} (a) P_{ij} \begin{bmatrix}
            U_1  & & \\
            & \ddots & \\
            & & U_n  
        \end{bmatrix}\\
        &= \begin{bmatrix}
            P_{11}\big|_{\mathcal{H}_1}  & & \\
            & \ddots & \\
            & & P_{1n} \big|_{\mathcal{H}_n} 
        \end{bmatrix} \begin{bmatrix}
           0 & 0 & 0 \\
           0  & 0 & \pi_{ii} (a) {P_{ij}\big|_{\mathcal{H}_j}} \\
            0 & 0 & 0 
        \end{bmatrix}
        \begin{bmatrix}
            P_{11}\big|_{\mathcal{H}_1}  & & \\
            & \ddots & \\
            & & P_{n1} \big|_{\mathcal{H}_1} 
        \end{bmatrix}\\
        &= \begin{bmatrix}
            0 & 0 & 0 \\
            0  & 0 & (P_{1i} \big|_{\mathcal{H}_i}) \pi_{ii} (a){(P_{ij}\big|_{\mathcal{H}_j})}(P_{j1} \big|_{\mathcal{H}_1}) \\
             0 & 0 & 0 
         \end{bmatrix}\\
         &= \begin{bmatrix}
            0 & 0 & 0 \\
            0  & 0 & (P_{1i} \big|_{\mathcal{H}_i}) \pi_{ii} (a)(P_{i1} \big|_{\mathcal{H}_1}) \\
             0 & 0 & 0 
         \end{bmatrix}\\
         &= \begin{bmatrix}
            0 & 0 & 0 \\
            0  & 0 & \pi(a)  \\
             0 & 0 & 0 
         \end{bmatrix},
    \end{align*}
    sledi $U^* \pi([a_{ij}]_{i, j}) U = [\pi(a_{ij})]_{i, j}$.
\end{proof}

Če v zgornji lemi vzamemo $\mathfrak{A} = \C$, potem za vsak $*$-homomorfizem 
$\pi: M_n (\C) \to \bk$ obstaja Hilbertov prostor $\mathcal{H}$ in unitaren operator 
$U: \mathcal{H}^n \to \mathcal{K}$, tako da za vsako matriko
$A = [a_{ij}]_{i, j} \in M_n (\mathfrak{A})$ velja
$$U^* \pi(A) U = A \otimes \id_{\mathcal{H}} \in M_n (\bh) = \mathcal{B} (\mathcal{H}^n).$$

Denimo, da je $\{v_i\}_{i \in I}$ ortonormalna baza $\mathcal{H}$.
Potem imamo dekompozicijo $$\mathcal{H}^n = \oplus _{i \in I} \mathcal{V}_i,$$
pri čemer je $\mathcal{V}_i = \lin \{u_i\}^n$. Zgornja lema nam potem pove, da je
$$U^* \pi(A) U \big|_{\mathcal{V}_i} = A : \mathcal{V}_i \to \mathcal{V}_i,$$
od koder sledi $U^* \pi(A) U = \bigoplus_{i \in I} A$.

\begin{dokaz}[Dokaz izreka \ref{izr:1.2}]
    Naj bo $\varphi: M_n (\C) \to M_m (\C) = \mathcal{B}(\C^m)$ povsem pozitivna preslikava.
    Po Stinespringovem izreku obstaja upodobitev $\pi: M_n (\C) \to \bk$ in $V : \C^m \to \mathcal{K}$, 
    tako da je
    $$\varphi(A) = V^* \pi(A) V$$ za vse $A \in M_n (\C)$.
    Iz konstrukcije Stinespringovega izreka sledi $$\dim (\mathcal{K}) \leq \dim (M_n (\C) \otimes \C^m) = n^2 m.$$
    Po prejšnjem odstavku lahko vzamemo $\mathcal{K} = \bigoplus_{k = 1} ^r \C ^{n} _k$,
    pri čemer je $r \leq nm$ in $\pi(A) \big|_{\C_k ^n} = A$ za vsak $k \leq r$. V bločno matrični obliki preslikavo $V: \C^m \to \bigoplus_{k = 1} ^r \C ^{n} _k$ sedaj zapišemo
    $$V = \begin{bmatrix}
        V_1 \\ \vdots \\ V_r
    \end{bmatrix},\quad  V_i \in \C^{n \times m}.$$
    Če vse to združimo, dobimo
    \begin{align*}
        \varphi(A) &= V^* \pi(A) V\\
        &= \begin{bmatrix}
            V_1 ^* & \cdots & V_r ^*
        \end{bmatrix} \begin{bmatrix}
            A &  & \\
             & \ddots &  \\
             &  & A
        \end{bmatrix}
        \begin{bmatrix}
            V_1 \\ \vdots \\ V_r
        \end{bmatrix} \\
        &= \sum_{k = 1} ^r V_k ^* A V_k
    \end{align*}
    za vsako matriko $A \in M_n (\C)$.
\end{dokaz}

\end{comment}

\section{Povsem pozitivne preslikave na $M_n (\C)$}\label{sec:7}


V tem oglavju karakteriziramo povsem pozitivne preslikave $\varphi: \mathfrak{A} \to \mathfrak{B}$,
kjer je $\mathfrak{A}$ ali $\mathfrak{B}$ oblike $M_n (\C)$. 
Kot posledica izpeljane teorije bo sledilo, da lahko vsako povsem pozitivno preslikavo $\varphi: \mathcal{S} \to M_n (\C)$ 
na poljubnem operatorskem sistemu $\mathcal{S} \subseteq \mathfrak{A}$ razširimo do povsem pozitivne preslikave na celotni $C^*$-algebri $\mathfrak{A}$.
To je ključni rezultat, ki ga bomo potrebovali za dokaz Arvesonovega razširitvenega izreka v naslednjem poglavju.

Začnimo najprej s karakterizacijo povsem pozitivnih preslikav $\varphi: M_n (\C) \to M_m (\C)$.
Ključen pojem bo \emph{Choijeva matrika} $\varphi^{(n)}([E_{ij}]_{i, j})$, kjer $E_{ij}$ označujejo standardne matrične bazne enote v $M_n (\C)$.
Če definiramo vektor 
$$e = \begin{bmatrix}
    e_1 \\ \vdots \\ e_n
\end{bmatrix} \in \C^{n^2},$$
kjer so $e_i$ standardni bazni vektorji v $\C^n$, potem vidimo, da je 
$$e e^* = [E_{ij}]_{i, j} \in M_n (M_n (\C))$$
pozitivna matrika. 
Od tod sledi, da ima povsem pozitivna preslikava $\varphi: M_n (\C) \to M_m(\C)$ vedno pozitivno Choijevo matriko.
Naslednja trditev, ki je bila dokazana v \cite{kraus1971} in \cite{choi1975}, nam pove, da velja tudi obratno.


\begin{comment}
\begin{opomba}
    V zgornjem dokazu smo dokazali, da za vsako pozitivno matriko $A \in M_{mn} (\C)$ obstajajo matrike 
    $V_1, \dots, V_{mn} \in \C^{n \times m}$, tako da je 
    $$A = \sum_{k = 1} ^{mn} (I_n \otimes V_k)^* [E_{ij}]_{i, j} (I_n \otimes V_k).$$
\end{opomba}
\end{comment}

\begin{izrek}[Choi-Kraus]\label{izr:1.2}
    Za linearno preslikavo $\varphi: M_n (\C) \to M_m (\C)$ so naslednje trditve ekvivalentne:
    \begin{enumerate}
        \item $\varphi$ je povsem pozitivna;
        \item $\varphi$ je $n$-pozitivna;
        \item {Choijeva matrika} $\varphi^{(n)}([E_{ij}]_{i, j})$ je pozitivna;
        \item $\varphi$ je oblike 
        \begin{equation}
            \varphi(A) = \sum_{k = 1} ^r V_k ^* A V_k,\quad \forall A \in M_n (\C). \label{eq:1.6}
        \end{equation}
    \end{enumerate}
\end{izrek}


\begin{dokaz}
    Edina netrivialna implikacija je $(3.) \Rightarrow (4.)$, zato predpostavimo,
    da je Choijeva matrika $\varphi^{(n)}([E_{ij}]_{i, j})$ pozitivna.
    Po lemi \ref{lem:1.2} obstajajo vrstični vektorji $v_1, \dots, v_{mn} \in \C^{mn}$, tako da je 
    $$[\varphi (E_{ij})]_{i, j} = \sum_{k = 1} ^{mn} v_k v_k ^* \in M_n (M_m (\C)) = M_{mn} (\C).$$
    Vsakemu stolpcu
    $$v_k = \begin{bmatrix}
        x_1 ^{(k)}  \\ \vdots \\ x_n ^{(k)}
    \end{bmatrix} \in \C^{mn},\quad x_j ^{(k)} \in \C^{ m},$$
    priredimo matriko 
    $$V_k = \begin{bmatrix}
        {x_1 ^{(k)}} ^* \\
         \vdots \\
         {x_n ^{(k)} }^*
    \end{bmatrix} \in \C^{n \times m}.$$
    Opazimo, da velja
    $$V_k ^* E_{ij} V_k  = (V_k ^* e_i) (e_j^* V_k)  = {x_i ^{(k)}} {x_j ^{(k)}} ^* = v_k v_k ^*,$$
    torej
    $$\varphi^{(n)} ([E_{ij}]_{i, j}) = \sum_{k = 1} ^{mn} (I_n \otimes V_k)^* [E_{ij}]_{i, j} (I_n \otimes V_k).$$ 
    Sedaj za vsako matriko $A = \sum_{i, j = 1} ^n a_{ij} E_{ij} \in M_{n} (\C)$ velja
    \begin{align*}
        \varphi (A) &= \varphi \left(\sum_{i, j = 1} ^n a_{ij} E_{ij}\right) = \sum_{i, j = 1} ^n a_{ij} \varphi \left(E_{ij}\right)\\
        &= \sum_{i, j = 1} ^n a_{ij} \sum_{k = 1}^{mn} V_k ^* E_{ij} V_k = \sum_{k = 1}^{mn} \sum_{i, j = 1} ^n a_{ij} V_k ^* E_{ij} V_k\\
        &= \sum_{k = 1}^{mn} V_k ^* \left( \sum_{i, j = 1} ^n a_{ij} E_{ij} \right) V_k = \sum_{k = 1}^{mn} V_k ^* A V_k. \qedhere
    \end{align*}
\end{dokaz}


Naj bo $\mathfrak{B}$ sedaj poljubna $C^*$-algebra.
S pomočjo Choijeve matrike lahko karakteriziramo povsem pozitivne preslikave $\varphi: M_n (\C) \to \mathfrak{B}$.

\begin{trditev}\label{trd:1.6}
    Za linearno preslikavo $\varphi: M_n (\C) \to \mathfrak{B}$ so naslednje
    trditve ekvivalentne:
    \begin{enumerate}
        \item $\varphi$ je povsem pozitivna;
        \item $\varphi$ je $n$-pozitivna;
        \item Choijeva matrika $[\varphi (E_{ij})]_{i, j} \in M_n (\mathfrak{B})$ je pozitivna.
    \end{enumerate}
\end{trditev}

\begin{dokaz}
    Dovolj je dokazati implikacijo $(3.) \Rightarrow (1.)$.
    Brez škode za splošnost predpostavimo, da je $\mathfrak{B} = \bh$ za nek Hilbertov prostor $\mathcal{H}$.
    Naj bo $[A_{ij}]_{i, j} \in M_m (M_n (\C))$ pozitivna matrika.
    Vzemimo poljubne vektorje $x_1, \dots, x_m \in \mathcal{H}$ in dokazujemo, da je 
    $$\sum_{i, j} ^m \langle \varphi (A_{ij}) x_j, x_i \rangle \geq 0.$$
    Definiramo $$V: \C^m \to \mathcal{H},\quad V e_i = x_i$$ in
    $$\psi : M_n (\C) \to M_m (\C),\quad \psi (A) = V^* \varphi(A) V.$$
    Ker je Choijeva matrika 
    $$\psi^{(n)} ([E_{ij}]_{i, j}) = [V^* \varphi (E_{ij}) V]_{i, j} = (I_n \otimes V)^* [\varphi (E_{ij})]_{i, j} (I_n \otimes V)$$
    po predpostavki pozitivna, je po Choijevem izreku $\psi$ povsem pozitivna preslikava. 
    Od tod pa sledi 
    \begin{equation*}
        \sum_{i, j} ^m \langle \varphi (A_{ij}) x_j, x_i \rangle = \sum_{i, j} ^m \langle V^* \varphi (A_{ij}) V e_j, e_i \rangle = \sum_{i, j} ^m \langle \psi (A_{ij}) e_j, e_i \rangle \geq 0. \qedhere
    \end{equation*}
\end{dokaz}

Naš cilj za preostanek tega poglavja je pridobiti podobno karakterizacijo povsem pozitivnih preslikav $\varphi: \mathfrak{A} \to M_n (\C)$.

Naj bo $\mathcal{M} \subseteq \mathfrak{A}$ operatorski prostor.
Za poljubno linearno preslikavo $\varphi \in \mathcal{L} (\mathcal{M}, M_n (\C))$
lahko definiramo $$s_{\varphi} \in \mathcal{L}(M_n (\mathcal{M}), \C),\quad s_{\varphi} (A) = \frac{1}{n} \sum_{i, j} ^n \varphi(a_{ij})_{ij}.$$

Obratno lahko vsakemu linearnemu funkcionalu $s \in \mathcal{L}(M_n(\mathcal{M}), \C) $ priredimo preslikavo 
$$\varphi_s \in \mathcal{L} (\mathcal{M}, M_n (\C)),\quad \varphi_s (a) = n \cdot s^{(n)} ([aE_{ij}]_{i, j}).$$
Faktor $\frac{1}{n}$ v definiciji $s_{\varphi}$ smo dodali zato, da je za enotsko preslikavo $\varphi$ tudi $s_{\varphi}$ enotski funkcional.
Preslikavi 
$$\Phi: \mathcal{B} (\mathcal{M}, M_n (\C)) \to (M_n (\mathcal{M}))^*,\quad \varphi \mapsto s_{\varphi}$$
in 
    $$\Psi: (M_n (\mathcal{M}))^* \to \mathcal{B} (\mathcal{M}, M_n (\C)),\quad s \mapsto {\varphi}_s$$
sta linearni in druga drugi inverzni.

    Linearnost je očitna, zato dokažimo, da je $\Psi \circ \Phi = \id_{\mathcal{B} (\mathcal{M}, M_n (\C))}$.
    Vzemimo poljuben $\varphi \in \mathcal{B} (\mathcal{M}, M_n (\C))$. Potem je za vsak $a \in \mathcal{M}$
    \begin{align*}
        \varphi_{s_{\varphi}} (a) &= n \cdot (s_{\varphi}) ^{(n)} ([a E_{ij}]_{i, j}) = n \cdot [s_{\varphi}(a E_{ij})]_{i, j} \\
        &= n \cdot \frac{1}{n}[\varphi (a)_{ij} ]_{i, j} = [\varphi (a)_{ij} ]_{i, j}\\
        &= \varphi(a).
    \end{align*}
    Dokažimo še $\Phi \circ \Psi = \id_{M_n (\mathcal{M})^*}$. Za $A = [a_{ij}]_{i, j} \in M_n (\mathcal{M})$ je 
    \begin{align*}
        s_{\varphi_s} (A) &= \frac{1}{n} \sum_{i, j} ^n \varphi_s (a_{ij})_{ij} = \frac{1}{n} \sum_{i, j} ^n n \cdot (s^{(n)} ([a_{ij} E_{ij}]_{i, j}))_{ij}\\
        &= \sum_{i, j} ^n ([ s(a_{ij} E_{ij})]_{i, j})_{ij} = \sum_{i, j} ^n  s(a_{ij} E_{ij}) = s(A).
    \end{align*}
    Sedaj lahko dokažemo naslednjo trditev.

\begin{trditev}\label{trd:1.7}
    Za linearno preslikavo $\varphi: \mathcal{S} \to M_n (\C)$ so naslednje trditve ekvivalentne:
    \begin{enumerate}
        \item $\varphi$ je povsem pozitivna;
        \item $\varphi$ je $n$-pozitivna;
        \item $s_\varphi$ je pozitiven funkcional.
    \end{enumerate}
\end{trditev}

\begin{dokaz}
    Dokažimo najprej implikacijo $(2.) \Rightarrow (3.)$.
    Naj bo $$e = \begin{bmatrix}
        e_1 \\ \vdots \\ e_n
    \end{bmatrix} \in \C^{n^2}$$
    kot prej. Potem je $$s_{\varphi} (A) = \frac{1}{n} \sum_{i, j} ^n \varphi(a_{ij})_{ij} = \frac{1}{n} e^* \varphi^{(n)}(A) e$$
    za poljubno matriko $A = [a_{ij}]_{i, j} \in M_n (\mathcal{S})$.
    Če je $\varphi^{(n)}$ pozitivna preslikava, potem je $s_\varphi$ pozitiven funkcional.
    
    Sedaj dokažimo še implikacijo $(3.) \Rightarrow (1.)$.
    Če je $s$ pozitiven funkcional, potem je povsem pozitivna preslikava in posledično je tudi 
    $$\varphi_{s} (a) = n \cdot s^{(n)} ([aE_{ij}]_{i, j}) = n \cdot s^{(n)} (e a e^*)$$
    povsem pozitivna. Če je $s_\varphi$ pozitiven funkcional, potem mora biti $\varphi = \varphi_{s_\varphi}$ povsem pozitivna.
\end{dokaz}

\begin{posledica}
    Naj bo $\mathcal{M} \subseteq \mathfrak{A}$ operatorski sistem.
    Za enotsko linearno preslikavo $\varphi: \mathcal{M} \to M_m (\C)$ so naslednje trditve ekvivalentne:
    \begin{enumerate}
        \item $\varphi$ je povsem skrčitev;
        \item $\varphi$ je $n$-skrčitev;
        \item $s_\varphi$ je skrčitev.
    \end{enumerate}
\end{posledica}

\begin{dokaz}
    Implikacija $(2.) \Rightarrow (3.)$ sledi, ker je 
    $$\| s_\varphi (A) \| = \frac{1}{n} \|e^* \varphi^{(n)}(A) e\| \leq \frac{1}{n} \| e\|^2 \| \varphi^{(n)} (A)\| \leq \frac{1}{n} \cdot n \cdot \| A\|  = \| A\|.$$

    Za implikacijo $(3.) \Rightarrow (1.)$ opazimo, da je $s_\varphi$ enotska skrčitev na $M_n (\mathcal{M})$,
    zato jo lahko razširimo do pozitivnega enotskega funkcionala $\widetilde{s_\varphi}: M_n (\mathcal{M}) + M_n (\mathcal{M})^* \to \C$.
    Po trditvi \ref{trd:1.7} je $\varphi_{\widetilde{s_\varphi}}: \mathcal{M} + \mathcal{M}^* \to M_n (\C)$ povsem pozitivna preslikava.
    Dokažimo, da je razširitev $\varphi$.
    Za poljuben $a \in \mathcal{M}$ je 
    $$\varphi_{\widetilde{s_\varphi}} (a) = n \cdot \widetilde{s_{\varphi}} ^{(n)} ([a E_{ij}]_{i, j}) = n \cdot {s_{\varphi}} ^{(n)} ([a E_{ij}]_{i, j}) = \varphi_{s_\varphi} (a) = \varphi(a),$$
    torej je $\varphi_{\widetilde{s_\varphi}}$ v posebnem enotska preslikava.
    Dokazali smo, da je $\varphi_{\widetilde{s_\varphi}}$ enotska povsem pozitivna preslikava, zato je tudi povsem skrčitev.
    Od tod pa sledi, da velja enako za njeno zožitev $\varphi$.
\end{dokaz}

Trditev \ref{trd:1.7} nam omogoča, da podamo trditev \ref{trd:1.6} za preslikave na operatorskem sistemu.

\begin{posledica}
    Naj bo $\mathcal{S} \subseteq M_n (\C)$ operatorski sistem. 
    Za linearno preslikavo $\varphi: \mathcal{S} \to \mathfrak{B}$ sta naslednji
    trditvi ekvivalentni:
    \begin{enumerate}
        \item $\varphi$ je povsem pozitivna;
        \item $\varphi$ je $n$-pozitivna.
    \end{enumerate}
\end{posledica}

\begin{dokaz}
    Denimo, da je $\varphi$ $n$-pozitivna preslikava. Brez škode za splošnost vzemimo $\mathfrak{B} = \bh$.
    Naj bo $[A_{ij}]_{i, j} \in M_m (M_n (\C))$ pozitivna matrika in $x_1, \dots, x_m \in \mathcal{H}$ poljubni vektorji.
    Ponovno definiramo $$V: \C^m \to \mathcal{H},\quad V e_i = x_i$$ in
    $$\psi : \mathcal{S} \to M_m (\C),\quad \psi (A) = V^* \varphi(A) V.$$
    Ker je $\varphi$ $n$-pozitivna, je $\psi$ prav tako $n$-pozitivna in zato po trditvi \ref{trd:1.7} tudi povsem pozitivna.
    Od tod pa sledi 
    \begin{equation*}
        \sum_{i, j} ^m \langle \varphi (A_{ij}) x_j, x_i \rangle = \sum_{i, j} ^m \langle V^* \varphi (A_{ij}) V e_j, e_i \rangle = \sum_{i, j} ^m \langle \psi (A_{ij}) e_j, e_i \rangle \geq 0. \qedhere
    \end{equation*}
\end{dokaz}

Naj bo $\mathcal{S} \subseteq \mathfrak{A}$ poljuben operatorski sistem.
Dokazujemo glavni izrek tega razdelka, ki nam pove, da lahko povsem pozitivno preslikavo $\varphi: \mathcal{S} \to M_n (\C)$
razširimo do povsem pozitivne preslikave $\psi: \mathfrak{A} \to M_n (\C)$.

Najprej to dokažemo za $n = 1$: naslednji rezultat je zgolj različica Hahn-Banachovega izreka za pozitivne funkcionale.

\begin{lema}[Krein-Riesz]\label{lem:1.4}
    Če je $\varphi: \mathcal{S} \to \C$
    pozitiven funkcional, potem ima pozitivno razširitev $\psi: \mathfrak{A} \to \C$, 
    tako da je $\| \varphi \| = \| \varphi \| =  \varphi (1)$.
\end{lema}

\begin{dokaz}
    Ker je $\varphi: \mathcal{S} \to \C$ pozitiven funkcional, velja po lemi \ref{lem:1.1} $\| \varphi \| = \varphi (1)$.
    Po Hahn-Banachovem izreku obstaja razširitev $\varphi$ do $\varphi: \mathfrak{A} \to \C$,
    tako da je $\| \varphi \| = \| \varphi \|$. Od tod sledi 
    $$\| \psi \| = \| \varphi \| = \varphi (1) = \psi(1),$$
    torej je po lemi \ref{lem:1.1} tudi $\psi$ pozitiven funkcional.
\end{dokaz}

\begin{izrek}\label{izr:aux}
    Če je $\varphi: \mathcal{S} \to M_n (\C)$ povsem pozitivna preslikava,
    potem ima povsem pozitivno razširitev $\psi: \mathfrak{A} \to M_n (\C)$,
    tako da velja $\| \psi \|_{\po} = \| \varphi\|_{\po} = \| \varphi (1)\|$.
\end{izrek}

\begin{dokaz}
    Ker je $\varphi$ povsem pozitivna preslikava, je $s_{\varphi}: M_n(\mathcal{S}) \to \C$
    pozitiven funkcional, zato ga lahko po lemi \ref{lem:1.4} razširimo do pozitivnega funkcionala $s: M_n(\mathfrak{A}) \to \C$.
    Potem je $\psi := \varphi_s : \mathfrak{A} \to M_n (\C)$ povsem pozitivna preslikava.

    Dokazati moramo še, da je razširitev $\varphi$. Res, za poljuben $a \in \mathcal{S}$ imamo 
    $$\psi (a) = \varphi_s (a) = n \cdot s^{(n)} ([a E_{ij}]_{i, j}) = n \cdot s_{\varphi}^{(n)} ([a E_{ij}]_{i, j}) = \varphi_{s_{\varphi}} (a) = \varphi(a).$$
    Ker sta $\varphi$ in $\psi$ povsem pozitivna, sledi še 
    \begin{equation*}
        \| \varphi \|_{\po} \leq \| \psi \|_{\po} = \| \psi (1)\| = \|\varphi(1)\| \leq \| \varphi\|_{\po},
    \end{equation*}
    torej $\| \varphi\|_{\po} = \| \psi \|_{\po} = \| \varphi (1)\|$.
\end{dokaz}

\begin{posledica}
    Naj bo $\mathcal{M} \subseteq \mathfrak{A}$ operatorski prostor.
    Če je $\varphi: \mathcal{M} \to M_n (\C)$ enotska povsem skrčitev,
    potem ima povsem pozitivna razširitev $\psi: \mathfrak{A} \to M_n (\C)$.
\end{posledica}

\begin{dokaz}
    Enotsko povsem skrčitev $\varphi: \mathcal{M} \to M_n (\C)$ lahko razširimo do 
    povsem pozitivne preslikave $\tilde{\varphi}: \mathcal{M} + \mathcal{M}^* \to M_n (\C)$,
    ki jo nato lahko razširimo do povsem pozitivne preslikave $\psi: \mathfrak{A} \to M_n (\C)$ po trditvi \ref{trd:1.7}.
\end{dokaz}

Poglavje zaključimo z verzijo izreka \ref{izr:1.2} za operatorske sisteme.

\begin{trditev}
    Naj bo $\mathcal{S} \subseteq M_n (\C)$ operatorski sistem.
    Za linearno preslikavo $\varphi: \mathcal{S} \to M_m (\C)$ so naslednje trditve ekvivalentne:
    \begin{enumerate}
        \item $\varphi$ je povsem pozitivna;
        \item $\varphi$ je $n$-pozitivna;
        \item $\varphi$ je oblike \eqref{eq:1.5}.
    \end{enumerate}
\end{trditev}

\begin{dokaz}
    Dovolj je dokazati implikacijo $(2.) \Rightarrow (3.)$. Če je $\varphi$ $n$-pozitivna,
    je po trditvi \ref{trd:1.7} povsem pozitivna,
    zato jo lahko po trditvi \ref{izr:aux} razširimo do povsem pozitivne preslikave $\psi: M_n (\C) \to M_m (\C)$, ki mora biti oblike \eqref{eq:1.5}.
    Torej je tudi $\varphi$ oblike \eqref{eq:1.5}.
\end{dokaz}

\section{Arvesonov razširitveni izrek}\label{sec:8}

V poglavju \ref{sec:2} smo dokazali, da lahko vsako povsem pozitivno preslikavo
na operatorskem prostoru razširimo do povsem pozitivne preslikave na pripadajočem operatorskem sistemu.
Iz tamkajšnje diskusije je sledilo, da so operatorski sistemi naraven objekt, na katerem so morfizmi povsem pozitivne preslikave.

Naravno vprašanje je, če imamo 
izrek Hahn-Banachovega tipa, ki bi nam povedal, da 
lahko povsem pozitivno preslikavo $$
\varphi : \mathcal{S} \to \bh
$$ na operatorskem sistemu $\mathcal{S} \subseteq \mathfrak{A}$ razširimo do povsem pozitivne preslikave $$
\psi : \mathfrak{A} \to \bh.
$$ na celotni $C^*$-algebri.
To je vsebina Arvesonovega razširitvenega izreka (\cite{arveson1969}).

V končno dimenzionalnem primeru smo to že dokazali v izreku \ref{izr:aux}.
Natančneje, če je $\mathcal{H}$ končno dimenzionalen,
lahko vsako povsem pozitivno preslikavo
$$
\varphi : \mathcal{S} \to \bh
$$
razširimo do povsem pozitivne preslikave na $\mathfrak{A}$.

Neskončno dimenzionalni primer je bistveno zahtevnejši.
Dokaz temelji na konstrukciji ustrezne topologije na prostoru
$\mathcal{B}(\mathcal{S}, \bh)$,
v kateri bo množica
$$
\PP_r(\mathcal{S}, \bh)
:=
\{L \in \mathcal{B}(\mathcal{S}, \bh)
\mid
\text{$L$ je p.p. in $\|L\| \le r$}\}
$$
kompaktna za vsak $r \ge 0$.
Ta topologija bo izhajala iz šibke-$*$ topologije,
zato bo ključna uporaba Banach-Alaoglujevega izreka iz funkcionalne analize (\cite[Theorem V.3.1., str. 130]{conway1990}).

Naš prvi cilj je zato opisati ustrezno preddualno strukturo prostora
$\mathcal{B}(X,Y^*)$.
Naj bosta $X$ in $Y$ Banachova prostora.
Za $x \in X$ in $y \in Y$ definiramo linearen funkcional
$$
x \otimes y : \mathcal{B}(X,Y^*) \to \mathbb{C},
\qquad
(x \otimes y)(L) := L(x)(y).
$$
\begin{lema}
    Za vsaka $x \in X$ in $y \in Y^*$ velja $\| x \otimes y\| = \| x\| \cdot \| y\|$ in posledično
    $x \otimes y \in \mathcal{B}(X, Y^*)^*$.
\end{lema}

\begin{proof}
    Opazimo, da velja $$| x \otimes y (L)| \leq \| L\| \cdot \| x\| \cdot \| y\|,$$
    od koder sledi $\| x \otimes y\| \leq \| x\| \cdot \| y\|$ in $x \otimes y \in \mathcal{B} (X, Y^*)^*$.
    Dokazujemo, da velja tudi $\| x \otimes y \| = \| x\| \cdot \| y\|$. 
    Preko Hahn-Banachovega izreka lahko konstruiramo funkcional 
    $\varphi_x \in X^*$, tako da je $\varphi_x (x) = \| x\|$ in $\| \varphi_x \| = 1$. 
    Podobno lahko definiramo $\varphi_y \in Y^*$, tako da velja $\varphi_y (y) = \| y\|$ in $\| \varphi_y \| = 1$. 
    Nato definiramo preslikavo $$L_{x, y}: X \to Y^*,\quad L_{x, y} (z) = \varphi_x (z) \varphi_y,$$
    za katero velja
    $$\| L_{x, y} (z)\| \leq | \varphi_x (z) | \cdot \| \varphi_y \| \leq \| \varphi_x\| \cdot \| \varphi_ y\| \cdot \| z\| = \| z\|.$$
    To pomeni, da je $\| L_{x, y} \| \leq 1$ in $L_{x, y} \in \mathcal{B}(X, Y^*)$. Po drugi strani pa je $\| L_{x, y} (x)\| = \| x\| \cdot \| \varphi_y \| = \| x\|$
    in zato $\| L_{x, y} \| = 1$.
    Naposled je
    $$x \otimes y (L_{x, y}) = \| x\| \cdot \| y\| = \| x\| \cdot \| y\| \cdot \| L_{x, y} \|,$$
    zato $\| x \otimes y\| = \| x\| \cdot \| y\|$. 
\end{proof}
 
Definiramo prostor
$$Z := \overline{\lin \{x \otimes y\ |\ x \in X, y \in Y\}} \leq \mathcal{B} (X, Y^*)^*.$$
\begin{lema}
    Banachov prostor $\mathcal{B} (X, Y^*)$ je izometrično izomorfen $Z^*$ s preslikavo
    $$\Phi: \mathcal{B}(X, Y^*) \to Z^*,\quad L \mapsto (x \otimes y \mapsto x \otimes y (L)).$$
\end{lema}

\begin{dokaz}
    Najprej dokažimo, da je $\Phi$ izometrija. Vzemimo poljuben $L \in \mathcal{B}(X, Y^*)$.
    Za vsak $x \otimes y \in Z$ velja
    \begin{align*}
        | \Phi (L) (x \otimes y)| &= | x \otimes y (L)|\\
        &= | L(x) (y) |\\
        &\leq \| L(x)\| \cdot \| y\|\\
        &\leq \| L\| \cdot \| x\| \cdot \| y\|\\
        &= \| L \| \cdot \| x \otimes y\|,
    \end{align*}
    od koder sledi $\| \Phi(L) \| \leq \| L\|$. Po drugi strani pa je
    \begin{align*}
        \| L\| &= \sup_{\| x\| = 1} \| L(x) \| \\
        &= \sup_{\| x\| = 1} \sup_{\| y\| = 1} | L(x) (y)|\\
        &= \sup_{\| x\| = 1} \sup_{\| y\| = 1} | x \otimes y (L)|\\
        &= \sup_{\| x\| = 1} \sup_{\| y\| = 1} | \Phi (L) (x \otimes y)|\\
        &\leq \sup_{\| x \otimes y\| = 1} |\Phi (L) (x \otimes y)\| = \| \Phi (L)\|. 
    \end{align*}
    Nazadnje dokažimo še surjektivnost. Fiksirajmo nek $f \in Z^*$.
    Za vse $x \in X$ lahko definiramo funkcional
    $$f_x : Y \to \C,\quad y \mapsto f(x \otimes y).$$
    Ker je $| f_x (y)| \leq \| f\| \cdot \| x\| \cdot \| y\|$, sledi $f_x \in Y^*$.
    Potem je $$L : X \to Y^*,\quad L(x) = f_x.$$
    omejena preslikava in $\Phi(L) = f$. 
\end{dokaz}

Če identificiramo prostor $\mathcal{B} (X, Y^*)$ z $Z^*$, ga lahko opremimo s šibko-$*$ topologijo.
Tej topologiji pravimo \emph{omejeno šibka (OŠ) topologija}.

\begin{lema}
    Naj bo $(L_\lambda)_{\lambda}$ omejena mreža v $\mathcal{B}(X, Y^*)$.
    Potem konvergira $L_\lambda \xrightarrow{} L$ v OŠ topologiji natanko tedaj, ko $L_\lambda (x) \xrightarrow{} L(x)$ v šibki-$*$ topologiji za vse $x \in X$.
\end{lema}

\begin{dokaz}
    Najprej dokažemo implikacijo $(\Rightarrow)$. Če $L_\lambda \to L$ v OŠ topologiji, potem 
    $$L_\lambda (x) (y) = \Phi (L_\lambda) (x \otimes y) \to \Phi (L) (x \otimes y) = L(x) (y)$$
    za vse $x \in X$ in $y \in Y$, torej $L_\lambda (x) \to L (x)$ v šibki-$*$ topologiji na $Y^*$. 
    
    Sedaj dokažimo še implikacijo $(\Leftarrow)$. Ker je mreža $(L_\lambda)_\lambda$ omejena z normo, obstaja $M$, tako da je $\| L_\lambda \| \leq M$ za vsak $\lambda \in \Lambda$.
    Če konvergira $L_\lambda (x) \xrightarrow{} L(x)$ v šibki-$*$ topologiji za vse $x \in X$, tedaj
    $$\Phi (L_\lambda) (x \otimes y) = L_\lambda (x) (y) \to L(x) (y) = \Phi (L) (x \otimes y)$$
    za vse $x \otimes y \in \mathcal{B}(X, Y^*)^*$. Posledično $\Phi (L_\lambda)(z) \to \Phi(L) (z)$ za vse $z \in \lin \{x \otimes y\ |\ x \in X, y \in Y\}$. 
    Po definiciji je $Z = \overline{\{x \otimes y\ |\ x \in X, y \in Y\}}$. Vzemimo poljuben $z \in Z$.
    Potem obstaja $z_0 \in \lin \{x \otimes y\ |\ x \in X, y \in Y\}$, tako da je $\| z - z_0 \| \leq \frac{\varepsilon}{3}$.
    Ker je $\Phi (L_\lambda) (z_0) \to \Phi (L) (z_0)$, obstaja $\lambda_0 \in \Lambda$,
    tako da je $\| \Phi (L_\lambda) (z_0) - \Phi (L) (z_0)\| \leq \frac{\varepsilon}{3M}$ za vse $\lambda \geq \lambda_0$. Potem za vsak $\lambda \geq \lambda_0$ velja tudi 
    \begin{align*}
        \| \Phi (L_\lambda) (z) - \Phi (L) (z)\| &\leq \| \Phi (L_\lambda) (z) - \Phi (L_\lambda) (z_0)\| + \| \Phi (L_\lambda) (z_0) - \Phi (L) (z_0)\|\\
         &+ \|  \Phi (L) (z_0) - \Phi (L) (z)\|\\
        &\leq \| \Phi (L_\lambda) \| \cdot \| z - z_0\| + \frac{\varepsilon}{3} + \| \Phi (L) \| \cdot  \| z_0 - z\|\\
        &\leq \| L_\lambda \| \cdot \| z - z_0\| + \frac{\varepsilon}{3} + \| L \| \cdot  \| z_0 - z\|\\
        &\leq M \cdot \| z - z_0\| + \frac{\varepsilon}{3} + M \cdot  \| z_0 - z\|\\
        &\leq \frac{\varepsilon}{3} + \frac{\varepsilon}{3} + \frac{\varepsilon}{3} = \varepsilon,
    \end{align*}
    torej $\Phi (L_\lambda) (z) \to \Phi (L) (z)$.
\end{dokaz}

Spomnimo, da za poljuben Hilbertov prostor $\mathcal{H}$ obstaja izometrični izomorfizem  
$$\bh \cong L^1 (\mathcal{H})^*,\quad T \mapsto (S \mapsto \tr(TS))$$
pri čemer $L^1 (\mathcal{H})$ označuje Banachov prostor operatorjev s sledjo, opremljen z normo $\| S \|_1 = \tr (|S|)$.

\begin{lema}
    Če je $\{L_\lambda\}_\lambda \subseteq \mathcal{B} (X, \bh)$ omejena mreža,
    potem $L_\lambda \to L$ konvergira v OŠ topologiji natanko tedaj, ko za vse $x \in X$ in $u, v \in \mathcal{H}$ velja
    $\langle L_\lambda (x) u, v \rangle \to \langle L(x) u, v \rangle$.
\end{lema}

\begin{dokaz}
    Za vsaka $u, v \in \mathcal{H}$ definiramo operator ranga $1$ 
    $$S_{u, v} \in \bh,\quad S_{u, v} (\alpha) = \langle \alpha, v \rangle \cdot u.$$
    Opazimo, da velja $\tr (S_{u, v}) = \langle u, v\rangle$.
    Linearne kombinacije operatorjev ranga $1$ tvorijo množico operatorjev končnega ranga,
    ki so gosti v $L^1 (\mathcal{H})$. 

    Sedaj imamo naslednji niz ekvivalenc:
    \begin{align*}
        L_\lambda \to L &\Leftrightarrow L_\lambda (x) \to L (x),\ \forall x \in X \\
        &\Leftrightarrow \tr (L_\lambda (x) S) \to \tr (L (x) S),\ \forall x \in X,\ \forall S \in L^1 (\mathcal{H})\\
        &\Leftrightarrow \tr (L_\lambda (x) S_{u, v}) \to \tr (L (x) S_{u, v}),\ \forall x \in X,\ \forall u, v \in \mathcal{H}\\
        &\Leftrightarrow \langle L_\lambda (x) u, v \rangle \to \langle L (x) u, v \rangle,\ \forall x \in X,\ \forall u, v \in \mathcal{H}. \qedhere
    \end{align*}
\end{dokaz}

\begin{opomba}
    Topologiji na $\bh$, inducirani z družino polnorm 
    $$T \mapsto |\langle Tu, v\rangle|,\quad u, v \in \mathcal{H},$$
    pravimo \emph{šibka operatorska topologija} (od tod naprej ŠOT).
    Prejšnja lema nam torej pove, da omejena mreža $\{L_\lambda\}_{\lambda \in \Lambda} \subseteq \mathcal{B}(X, \bh)$
    konvergira v OŠ topologiji natanko tedaj, ko $\{L_\lambda (x)\}_{\lambda \in \Lambda}$
    konvergira v ŠOT za vsak $x \in X$.
\end{opomba}

\begin{lema}\label{lem:1.5}
    Naj bo $\mathfrak{A}$ $C^*$-algebra in $\mathcal{S} \subseteq \mathfrak{A}$ zaprt operatorski sistem. Potem je 
    $$\PP_r (\mathcal{S}, \bh) := \{L \in \mathcal{B} (\mathcal{S}, \bh)\ |\ \textrm{$L$ p.p.~in$\| L\| \leq r$}\}$$
    kompaktna množica v OŠ topologiji.
\end{lema}

\begin{proof}
    Po Banach-Alaoglujevem izreku je
    $$B_r (\mathcal{S}, \bh) := \{L \in \mathcal{B} (\mathcal{S}, \bh)\ |\ \| L\| \leq r\}$$
    kompaktna množica v OŠ topologiji. Preostane le še pokazati, da je $\PP_r$ zaprt podprostor $B_r$.
    Naj bo $(L_\lambda)_\lambda$ mreža v $\PP_r$, ki konvergira v OŠ topologiji k nekemu $L \in B_r$.
    Dokazati moramo, da je $L$ povsem pozitivna. Naj bo $n \in \N$ in $[a_{ij}]_{i, j} \in M_n (\mathcal{S})$ pozitivna matrika.
    Potem za poljubne $x_1, \dots, x_n \in \mathcal{H}$ velja
    $$\sum_{i, j} ^n \langle L_\lambda  ([a_{ij}]_{i, j})x_j, x_i \rangle \to \sum_{i, j} ^n \langle L  ([a_{ij}]_{i, j}) x_j, x_i \rangle.$$
    Ker je vsota na levi nenegativna za vsak $\lambda$, mora veljati enako za vsoto na desni.
\end{proof}

Sedaj lahko dokažemo glavni izrek poglavja.

\begin{izrek}[Arveson]\label{izr:1.3}
    Naj bo $\mathfrak{A}$ $C^*$-algebra, $\mathcal{S} \subseteq \mathfrak{A}$ operatorski sistem in $\varphi: \mathcal{S} \to \bh$ povsem pozitivna.
    Potem lahko $\varphi$ razširimo do povsem pozitivne preslikave $\psi: \mathfrak{A} \to \bh$, tako da velja
    $\| \psi \|_{\po} = \| \varphi \|_{\po} = \| \varphi (1)\|$.
\end{izrek}

\begin{dokaz}
    Brez škode za splošnost lahko predpostavimo, da je $\mathcal{S}$ zaprt podprostor in zato Banachov.
    Za vsak končno dimenzionalen podprostor $\mathcal{F} \leq \mathcal{H}$
    definiramo povsem pozitivno preslikavo $$\varphi_{\mathcal{F}}: \mathcal{S} \to \mathcal{B}(\mathcal{F}),\quad a \mapsto P_{\mathcal{F}} \varphi(a)\big|_{\mathcal{F}},$$
    kjer je $P_\mathcal{F}: \mathcal{H} \to \mathcal{F}$ projekcija.
    Ker je $\dim (\mathcal{F}) = n < \infty$, lahko identificiramo $\mathcal{B} (\mathcal{F}) \equiv M_n (\C)$.
    Po trditvi \ref{izr:aux} lahko $\varphi_{\mathcal{F}}$ razširimo do povsem pozitivne preslikave $\varphi_\mathcal{F} ' : \mathfrak{A} \to \mathcal{B}(\mathcal{F})$,
    tako da je $\| \varphi_{\mathcal{F}} '\| = \| \varphi_\mathcal{F}\|$.
    Nazadnje za razcep $\mathcal{H} = \mathcal{F} \oplus \mathcal{F}^\perp$ definiramo povsem pozitivno preslikavo
    $$\psi_{\mathcal{F}}: \mathfrak{A} \to \bh,\quad \psi_{\mathcal{F}} (a) = \begin{bmatrix}
        \varphi_{\mathcal{F}} ' (a) & 0\\
        0 & 0
    \end{bmatrix}$$
    in vidimo, da velja $\| \psi_{\mathcal{F}} \| = \| \varphi_{\mathcal{F}} ' \| = \| \varphi_{\mathcal{F}} \| \leq \| \varphi\|$.

    Opazimo, da je $\{\psi_{\mathcal{F}}\}$ mreža v $\PP_{\| \varphi\|} (\mathfrak{A}, \bh)$,
    pri čemer $\mathcal{F}$ teče po vseh končnih podprostorih v $\mathcal{H}$.
    Ker je $\PP_{\| \varphi\|} (\mathfrak{A}, \bh)$ kompaktna množica, ima mreža $\{\psi_{\mathcal{F}}\}$
    stekališče v $\psi \in \PP_{\| \varphi\|} (\mathfrak{A}, \bh)$. 

    Dokazujemo, da je $\psi\big|_{\mathcal{S}} = \varphi$. Fiksirajmo $a \in \mathcal{S}$. 
    Nato vzemimo poljubna vektorja $x, y \in \mathcal{H}$
    in definirajmo prostor $\mathcal{F}_{x, y} := \lin \{x, y\} \leq \mathcal{H}$. Potem za vsak končnorazsežen prostor $\mathcal{K} \subseteq \mathcal{H}$,
    za katerega velja $\mathcal{F} \geq \mathcal{F}_{x, y}$, dobimo 
    $$\langle \psi_{\mathcal{F}}(a) x, y \rangle  = \langle P_{\mathcal{F}} \varphi(a)P_{\mathcal{F}} x, y\rangle = \langle \varphi (a) x, y  \rangle,$$
    torej $\psi_{\mathcal{F}} (a) \to \varphi (a)$ v ŠOT. Po drugi strani pa ima 
    $\{\psi_{\mathcal{F}} (a)\}_{\mathcal{F}}$ stekališče v $\psi(a)$, od koder sledi $\varphi(a) = \psi(a)$.

    Oglejmo si še normo razširitve $\psi$. Ker sta $\varphi$ in $\psi$ povsem pozitivna, velja
    \begin{equation*}
        \| \psi \|_{\po} = \| \psi \| = \| \psi (1) \| = \| \varphi(1) \| = \| \varphi \| = \| \varphi \|_{\po}. \qedhere 
    \end{equation*}
\end{dokaz}

\begin{opomba}
    Predpostavka o povsem pozitivnosti v Arvesonovem razširitvenem izreku je ključna,
    saj v splošnem ne moremo razširiti pozitivne preslikave na operatorskem sistemu do pozitivne preslikave na celotni $C^*$-algebri.
    Vzemimo na primer pozitivno preslikavo $\varphi$ iz zgleda \ref{zgl:1.6}, za katero velja $\| \varphi \| = 2 \cdot \| \varphi(1)\|$.
    Po trditvi \ref{pos:russo-dye} pa za vsako pozitivno preslikavo na $C^*$-algebri velja $\| \varphi \| = \| \varphi(1)\|$.
\end{opomba}

\begin{posledica}
    Naj bo $\mathcal{M} \subseteq \mathfrak{A}$ operatorski prostor in $\varphi: \mathcal{M} \to \bh$ enotska povsem skrčitev.
    Potem lahko $\varphi$ razširimo do enotske povsem pozitivne preslikave $\psi: \mathfrak{A} \to \bh$.
\end{posledica}




\section{Abstraktni operatorski sistemi}\label{sec:9}

V funkcionalni analizi imajo konkretni objekti pogosto tudi abstraktno karakterizacijo.
Konkretne $C^*$-algebre (oziroma von Neumannove algebre) operatorjev na Hilbertovem prostoru lahko abstraktno karakteriziramo 
kot $C^*$-algebre (oziroma $W^*$-algebre).
Choi in Effros sta pokazala (\cite{choieffros1977}), da tudi operatorski sistemi, ki jih definiramo kot konkretne prostore operatorjev v $C^*$-algebri, 
premorejo abstraktno karakterizacijo. To nam omogoča, da lahko na operatorske sisteme gledamo kot na abstrakten vektorski prostor
ali pa na neko njegovo konkretno realizacijo v $C^*$-algebri.

Začnimo z nekaterimi potrebnimi lastnostmi, ki jih pričakujemo od abstraktnih operatorskih sistemov.
Vzemimo kompleksen vektorski prostor $\mathcal{S}$, opremljen z involucijo, torej s poševno linearno preslikavo 
$*: \mathcal{S} \to \mathcal{S}$, tako da velja $a^{**} = a$ za vse $a \in \mathcal{S}$.
Takšnemu vektorskemu prostoru pravimo \emph{$*$-vektorski prostor}. 
Če je $\mathcal{S}$ $*$-vektorski prostor, velja enako tudi za $M_n (\mathcal{S})$.
Množico sebiadjungiranih elementov v $*$-vektorskem sistemu označimo s 
$\mathcal{S}_{\sa} := \{s \in \mathcal{S}\ |\ s = s^*\}$.

Tudi v abstraktnem operatorskem sistemu želimo imeti ustrezne elemente, ki imajo vlogo pozitivnih elementov.
Zato mora za vsak $n \in \N$ obstajati konveksen stožec $\mathcal{C}_n \subseteq M_n (\mathcal{S})_h$.
Elementi stožcev $\mathcal{C}_n$ bodo imeli vlogo pozitivnih operatorjev, zato morajo biti ti stožci med seboj na ustrezen način usklajeni:
za poljubna $m, n \in \N$ in $A \in \C^{m \times n}$ mora veljati $A^* \mathcal{C}_m A \subseteq \mathcal{C}_n$.
Za $*$-vektorski prostor $\mathcal{S}$, opremljen s stožci $\{\mathcal{C}_n\}_{n \in \N}$,
pravimo, da je \emph{matrično urejen}. Če imamo matrično urejenima $*$-vektorska prostora $\mathcal{S}, \mathcal{S}'$ s stožci $\{\mathcal{C}_n\}_{n \in \N}, \{\mathcal{C}_n '\}_{n \in \N}$,
potem je preslikava $\varphi: \mathcal{S} \to \mathcal{S}'$ \emph{$n$-pozitivna}, če $\varphi(\mathcal{C}_n) \subseteq \mathcal{C}' _n$,
in \emph{povsem pozitivna}, če je $n$-pozitivna za vsak $n \in \N$. Kot v konkretnih operatorskih sistemih so 
povsem pozitivne preslikave naravni morfizem med abstraktnimi operatorski sistemi. Množico (enotskih) povsem pozitivnih preslikav iz $\mathcal{S}$ v $\mathcal{S}'$ 
ponovno označujemo s $\PP (\mathcal{S}, \mathcal{S}')$
($\EPP (\mathcal{S}, \mathcal{S}')$).

Nazadnje v abstraktnem operatorskem prostoru potrebujemo še enico.
Element $e \in \mathcal{S}_{\sa}$ je \emph{enota za urejenost} v $\mathcal{S}$, 
če za vsak $a \in \mathcal{S}_{\sa}$ obstaja pozitivno realno število $r > 0$, tako da je 
$er + a \in \mathcal{C}_1$. Če velja, da iz $re + a \in \mathcal{C}_1$ za vse $r > 0$
sledi, da je $a \in \mathcal{C}_1$, potem $e$ rečemo \emph{arhimedska enota}.
Če je $I_n \in M_n (\mathcal{S})$ enota za urejenost za vse $n \in \N$, potem je $e$ \emph{enota za matrično urejenost}.
V primeru, ko je $I_n $ Arhimedova za vsak $n \in \N$, je $e$ \emph{arhimedska enota za matrično urejenost} za $\mathcal{S}$.

Opazimo lahko, da obstaja $r > 0$, tako da je $re + e \in \mathcal{C}_1$,
zato je $e = (r + 1)^{-1} (re + e) \in \mathcal{C}_1$. Na enak način dobimo tudi $I_n  \in \mathcal{C}_n$
za vsak $n \in \N$. Če velja $re + a \in \mathcal{C}_1$, potem za vsak $s > r$ velja 
$se + a = (s - r)e + (re + a) \in \mathcal{C}_1 $. Če je $a \in \mathcal{S}_{\sa}$, potem obstaja dovolj velik $r > 0$,
tako da je $re \pm a \in \mathcal{C}_1$ in posledično 
$$a = \frac{re + a}{2} - \frac{re - a}{2} \in \mathcal{C}_1 - \mathcal{C}_1.$$
Od tod sledi, da je $\mathcal{S}_{\sa} = \mathcal{C}_1 - \mathcal{C}_1$.

\begin{definicija}
    Matrično urejenemu $*$-vektorskemu prostoru $\mathcal{S}$ z arhimedsko 
    enoto za matrično urejenost $e \in \mathcal{S}$ pravimo \emph{abstraktni operatorski sistem}.
\end{definicija}

Ker so konkretni operatorski sistemi vedno normirani, bi enako pričakovali tudi od abstraktnih operatorskih prostorov.

\begin{trditev}
    Za abstrakten operatorski sistem $\mathcal{S}$ lahko definiramo normo na $M_n (\mathcal{S})$ s predpisom
    $$\| A\|_n := \inf \left\lbrace r\ |\ \begin{bmatrix}
        r I_n & A\\
        A^* & rI_n
    \end{bmatrix} \in \mathcal{C}_{2n} \right\rbrace,$$
    za katero velja $\| I_n\|_n = 1$ in $\| A\|_n = \| A^*\|_n$ za vse $A \in M_n(\mathcal{S})$.
    V topologiji, inducirani s to normo, je $\mathcal{C}_n \subseteq M_n (\C)$ zaprta množica.
\end{trditev}


\begin{dokaz}
    \begin{enumerate}
        \item Najprej dokazujemo, da je $\| a\|_1 \geq 0$ za vsak $a \in \mathcal{S}$.
        Če je $$\begin{bmatrix}
            r e & a\\
            a^* & re
        \end{bmatrix} \in \mathcal{C}_2,$$
        potem je tudi 
        $$\begin{bmatrix}
            1 & 0\\
            0 & -1
        \end{bmatrix}\begin{bmatrix}
            r e & a\\
            a^* & re
        \end{bmatrix}\begin{bmatrix}
            1 & 0\\
            0 & -1
        \end{bmatrix} = \begin{bmatrix}
            re & -a\\
            -a^* & re
        \end{bmatrix} \in \mathcal{C}_2.$$
        To pomeni, da je tudi njuna vsota $2r I_2 \in \mathcal{C}_2$, zato je $2r \geq 0$ in $\| a\|_1 \geq 0$.
        \item Nato dokazujemo, da iz $\| a\|_1 = 0$ sledi $a = 0$.
        Če res velja $\| a\|_1$, potem je za vsak $r > 0$
        $$\begin{bmatrix}
            re & a\\
            a^* & re
        \end{bmatrix} \in \mathcal{C}_2$$
        in za vsak $\lambda \in \C$ 
        $$\begin{bmatrix}
            1 & \overline{\lambda}
        \end{bmatrix}\begin{bmatrix}
            re & a\\
            a^* & re
        \end{bmatrix}\begin{bmatrix}
            1 \\ \lambda
        \end{bmatrix} = r (1 + |\lambda|^2)e + \lambda a + (\lambda a)^* \in \mathcal{C}_1.$$
        Po arhimedski lastnosti je $\lambda a + (\lambda a)^* \in \mathcal{C}_1$ za vsak $\lambda \in \C$.
        Če vstavimo $\lambda = 1$, dobimo $a + a^* = 0$, če pa $\lambda = i$, pa dobimo $ia - ia^* = 0$.
        Ko ti enačbi seštejemo, dobimo $a = 0$.
        \item Naprej dokazujemo, da za vsak $\lambda \in \C$ velja $\| \lambda a\|_1 = |\lambda| \| a \|_1$.
        Razpišimo $\lambda = |\lambda| \xi$, kjer je $| \xi| = 1$. Potem je 
        $$\begin{bmatrix}
            \frac{re}{|\lambda|} & a\\
            a^* & \frac{re}{|\lambda|}
        \end{bmatrix} \in \mathcal{C}_2$$
        natanko tedaj, ko je 
        \begin{align*}
            |\lambda| \cdot 
            \begin{bmatrix}
                1 & 0\\
                0 & \xi^*
            \end{bmatrix}  
            \begin{bmatrix}
                \frac{re}{|\lambda|} & a\\
                a^* & \frac{re}{|\lambda|}
            \end{bmatrix}
            \begin{bmatrix}
                1 & 0\\
                0 & \xi
            \end{bmatrix} 
            = \begin{bmatrix}
                re & \xi |\lambda| a\\
                \xi^* |\lambda| a^* &  \xi^* \xi re
            \end{bmatrix} = \begin{bmatrix}
                re & \lambda a\\
                \lambda^* a^* &  re
            \end{bmatrix} \in \mathcal{C}_2.
        \end{align*}
        Od tod sledi, da je 
        $$|\lambda| \| a\|_1 = \inf \left\lbrace r\ |\ \begin{bmatrix}
            \frac{re}{|\lambda|} & a\\
            a^* & \frac{re}{|\lambda|}
        \end{bmatrix} \in \mathcal{C}_2 \right\rbrace= \inf  \left\lbrace r\ |\ \begin{bmatrix}
            re & \lambda a\\
            \lambda^* a^* &  re
        \end{bmatrix} \in \mathcal{C}_2 \right\rbrace = \|\lambda a\|_1.$$
        \item Dokažimo, da je $\| a + b\|_1 \leq \| a\|_1 + \| b\|_1$.
        Nastavimo $r_1 := \| a\|_1$ in $r_2 = \| b\|_1$. Potem je 
        $$\begin{bmatrix}
            r_1 e & a\\
            a^* & r_1 e
        \end{bmatrix}, \ \begin{bmatrix}
            r_2 e & a\\
            a^* & r_2 e
        \end{bmatrix} \in \mathcal{C}_2$$
        in zato
        $$ \begin{bmatrix}
            (r_1 + r_2) e & (a + b)\\
            (a + b)^* & (r_1 + r_2) e
        \end{bmatrix} \in \mathcal{C}_2.$$
        Zato je po definiciji $\| a + b\|_1 \leq r_1 + r_2 = \| a\|_1 + \| b\|_1$.
        \item V tem koraku pokažemo, da velja $\| e\|_1 = 1$. Ker je $e \in \mathcal{C}_1$ in 
        $$\begin{bmatrix}
            e & e\\
             e & e
        \end{bmatrix} = \begin{bmatrix}
            1 \\ 1
        \end{bmatrix} e \begin{bmatrix}
            1 & 1
        \end{bmatrix} \in \mathcal{C}_2,$$
        velja $\| e\|_1 \leq 1$. Za $r < 1$ je 
        $$\begin{bmatrix}
            1 & -1
        \end{bmatrix} \begin{bmatrix}
            re & e\\
             e & re
        \end{bmatrix} \begin{bmatrix}
            1 & -1
        \end{bmatrix} = 2 (r - 1) e \notin \mathcal{C}_1,$$
        zato $$\begin{bmatrix}
            re & e\\
             e & re
        \end{bmatrix} \notin \mathcal{C}_2$$
        in $\| e\|_1 = 1$.
        \item Dokažimo še, da je $\| a\| = \| a^*\|$.
        Velja
        $$\begin{bmatrix}
            re & a\\
            a^* & re
        \end{bmatrix} \in \mathcal{C}_2$$
        natanko tedaj, ko je 
        \begin{align*}
            \begin{bmatrix}
                0 & 1\\
                1 & 0
            \end{bmatrix}  
            \begin{bmatrix}
                re & a\\
                a^* & re
            \end{bmatrix}
            \begin{bmatrix}
                0 & 1\\
                1 & 0
            \end{bmatrix}
            = \begin{bmatrix}
                re & a^*\\
                a & re
            \end{bmatrix} \in \mathcal{C}_2.
        \end{align*}
        Od tod sledi $\| a\|_1 = \| a^*\|_1$. 
        \item Dokažimo še, da je $\mathcal{C}_1$ zaprta množica.
        Vzemimo zaporedje $\{a_n\}_{n \in \N} \subseteq \mathcal{C}_n$, ki konvergira v normi k nekemu $a \in \mathcal{S}$.
        Najprej dokažimo, da je $a \in \mathcal{S}_{\sa}$. Ker je $\lim_{n \to \infty} \| a - a_n\|_1 = 0$,
        je po prejšnji točki tudi $\lim_{n \to \infty} \| (a - a_n)^*\|_1 = \lim_{n \to \infty} \| a^* - a_n\| = 0$,
        zato je $a = a^*$.
        Sedaj dokazujemo, da za poljuben $r > 0$ velja $a + re \in \mathcal{C}_1$.
        Vemo, da obstaja $n \in \N$, tako da je $\| a - a_n \| < r$.
        Potem je $$\begin{bmatrix}
            re & a - a_n\\
            a - a_n & re
        \end{bmatrix} \in \mathcal{C}_2$$
        in zato 
        $$\begin{bmatrix}
            1 & 1
        \end{bmatrix}\begin{bmatrix}
            re & a - a_n\\
            a - a_n & re
        \end{bmatrix} \begin{bmatrix}
            1 \\ 1
        \end{bmatrix} = 2 re + 2 (a - a_n) \in \mathcal{C}_1.$$
        Od tod pa sledi $2re + 2a \in \mathcal{C}_1$ in $re + a \in \mathcal{C}_1.$
        \qedhere
    \end{enumerate}
\end{dokaz}

Do sedaj nismo spoznali še nobenega primera abstraktnega operatorskega sistema,
vendar pa bi po konstrukciji pričakovali, da lahko vsakemu 
konkretnemu operatorskemu sistemu $\mathcal{S} \subseteq \mathfrak{A}$ na kanoničen način podamo strukturo abstraktnega operatorskega sistema.
Involucija $*$ je seveda inducirana s $C^*$-algebro, konveksni stožci pa so $\mathcal{C}_n = M_n (\mathcal{S})_+$.
Dokažimo še, da je $1 \in \mathcal{S}$ arhimedska enota za matrično urejenost.
Za poljubno sebiadjungirano matriko $A \in M_n (\mathcal{S})$ je $\| A\| \cdot I_n - A \geq 0$, torej je $1$ enota za matrično urejenost.
Dokazati moramo še, da je $1 \in \mathcal{S}$ arhimedska enota.
Denimo, da velja $r \cdot I_n + A \geq 0$ za vse $r > 0$. 
Potem je $A$ sebiadjungirana, zato ima realen spekter: $\sigma(A) \subseteq \R$.
Ker za vsak $r > 0$ velja $\sigma(r \cdot I_n + A) = r + \sigma(A) \subseteq [0, \infty)$,
sledi $\sigma(A) \subseteq [0, \infty)$ in zato je $A \in M_n (\mathfrak{A})_+$.

Takoj sledi, da je vsaka povsem pozitivna preslikava $\varphi: \mathcal{S}_1 \to \mathcal{S}_2$ med konkretnima operatorskima prostoroma tudi povsem pozitivna,
če $\mathcal{S}_1, \mathcal{S}_2$ obravnavamo kot abstraktna operatorska sistema z zgoraj opisano strukturo.
V nadaljevanju tega poglavja bomo torej vsak konkreten operatorski sistem implicitno obravnavali kot abstrakten operatorski sistem s strukturo, opisano v zgornjem odstavku.

\begin{zgled}
    Če je $\mathcal{S}$ abstrakten operatorski sistem in $\varphi: \mathcal{S} \to \C$
    pozitiven funkcional (torej $\varphi (\mathcal{C}_1) \subseteq [0, \infty)$),
    potem je $\varphi$ povsem pozitivna preslikava. Dokaz poteka popolnoma enako kot v zgledu \ref{zgl:1.3}.
\end{zgled}

Glavni izrek tega poglavja nam pove, da velja tudi delni obrat: vsak abstrakten operatorski sistem je izomorfen (kot abstrakten operatorski sistem)
nekemu konkretnemu operatorskemu sistemu.

\begin{izrek}[Choi-Effros]\label{izr:choi-effros}
    Za vsak abstrakten operatorski sistem $\mathcal{S}$ obstaja (konkreten) operatorski sistem $\mathcal{S}_1 \subseteq \mathfrak{A}$
    in izomorfizem $\varphi: \mathcal{S} \to \mathcal{S}_1$, ki slika $e \in \mathcal{S}$ v $1 \in \mathcal{S}_1$.
\end{izrek}

Za dokaz potrebujemo pomožno trditev.
Kot pri dokazu trditve \ref{trd:1.7} lahko tudi za
abstrakten operatorski prostor $\mathcal{S}$
poljubni preslikavi $\varphi \in \mathcal{L}(\mathcal{S}, M_n (\C))$
priredimo linearen funkcional 
$$s_\varphi \in \mathcal{L}(M_n (\mathcal{S}), \C),\quad s_\varphi ([a_{i, j}]_{i, j}) = \sum_{i, j} ^n \varphi(a_{ij})_{ij} = e^* \varphi^{(n)} ([a_{ij}]_{i, j}) e,$$
pri čemer je $$e = \begin{bmatrix}
    e_1 \\ \vdots \\ e_n
\end{bmatrix} \in \C^{n^2}$$
(tokrat izpustimo faktor $\frac{1}{n}$, saj ga ne potrebujemo).
Obratno lahko vsakemu linearnemu funkcionalu $s \in \mathcal{L}(M_n (\mathcal{S}), \C)$
priredimo $$\varphi \in \mathcal{L}(\mathcal{S}, M_n (\C)),\quad \varphi (a) = s^{(n)} ([a E_{ij}]_{i, j}) = n \cdot s^{(n)} (e a e^*).$$
Tudi tokrat sta preslikavi 
$\varphi \mapsto s_{\varphi}$ in $s \mapsto \varphi_s$ druga drugi inverzni.
Ponovno velja analog trditve \ref{trd:1.7}. 
Dokaz je popolnoma enak, zato ga izpustimo.

\begin{trditev}
    Naj bo $\mathcal{S}$ abstrakten operatorski sistem in $\varphi \in \mathcal{L}(\mathcal{S}, M_n (\C))$. Naslednje trditve so ekvivalentne:
    \begin{enumerate}
        \item $\varphi$ je povsem pozitivna;
        \item $\varphi$ je $n$-pozitivna;
        \item $s_\varphi$ je pozitiven funkcional.
    \end{enumerate}
\end{trditev}

Sedaj lahko dokažemo Choi-Effrosov izrek.

\begin{dokaz}[Dokaz izreka \ref{izr:choi-effros}]
    Za vsak $n \in \N$ označimo 
    $\mathcal{P}_n := \EPP (\mathcal{S}, M_n (\C))$ in definiramo preslikavo 
    $$i: \mathcal{S} \to \bigoplus_{n \in \N} M_n (\C) ^{\oplus \mathcal{P}_n},\quad \varphi = \bigoplus_{n \in \N} \bigoplus_{\psi \in \mathcal{P}_n} \psi.$$
    Označimo $C^*$-algebro $$\mathfrak{A} := \bigoplus_{n \in \N} M_n (\C) ^{\oplus \mathcal{P}_n}.$$
    Očitno je $i(e) = 1 \in \mathfrak{A}$.
    
    Če želimo dokazati, da je $i$ izomorfizem abstraktnih operatorskih prostorov na svojo sliko, mora veljati, 
    da je $A \in \mathcal{C}_n$ natanko tedaj, ko je $i^{(n)} (A) \geq 0$ v $M_n(\mathfrak{A})$.
    Implikacija v desno je očitna po definiciji $i$, zato vzamemo $A_0 \notin \mathcal{C}_n$
    in dokazujemo, da $i^{(n)} (A_0)$ ni pozitivna matrika. 
    Dovolj je dokazati, da obstaja $\varphi \in \EPP (\mathcal{S}, M_k (\C))$,
    tako da $\varphi^{(n)} (A_0)$ ni pozitivna.
    
    Najprej bomo dokazali, da obstaja pozitiven linearen funkcional $s: M_n (\mathcal{S}) \to \C$,
    tako da $s(A_0) \not\geq 0$.
    Po Hahn-Banachovem separacijskem izreku (\cite[Theorem IV.3.9., str. 111]{conway1990}) obstaja omejen realen linearen funkcional
    $s_0: M_n (\mathcal{S}) \to \R$, ki loči zaprto konveksno množico $\mathcal{C}_n$ od $A_0$, torej 
    $$ s_0(A_0) < \inf \{ s_0(B)\ |\ B \in \mathcal{C}_n\}.$$
    Najprej opazimo, da mora veljati $s_0 (A_0) < 0$.
    Nadalje je $s_0 (B) \geq 0$ za vsak $B \in \mathcal{C}_n$. Če bi namreč obstajal $B \in \mathcal{C}_n$,
    tako da je $s_0 (B) < 0$, potem je $\frac{s_0 (A_0)}{s_0(B)} > 0$ in zato $\frac{s_0 (A_0)}{s_0(B)} B \in \mathcal{C}_n$.
    To pa bi pomenilo, da je 
    $$s_0 (A_0) < s_0 \left(\frac{s_0 (A_0)}{s_0(B)} B\right) = s_0 (A_0),$$ kar vodi v protislovje.
    Sedaj pa lahko definiramo želen kompleksen funkcional
    $$s : M_n (\mathcal{S}) \to \C,\quad s(A) = s_0 (\real A) + i s_0 (\imag A).$$

    Ker je $s: M_n (\mathcal{S}) \to \C$ pozitiven funkcional, je $\varphi_s : \mathcal{S} \to M_n (\C)$ po prejšnji trditvi
    povsem pozitivna preslikava. Ker za $$e = \begin{bmatrix}
        e_1 \\ \vdots \\ e_n
    \end{bmatrix} \in \C^{n^2}$$ velja 
    $$e^* \varphi_s ^{(n)} (A_0) e = s_{\varphi_s} (A_0)= s(A_0) \not\geq 0,$$
    $\varphi_s ^{(n)} (A_0)$ ni pozitivna matrika. Če je $\varphi_s(e) = 1$, potem lahko nastavimo $\varphi := \varphi_s$ in smo končali.

    Oglejmo si pozitivno matriko $\varphi_s (e) = P = B^* B \in M_n (\mathcal{S})$. 
    Če je $P$ obrnljiva, potem lahko definiramo $\varphi := (B^{-1})^* \varphi_s B^{-1} \in \EPP(\mathcal{S}, M_n (\C))$.
    V primeru, ko $P$ ni obrnljiva, vzamemo nek $a \in \mathcal{S}$ z normo $\| a\| \leq 1$.
    Od tod sledi 
    $$\begin{bmatrix}
        e & a\\
        a ^* & e
    \end{bmatrix} \in \mathcal{C}_2$$
    in posledično 
    $$\begin{bmatrix}
        P & \varphi_s (a)\\
        \varphi_s (a) ^* & P
    \end{bmatrix}\geq 0.$$
    Če je $P x = 0$, potem je po lemi \ref{lem:random1} 
    $$\langle \varphi_s (a) x, \varphi_s (a) x \rangle ^2 \leq \langle \varphi_s (a) x, \varphi_s (a) x \rangle \cdot \langle P x, P x \rangle = 0,$$
    torej je $\varphi_s (a) x = 0$. Od tod sledi $ \ker P \subseteq \ker \varphi_s (a)$ in zato $ \ker \varphi_s (a) ^\perp \subseteq\ker P^\perp$.
    Naj bo $Q \in M_n (\C)$ ortogonalna projekcija na $\ker P^\perp$, torej $Q \varphi_s (A) Q = \varphi_s (A)$.
    Denimo, da je $Q$ matrika ranga $k < n$ (in zato tudi $P$ ranga $k$). Ker je $P \geq 0$, obstaja po spektralnem izreku 
    unitarna matrika $U \in M_n (\C)$, tako da je 
    $$P = U^* \begin{bmatrix}
        \diag (\lambda_1, \dots, \lambda_k) &\\
        & 0
    \end{bmatrix} U,\quad \lambda_1 \geq \cdots \geq \lambda_k > 0.$$
    Definiramo 
    $$V := U^* \begin{bmatrix}
        \diag (\lambda_1 ^{-\frac{1}{2}}, \dots, \lambda_k ^{-\frac{1}{2}})\\ 
        0
    \end{bmatrix} \in \C^{n \times k},$$
    tako da je $V^* P V = I_k$. Nadalje definiramo 
    $$W := \begin{bmatrix}
        \diag (\lambda_1 ^{\frac{1}{2}}, \dots, \lambda_k ^{\frac{1}{2}}) & 0
    \end{bmatrix} U \in \C^{k \times n},$$
    tako da je 
    $$V W = U^* \begin{bmatrix}
        I_k & 0\\ 
        0 & 0
    \end{bmatrix} U = Q.$$
    Potem je $\varphi := V^* \varphi_s V \in \EPP (\mathcal{S}, M_k (\C))$ prelikava, za katero velja
    $$e^* (I_n \otimes W)^* \varphi^{(n)} (A_0) (I_n \otimes W) e = e^* (I_n \otimes Q)^* \varphi_s ^{(n)} (A_0) (I_n \otimes Q) e =e^* \varphi_s ^{(n)} (A_0) e < 0,$$
    torej $\varphi^{(n)} (A_0)$ ni pozitivna matrika.
\end{dokaz}

\section{$C^*$-ogrinjače}\label{sec:10}



Izrek \ref{izr:choi-effros} nam pove, da vsak abstrakten operatorski sistem $\mathcal{S}$ lahko `vložimo' 
v neko $C^*$-algebro $\mathfrak{A}$.
Zato se naravno pojavi vprašanje, ali obstaja $C^*$-algebra, 
v katero lahko vložimo abstrakten operatorski sistem $\mathcal{S}$, in ki je na nek način `minimalna'. 
V ta namen bomo za vsako vložitev $i: \mathcal{S} \to \mathfrak{A}$ zahtevali, da je $\mathfrak{A} = C^*(i(\mathcal{S}))$.

To vprašanje lahko ekvivalentno prevedemo na (konkretne) operatorske sisteme. 
Od sedaj naprej bomo delali v kategoriji $\os_1$, v kateri so izomorfizmi enotske povsem izometrične bijekcije med operatorskimi sistemomi.
Za konkreten operatorski sistem $\mathcal{S} \subseteq C^* (\mathcal{S})$ torej iščemo `minimalno' $C^*$-algebro $\mathfrak{A}$, za katero obstaja 
povsem izometrična enotska preslikava $i: \mathcal{S} \to \mathfrak{A}$, tako da je $\mathfrak{A} = C^*(i(\mathcal{S}))$.

\begin{definicija}
    \emph{$C^*$-pokritje} operatorskega sistema $\mathcal{S}$ je par $(i, \mathfrak{A})$,
    kjer je $i: \mathcal{S} \to \mathfrak{A}$ enotska povsem izometrija in  
    $\mathfrak{A} = C^* (i(\mathcal{S}))$.

    Če sta $(i_1, \mathfrak{A}_1)$ in $(i_2, \mathfrak{A}_2)$
    $C^*$-pokritji $\mathcal{S}$ in obstaja takšen $*$-homomorfizem $\pi: \mathfrak{A}_1 \to \mathfrak{A}_2$,
    da velja $\pi \circ i_1 = i_2$, potem je $\pi$ \emph{morfizem $C^*$-pokritij $\mathcal{S}$}.
\end{definicija}

\begin{opomba}
    Morfizem $C^*$-pokritij $(i_1, \mathfrak{A}_1)$ in $(i_2, \mathfrak{A}_2)$,
    če ta obstaja, je enolično določen na generatorjih $i_1 (\mathcal{S})$ $C^*$-algebre $\mathfrak{A}_1$,
    zato je enoličen.

    Ker $i_2(\mathcal{S})$ generira $C^*$-algebro $\mathfrak{A}_2$ in velja $i_2(\mathcal{S}) = \pi (i_1 (\mathcal{S})) \subseteq \pi(\mathfrak{A}_1)$,
    je $\pi(\mathfrak{A}_1) = \mathfrak{A}_2$, zato je $\pi: \mathfrak{A}_1 \to \mathfrak{A}_2$ surjektivna.
\end{opomba}


\begin{definicija}
    \emph{$C^*$-ogrinjačo} operatorske algebre $\mathcal{S}$ definiramo kot končni objekt $(i, C_{\env} ^* (\mathcal{S}))$ v kategoriji $C^*$-pokritij $\mathcal{S}$:
    če je $(j, \mathfrak{A})$ poljubno $C^*$-pokritje $\mathcal{S}$, potem obstaja $*$-homomorfizem $\pi: \mathfrak{A} \to C_{\env} ^* (\mathcal{S})$,
    tako da je $\pi \circ j = i$.
    \begin{equation*}
        \begin{tikzpicture}[>=stealth]

        \node (S) at (0,2) {$\mathcal{S}$};
        \node (A) at (4,2) {$\mathfrak{A}$};
        \node (Q) at (4,0) {$C^* _{\env} (\mathcal{S})$};
        
        % top arrow (unlabeled)
        \draw[->]  (S) -- node[above] {$j$} (A);
        
        % left diagonal arrow
        \draw[->] (S) -- node[below left] {$i$} (Q);
        
        % right vertical arrow
        \draw[->] (A) -- node[right] {$\pi$} (Q);
        
        \end{tikzpicture}    
    \end{equation*}
\end{definicija}

V tem podrazdelku dokažemo, da $C^*$-ogrinjača res obstaja za vsak operatorski sistem $\mathcal{S} \subseteq C^*(\mathcal{S})$.
Pri tem sledimo argumentom Dritschel-McCollougha (\cite{dritschel2005}) in Arvesona (\cite{arveson2003}).
Začnemo z naslednjo definicijo.

\begin{definicija}
    Naj bo $\varphi: \mathcal{S} \to \bh$ enotska povsem pozitivna preslikava.
    Pravimo, da je neka druga enotska povsem pozitivna preslikava $\psi: \mathcal{S} \to \mathcal{B}(\mathcal{K})$ njena \emph{dilatacija},
    če velja $\mathcal{H} \subseteq \mathcal{K}$ ter 
    $$P_{\mathcal{H} } \psi (a) |_{\mathcal{H}} = \varphi (a),\quad \forall a \in \mathcal{S}.$$
    Rečemo tudi, da je $\varphi$ \emph{kompresija} za $\psi$. To označujemo s $\varphi \preceq \psi$.

    Če je $ \overline{C^* (\psi (\mathcal{S})) \mathcal{H}} = \mathcal{K}$,
    potem pravimo, da je $\psi$ \emph{nerazcepna} dilatacija $\varphi$.
\end{definicija}

\begin{opomba}
    Relacija $\preceq$ je delna urejenost na enotskih povsem pozitivnih prelikavah iz $\mathcal{S}$ v algebre omejenih operatorjev na Hilbertovih prostorih.
\end{opomba}

\begin{zgled}
    Vsaka enotska povsem pozitivna preslikava $\varphi: \mathcal{S} \to \bh$ ima trivialno dilatacijo $\varphi \oplus \psi$,
    kjer je $\psi: \mathcal{S} \to \mathcal{B}(\mathcal{K})$ poljubna enotska povsem pozitivna preslikava.
\end{zgled}

\begin{zgled}
    Če je $\varphi: \mathfrak{A}  \to \bh$ enotska povsem pozitivna preslikava, 
    potem po Stinespringovem izreku obstaja $*$-homomorfizem $\pi: \mathfrak{A} \to \mathcal{B}(\mathcal{K})$,
    ki je dilatacija $\varphi$.
\end{zgled}

\begin{definicija}
    Enotska povsem pozitivna preslikava $\varphi: \mathcal{S} \to \bh$ je \emph{maksimalna},
    če nima nobene netrivialne dilatacije. Ekvivalentno: če je $\psi \succeq \varphi$,
    potem je $\psi = \varphi \oplus \psi'$ za neko enotsko povsem pozitivno preslikavo $\psi'$.
\end{definicija}

\begin{lema}
    Enotska povsem pozitivna preslikava $\varphi: \mathcal{S} \to \bh$ je maksimalna natanko tedaj, ko je edina njena nerazcepna dilatacija
    $\varphi$ sama.
\end{lema}

\begin{dokaz}
    Denimo, da je $\varphi$ maksimalna in $\psi: \mathcal{S} \to \bk$ njena nerazcepna dilatacija.
    Potem je $\psi = \varphi \oplus \psi'$,
    kjer je $\psi': \mathcal{S} \to \mathcal{B}(\mathcal{K}')$ enotska povsem pozitivna in $\mathcal{K} = \mathcal{H} \oplus \mathcal{K}'$.
    Od tod sledi, da je $\psi (\mathcal{S}) \mathcal{H} = \varphi(\mathcal{S}) \mathcal{H} = \mathcal{H}$ in posledično 
    $\overline{C^* (\psi(\mathcal{S})) \mathcal{H}} = \mathcal{H}$.
    Ker pa je $\psi$ nerazcepna, imamo $\overline{C^* (\psi(\mathcal{S})) \mathcal{H}} = \mathcal{K}$. To pomeni, da je $\mathcal{K} = \mathcal{H}$
    in $\psi' = 0$, torej $\psi = \varphi$.
    
    V obratno smer predpostavimo, da je $\varphi$ sama sebi edina nerazcepna dilatacija.
    Vzemimo poljubno dilatacijo $\psi : \mathcal{S} \to \bh$, tako da $\varphi \neq \psi$.
    Definiramo Hilbertov podprostor $\mathcal{H}_1 := \overline{C^* (\psi (\mathcal{S})) \mathcal{H}}  \subseteq \mathcal{K}$.
    Potem je
    $P_{\mathcal{H}_1} \psi |_{\mathcal{H}_1}$ nerazcepna dilatacija $\varphi$, zato je $\mathcal{H}_1 = \mathcal{H}$.
    Od tod sledi, da je $\mathcal{H}$ invarianten prostor za $C^* (\psi (\mathcal{S}))$. 
    Fiksirajmo poljuben element $a \in \mathcal{S}$.
    Potem je $\mathcal{H}$ invarianten prostor za $\psi (a)$.
    Ker je tudi $a^* \in \mathcal{S}$, je $\mathcal{H}$ invarianten za $\psi (a)^* = \psi(a^*)$.
    To pomeni, da je tudi $\mathcal{H}^\perp$ invarianten prostor za $\psi (a)$,
    zato lahko zapišemo 
    $$\psi (a) = \underbrace{P_{\mathcal{H}} \psi (a) |_{\mathcal{H}}}_{= \varphi (a)} \oplus P_{\mathcal{H}^\perp} \psi (a) |_{\mathcal{H}^\perp}.$$
    Če definiramo $$\psi': \mathcal{S} \to \mathcal{B}(\mathcal{H}^\perp),\quad \psi' (a) = P_{\mathcal{H}^\perp} \psi (a) |_{\mathcal{H}^\perp},$$
    potem je $\psi = \varphi \oplus \psi'$.
\end{dokaz}

\begin{definicija}
    Naj bo $\varphi: \mathcal{S} \to \bh$ enotska povsem pozitivna.
    Pravimo, da je $\varphi$ \emph{maksimalna v točki $(a, x) \in \mathcal{S} \times \mathcal{H}$}, če za vsako njeno dilatacijo $\psi \succeq \varphi$
    velja $\psi (a) x = \varphi(a) x.$        
\end{definicija}

\begin{opomba}
    Opazimo, da če je enotska povsem pozitivna $\varphi$ maksimalna v $(a, x)$, potem je tudi vsaka njena dilatacija $\psi \succeq \varphi$
    maksimalna v $(a, x)$.
\end{opomba}

\begin{lema}
    Enotska povsem pozitivna preslikava $\varphi: \mathcal{S} \to \bh$ je maksimalna natanko tedaj, ko je maksimalna za vse $(a, x) \in \mathcal{S} \times \mathcal{H}$.
\end{lema}

\begin{dokaz}
    Denimo, da je $\varphi: \mathcal{S} \to \bh$ maksimalna.
    Če je $\psi \succeq \varphi$, potem je $\psi = \varphi \oplus \psi'$,
    torej $\psi (a) |_{\mathcal{H}} = \varphi(a)$ za vsak $a \in \mathcal{S}$.
    Od tod sledi $\psi (a) x = \varphi(a) x$ za vsak $a \in \mathcal{S}$ ter $x \in \mathcal{H}$.

    Obratno, denimo, da je $\varphi$ maksimalna za vse $(a, x) \in \mathcal{S} \times \mathcal{H}$.
    Naj bo $\psi: \mathcal{S} \to \mathcal{B}(\mathcal{K})$ dilatacija $\varphi$. 
    Potem je $\psi (\mathcal{S}) \mathcal{H} = \mathcal{S}$. Na enak način kot v prejšnji lemi lahko od tod zaključimo,
    da $$\psi': \mathcal{S} \to \mathcal{B}(\mathcal{H}^\perp),\quad \psi' (a) = P_{\mathcal{H}^\perp} \psi (a) |_{\mathcal{H}^\perp}$$
    zadošča $\psi = \varphi \oplus \psi'$.
\end{dokaz}

\begin{posledica}
    Naj bo $\{\varphi_i: \mathcal{S} \to \mathcal{B}(\mathcal{H}_i)\}_{i \in I}$ množica maksimalnih preslikav.
    Potem je tudi $\bigoplus_{i \in I} \varphi_i : \mathcal{S} \to \mathcal{B}\left(\bigoplus_{i \in I} \mathcal{H}_i\right)$
    maksimalna preslikava.
\end{posledica}

\begin{dokaz}
    Očitno je $\bigoplus_{i \in I} \varphi_i$ enotska povsem pozitivna preslikava.
    Vzemimo nek $a \in \mathcal{S}$ in $h = (h_i)_{i \in I} \in \bigoplus_{i \in I} \mathcal{H}_i$.
    Denimo, da je $\psi$ neka njena dilatacija. Potem je $\psi$ tudi dilatacija $\varphi_i$ za vsak $i \in I$.
    Od tod sledi $$\psi (a) h = \psi (a) (h_i)_{i \in I} = (\psi(a) h_i)_{i \in I} = (\varphi_i (a) h_i)_{i \in I} = \left(\bigoplus_{i \in I} \varphi_i \right) (a) h,$$
    torej je $\varphi$ maksimalna na $\mathcal{S} \times \left(\bigoplus_{i \in I} \mathcal{H}_i \right)$.
\end{dokaz}

\begin{lema}
    Enotska povsem pozitivna preslikava $\varphi: \mathcal{S} \to \bh$ je maksimalna v točki $(a, x)$ natanko tedaj, ko velja 
    $$\| \varphi (a) x \| = \sup_{\psi \succeq \varphi} \| \psi(a) x \|.$$
\end{lema}

\begin{dokaz}
    Vzemimo poljubno dilatacijo $\psi \succeq \varphi$.
    Potem je $$\| \psi(a) x \| \leq \| \varphi (a)x \| = \| P_\mathcal{H} \psi(a) x \|,$$
    od koder sledi $\psi(a) x = P_{\mathcal{H}} \psi(a) x = \varphi(a) x$.
\end{dokaz}

Ključna lastnost maksimalnih preslikav je, da so invariantne za operatorske sisteme.

\begin{trditev}\label{trd:2.2}
    Naj bosta $\mathcal{S}_1 \subseteq C^*(\mathcal{S}_1),  \mathcal{S}_2 \subseteq C^*(\mathcal{S}_2)$ operatorska sistema in $\theta: \mathcal{S}_1 \to \mathcal{S}_2$
    izomorfizem. Če je $\varphi_1: \mathcal{S}_1 \to \bh$ maksimalna preslikava, potem je maksimalna
    tudi $\varphi_2 := \varphi_1 \circ \theta^{-1} : \mathcal{S}_2 \to \bh$.
\end{trditev}

\begin{dokaz}
    Naj bo $\psi_2: \mathcal{S}_2 \to \mathcal{B}(\mathcal{K})$ nerazcepna dilatacija $\varphi_2$, torej je 
    $$\overline{C^* (\psi_2 (\mathcal{S}_2)) \mathcal{H}} = \mathcal{K}.$$
    Definirajmo $\psi_1 := \psi_2 \circ \theta: \mathcal{S}_1 \to \mathcal{B}(\mathcal{K})$, ki je očitno enotska povsem pozitivna preslikava.
    Sedaj za vsak $a \in \mathcal{S}_1$ velja 
    $$P_{\mathcal{H}} \psi_1 (a) |_{\mathcal{H}} = P_{\mathcal{H}} \psi_2 (\theta (a)) |_{\mathcal{H}} = \varphi_2 (\theta(a)) = \varphi_1 (a),$$
    torej je $\psi_1: \mathcal{S}_1 \to \mathcal{B}(\mathcal{K})$ dilatacija $\varphi_1$.
    Ker je $\psi_1 (\mathcal{S}_1) = \psi_2 (\theta(\mathcal{S})_1) = \psi_2 (\mathcal{S}_2)$, velja 
    $\overline{C^* (\psi_1 (\mathcal{S}_1)) \mathcal{H}} = \mathcal{K}$, zato je $\psi_1$ nerazcepna dilatacija $\varphi_1$. 
    Ker je $\varphi_1$ maksimalna, sledi $\psi_1 = \varphi_1$
    ter $\mathcal{K} = \mathcal{H}$. Od tod pa sledi $\psi_2 = \varphi_2$.
\end{dokaz}

Maksimalne preslikave lahko karakteriziramo še na en način, ki je ključen za dokaz obstoja $C^*$-ogrinjače.

\begin{definicija}
    Naj bo $\mathcal{S} \subseteq C^* (\mathcal{S})$ operatorski sistem.
    Enotska povsem pozitivna preslikava $\varphi: \mathcal{S} \to \bh$ ima \emph{lastnost enolične razširitve}, če velja: 
    \begin{enumerate}
        \item $\varphi$ ima enolično povsem pozitivno razširitev $\varphi ': C^* (\mathcal{S}) \to \bh$; 
        \item $\varphi'$ je $*$-homomorfizem.
    \end{enumerate}
\end{definicija}

\begin{trditev}\label{trd:2.1}
    Preslikava $\varphi: \mathcal{S} \to \bh$ je maksimalna natanko tedaj, ko ima lastnost enolične razširitve.
\end{trditev}

Za dokaz trditve potrebujemo naslednjo lemo.

\begin{lema}
    Preslikava $\varphi: \mathcal{S} \to \bh$ ima lastnost enolične razširitve natanko tedaj, ko 
    je vsaka njena povsem pozitivna razširitev $\widetilde{\varphi}: C^* (\mathcal{S}) \to \bh$
    multiplikativna.
\end{lema}

\begin{dokaz}
    Denimo, da ima $\varphi: \mathcal{S} \to \bh$ samo multiplikativne
    povsem pozitivne razširitve na $C^* (\mathcal{S})$. 
    Vzemimo torej eno tako razširitev $\varphi': C^* (\mathcal{S}) \to \bh$.
    Ker je $\varphi'$ hkrati multiplikativna in $*$-linearna (kar sledi iz povsem pozitivnosti),
    je upodobitev $C^* (\mathcal{S})$ na $\mathcal{H}$.
    Če je $\varphi''$ še ena povsem pozitivna razširitev $\varphi$ na $C^* (\mathcal{S})$,
    je prav tako upodobitev $C^* (\mathcal{S})$ na $\mathcal{H}$.
    Ker se ti upodobitvi $C^*$-algebre $C^* (\mathcal{S})$ ujemata na njeni množici generatorjev $\mathcal{S}$,
    se ujemata povsod in velja $\varphi' = \varphi''$.
\end{dokaz}


\begin{dokaz}[Dokaz trditve \ref{trd:2.1}]
    Denimo najprej, da je $\varphi: \mathcal{S} \to \bh$ maksimalna.
    Predpostavimo, da je $\varphi': C^* (\mathcal{S}) \to \bh$ neka njena povsem pozitivna razširitev.
    Po Stinespringovemu izreku obstaja Hilbertov prostor $\mathcal{K} \supseteq \mathcal{H}$ ter upodobitev $\pi: C^* (\mathcal{S}) \to \mathcal{B}(\mathcal{K})$, tako da je 
    $\mathcal{K} = \overline{\pi (C^* (\mathcal{S})) \mathcal{H}}$ in
    $$\varphi' (a) = P_{\mathcal{H}} \pi (a) |_{\mathcal{H}},\quad \forall a \in C^* (\mathcal{S}).$$ 
    Potem je $\pi |_{\mathcal{S}}: \mathcal{S} \to \mathcal{B}(\mathcal{K})$ dilatacija 
    $\varphi: \mathcal{S} \to \bh$ in 
    $$\mathcal{K} = \overline{\pi (C^* (\mathcal{S})) \mathcal{H}} = \overline{C^*( \pi (\mathcal{S})) \mathcal{H}},$$
    zato je po maksimalnosti $\mathcal{K} = \mathcal{H}$ ter $\pi|_{\mathcal{S}} =  \varphi$.
    Od tod sledi $\varphi' = \pi$ in $\varphi'$ je multiplikativna.
    
    Za obratno implikacijo predpostavimo, da ima 
    $\varphi$ lastnost enolične razširitve. Naj bo $\psi : \mathcal{S} \to \mathcal{B}(\mathcal{K})$ njena poljubna nerazcepna dilatacija, torej 
    $\overline{C^* (\psi(\mathcal{S})) \mathcal{H}} = \mathcal{K}$.
    Po Arvesonovem razširitvenem izreku lahko $\psi$ razširimo do povsem pozitivne preslikave 
    $\psi': C^* (\mathcal{S}) \to \mathcal{B}(\mathcal{K})$. Ker velja 
    $P_\mathcal{H} \psi' (a) |_{\mathcal{H}} = \varphi(a)$ za vsak $a \in \mathcal{S}$,
    je $P_\mathcal{H} \psi' |_{\mathcal{H}}: C^*(\mathcal{S}) \to \bh$ enotska povsem pozitivna razširitev $\varphi$.
    Ker ima $\varphi$ lastnost enolične preslikave, je $P_\mathcal{H} \psi' |_{\mathcal{H}}$ multiplikativna.
    Po Cauchy-Schwarzevi neenakosti za povsem pozitivne preslikave za poljuben 
    $a \in \mathcal{S}$ velja 
    $$P_{\mathcal{H}} \psi' (a^*) P_{\mathcal{H}} \psi' (a) P_{\mathcal{H}} = P_{\mathcal{H}} \psi' (a^* a) P_{\mathcal{H}} \geq P_{\mathcal{H}} \psi' (a^*) \psi' (a) P_{\mathcal{H}}.$$
    Od tod sledi 
    \begin{align*}
        ((1 - P_{\mathcal{H}}) \psi'(a) P_{\mathcal{H}})^* ((1 - P_{\mathcal{H}}) \psi'(a) P_{\mathcal{H}}) = P_{\mathcal{H}} \psi'(a)^* (1 - P_{\mathcal{H}}) \psi'(a) P_{\mathcal{H}} \leq 0,
    \end{align*}
    zato $((1 - P_{\mathcal{H}}) \psi'(a) P_{\mathcal{H}})^* ((1 - P_{\mathcal{H}}) \psi'(a) P_{\mathcal{H}}) = 0$.
    Posledično je $\| (1 - P_{\mathcal{H}}) \psi'(a) P_{\mathcal{H}}\| = 0$, torej 
    $(1 - P_{\mathcal{H}}) \psi'(a) P_{\mathcal{H}}= 0$. To pomeni, da je $\psi' (a) \mathcal{H} \subseteq \mathcal{H}$ 
    in $\psi (\mathcal{S}) \mathcal{H} \subseteq \mathcal{H}$.
    Od tod sledi $C^* (\psi (\mathcal{S})) \mathcal{H} \subseteq \mathcal{H}$ in zato $\mathcal{K} = \mathcal{H}$. S tem smo dokazali, da je $\psi = \varphi$.
\end{dokaz}

Do sedaj še nismo spoznali nobenega primera maksimalnih preslikav.
Njihov obstoj nam zagotavlja naslednji izrek.

\begin{izrek}[Arveson-Dritschel-McCollough]\label{izr:2.2}
    Vsako enotsko povsem pozitivno preslikavo $\varphi: \mathcal{S} \to \bh$ lahko dilatiramo do 
    maksimalne preslikave $\psi: \mathcal{S} \to \bk$.
\end{izrek}

Dokaz izreka \ref{izr:2.2} je tehničen in temelji na uporabi Zornove leme.
Začnimo z naslednjo lemo, ki je posplošitev argumenta v dokazu Arvesonovega izreka.

\begin{lema}\label{lem:2.1}
    Naj bo $\{\varphi_{\lambda} : \mathcal{S} \to \mathcal{B} (\mathcal{H}_\lambda)\}_{\lambda \in \Lambda}$ 
    usmerjena množica enotskih povsem pozitivnih preslikav z relacijo $\preceq$.
    Definirajmo $\mathcal{H}_\infty$ kot napolnitev prostora $\bigcup_{\lambda \in \Lambda} \mathcal{H}_\lambda$.
    Potem obstaja enolično določena preslikava $\varphi_{\infty} \in \EPP (\mathcal{S}, \mathcal{B} (\mathcal{H}_\infty))$,
    tako da je $\varphi_\infty \succeq \varphi_{\lambda}$ za vsak $\lambda \in \Lambda$. 
\end{lema}

\begin{dokaz}
    Za vsak $\lambda$ na $\mathcal{H}_\infty = \mathcal{H}_\lambda \oplus \mathcal{H}_\lambda ^\perp$ definiramo povsem pozitivno preslikavo 
    $$\psi_\lambda: \mathcal{S} \to \mathcal{B}(\mathcal{H}_\infty),\quad \psi_\lambda (a) = \begin{bmatrix}
        \varphi_\lambda (a) & 0\\
        0 & 0
    \end{bmatrix},$$
    torej je $\{\psi_\lambda\}_{\lambda \in \Lambda}$ posplošeno zaporedje v $\PP_1 (\mathcal{S}, \mathcal{B}(\mathcal{H}_\infty))$.
    Ker je slednji prostor kompakten, ima posplošeno zaporedje neko stekališče $\varphi_\infty \in \PP_1 (\mathcal{S}, \mathcal{B}(\mathcal{H}_\infty))$.
    Dokazujemo, da je $\varphi_\infty$ želena dilatacija. 

    Fiksirajmo nek $\lambda \in \Lambda$ in $a \in \mathcal{S}$.
    Po definiciji za vsak $\kappa \succeq \lambda$ velja $P_{\mathcal{H}_\lambda} \psi_\kappa (a) \big|_{\mathcal{H}_\lambda} = \varphi_\lambda (a)$.
    Po drugi strani pa ima $\{\psi_\kappa (a)\}_{\kappa \succeq \lambda}$ stekališče v $\varphi_\infty (a)$ v ŠOT,
    zato ima $\{P_{\mathcal{H}_\lambda} \psi_\kappa (a) |_{\mathcal{H}_\lambda}\}_{\kappa \succeq \lambda}$ stekališče v $P_{\mathcal{H}_\lambda} \varphi_\infty (a) |_{\mathcal{H}_\lambda}$.
    Od tod sledi $P_{\mathcal{H}_\lambda} \varphi_\infty(a) |_{\mathcal{H}_\lambda} = \varphi(a)$ in $\varphi_\infty \succeq \varphi$.

    Dokažimo enoličnost. Denimo, da je $\varphi_\infty ': \mathcal{S} \to \mathcal{B} (\mathcal{H}_\infty)$ še ena povsem pozitivna preslikava,
    tako da je $\varphi_\lambda \preceq \varphi_\infty$ za vsak $\lambda \in \Lambda$. Vzemimo poljubna $a \in \mathcal{S}$ in $x, y \in \bigcup_{\lambda \in \Lambda} \mathcal{H}_\lambda$.
    Denimo, da je $x \in \mathcal{H}_\mu$ in $y \in \mathcal{H}_\nu$. Potem obstaja $\lambda \in \Lambda$, tako da je $\varphi_\mu \preceq \varphi_\lambda$ in $\varphi_\nu \preceq \varphi_\lambda$,    
    torej $x, y \in \mathcal{H}_\lambda$. Od tod sledi
    \begin{align*}
        \langle \varphi_\infty ' (a)x, y \rangle &= \langle P_{\mathcal{H}_\lambda}\varphi_\infty ' (a) P_{\mathcal{H}_\lambda}x, y \rangle\\
        &= \langle \varphi_\lambda (a) x, y \rangle\\
        &= \langle P_{\mathcal{H}_\lambda}\varphi_\infty (a) P_{\mathcal{H}_\lambda}x, y \rangle\\
        &= \langle \varphi_\infty (a) x, y \rangle.
    \end{align*} 
    Ker je $\overline{\bigcup_{\lambda \in \Lambda} \mathcal{H}_\lambda} = H_\infty$, velja 
    $$\langle \varphi_\infty ' (a) x, y \rangle = \langle \varphi_\infty (a) x, y \rangle,\quad \forall x, y \in \mathcal{H}_\infty$$
    in posledično $\varphi_\infty (a) = \varphi_\infty ' (a)$ za poljuben $a \in \mathcal{S}$. 

    Preostane le še dokazati, da je $\varphi_\infty$ enotska preslikava. Kot v prejšnjem odstavku 
    za poljubna $x, y \in \bigcup_{\lambda \in \Lambda} \mathcal{H}_\lambda$ obstaja tak $\lambda \in \Lambda$, da sta $x, y \in \mathcal{H}_\lambda$.
    Potem je 
    \begin{align*}
        \langle \varphi_\infty (1)x, y \rangle &= \langle P_{\mathcal{H}_\lambda}\varphi_\infty (1) P_{\mathcal{H}_\lambda}x, y \rangle\\
        &= \langle \varphi_\lambda (1) x, y \rangle\\
        &= \langle x, y \rangle.
    \end{align*} 
    Od tod ponovno sledi, da mora veljati $\varphi_\infty (1) = 1$.
\end{dokaz}

\begin{trditev}\label{trd:2.3}
    Naj bo $\varphi: \mathcal{S} \to \bh$ enotska povsem pozitivna preslikava in $(a_0, x_0) \in \mathcal{S} \times \mathcal{H}$.
    Potem obstaja dilatacija $\psi \succeq \varphi$, ki je maksimalna na $(a_0, x_0)$.
\end{trditev}

\begin{dokaz}
    Najprej opazimo, da je 
    $$\sup \{\psi (a_0) x_0\ |\ \psi \succeq \varphi\} \leq \| a_0 \| \cdot \| x_0 \| < \infty.$$
    Induktivno konstruiramo zaporedje dilatacij $\varphi \preceq \varphi_1 \preceq \varphi_2 \preceq \cdots$ 
    na Hilbertovih prostorih $\mathcal{H} \subseteq \mathcal{H}_1 \subseteq \mathcal{H}_2 \subseteq \cdots$, tako da 
    $$\| \varphi_{n + 1} (a_0) x_0 \| \geq \sup_{\psi \succeq \varphi_n} \| \psi(a_0) x_0 \| - \frac{1}{n}.$$
    Naj bo $\mathcal{H}_\infty$ napolnitev prostora $\bigcup_{n \in \N} \mathcal{H}_n$ 
    in $\varphi_\infty: \mathcal{S} \to \mathcal{B} (\mathcal{H}_\infty)$ iz leme \ref{lem:2.1}.
    Preostane še pokazati, da je $\varphi_\infty$ maksimalna na $(a_0, x_0)$.
    Če je $\psi \succeq \varphi_\infty$, potem je tudi $\psi \succeq \varphi_n$ za vsak $n \in \N$ in zato 
    $$\| \varphi_\infty (a_0) x_0 \| \geq \| P_{\mathcal{H}_{n + 1}} \varphi_\infty (a_0) x_0 \| = \| \varphi_{n + 1} (a_0) x_0 \| \geq \| \psi (a_0) x_0 \| - \frac{1}{n}.$$
    Od tod sledi $\| \varphi_\infty (a_0) x_0 \| \geq \| \psi (a_0) x_0 \|$, kar smo želeli dokazati.
\end{dokaz}

\begin{opomba}
    Brez škode za splošnosti lahko vzamemo $\mathcal{H}_\infty = \mathcal{H} \oplus \C$,
    saj je $\varphi_\infty$, zožena na $\lin \{\mathcal{H}, \varphi_\infty (a_0)x_0 \}$, še vedno maksimalna v $(a_0, x_0)$.
\end{opomba}

\begin{dokaz}[Dokaz izreka \ref{izr:2.2}]
    Po izreku o dobri urejenosti na množici $\mathcal{S} \times \mathcal{H} = \{(a_\lambda, x_\lambda)\}_{\lambda \in \Lambda}$ obstaja dobra urejenost.
    \begin{enumerate}
        \item V prvem koraku za vsak $\lambda \in \Lambda$ konstruiramo preslikavo $\varphi_\lambda \in \EPP( \mathcal{S}, \mathcal{B} (\mathcal{H}_\lambda))$,
        ki je maksimalna na $\{(a_\mu, x_\mu) \ |\ \mu \preceq \lambda\}$
        in za katero velja $\varphi_\lambda \succeq \varphi_\mu$ za vse $\mu \preceq \lambda$.
        Dokaz tega koraka poteka s transfinitno indukcijo, kar pomeni, da moramo indukcijsko predpostavko 
        posebej preveriti na minimumu, naslednikih in limitnih točkah dobre urejenosti $\{(a_\lambda, x_\lambda)\}_{\lambda \in \Lambda}$.

        Označimo z $(a_0, x_0)$ minimum dobre urejenosti. Po trditvi \ref{trd:2.3} obstaja dilatacija $\varphi_0: \mathcal{S} \to \mathcal{B}(\mathcal{H}_0)$,
        tako da je $\varphi_0 \succeq \varphi$ maksimalna na $(a_0, x_0)$.    

        Denimo sedaj, da je $(a_\lambda, x_\lambda)$ naslednik nekega $x_\nu \in \mathcal{S} \times \mathcal{H}$.
        Po indukcijski predpostavki obstaja enotska povsem pozitivna preslikava $\varphi_\nu: \mathcal{S} \to \mathcal{B}(\mathcal{H}_\nu)$,
        tako da je $\varphi_\nu$ maksimalna na množici $\{(a_\mu, x_\mu) \ |\ \mu \preceq \nu\}$ in $\lambda_\nu \succeq \lambda_\mu$
        za vsak $\nu \preceq \mu$.
        Tedaj po trditvi \ref{trd:2.3} obstaja dilatacija $\varphi_\lambda \succeq \varphi_\nu$,
        tako da je $\varphi_\lambda$ maksimalna na $(a_\lambda, x_\lambda)$. Od tod sledi, da je 
        $\varphi_\lambda$ maksimalna na množici $\{(a_\mu, x_\mu) \ |\ \mu \preceq \lambda\}$ in $\varphi_\mu \preceq \varphi_\lambda$
        za vse $\mu \preceq \lambda$.

        Nazadnje obravnavamo primer, ko je $(a_\lambda, x_\lambda)$ limitna točka v dobri urejenosti.
        Potem je $\{\varphi_\mu\ |\ \mu \preceq \lambda\}$ usmerjena množica enotskih povsem pozitivnih preslikav.
        Označimo s $\mathcal{H}_\lambda$ napolnitev prostora $\bigcup_{\mu \preceq \lambda} \mathcal{H}_\mu$.
        Potem po lemi
        \ref{lem:2.1} obstaja $\varphi_\lambda: \mathcal{S} \to \mathcal{B} (\mathcal{H}_\infty)$, tako da je $\varphi_\lambda \succeq \varphi_\mu$
        za vsak $\mu \preceq \lambda$. Od tod sledi tudi, da je $\varphi_\lambda$ maksimalna na množici $\{(a_\mu, x_\mu)\ |\ \mu \preceq \lambda\}$,
        s čimer smo zaključili transfinitno indukcijo. 

        \item V drugem koraku dokazujemo, da lahko $\varphi$ dilatiramo do enotske povsem pozitivne preslikave $\varphi_1 : \mathcal{S} \to \mathcal{B} (\mathcal{H}_1)$,
        ki je maksimalna na $\mathcal{S} \times \mathcal{H}$. 
        Vzemimo dilatacije $\varphi_\lambda: \mathcal{S} \to \mathcal{B}(\mathcal{H}_\lambda)$, ki smo jih konstruirali v prvem koraku.
        Naj bo $\mathcal{H}_1$ napolnitev prostora $\bigcup_{\lambda} \mathcal{H}_\lambda$. 
        Po lemi \ref{lem:2.1} obstaja enotska povsem pozitivna preslikava $\varphi_1: \mathcal{S} \to \mathcal{B}(\mathcal{H}_1)$,
        za katero velja $\varphi_1 \succeq \varphi_\lambda$ za vsak $\lambda$. S tem smo konstruirali dilatacijo $\varphi_1 \succeq \varphi$,
        ki je maksimalna na vseh točkah $(a_\lambda, x_\lambda) \in \mathcal{S} \times \mathcal{H}$.

        \item V tretjem koraku uporabimo argument iz drugega koraka, da konstruiramo zaporedje dilatacij
        $\varphi_n: \mathcal{S} \to \mathcal{B}(\mathcal{H}_n)$,
        tako da je $\varphi_n \succeq \varphi_{n - 1}$ in $\varphi_n$ maksimalna na $\mathcal{S} \times \mathcal{H}_{n - 1}$.
        Nazadnje še enkrat uporabimo lemo \ref{lem:2.1}, tako da na napolnitvi $\mathcal{K}$ prostora $\bigcup_{n \in \N} \mathcal{H}_n$
        dobimo dilatacijo $\psi: \mathcal{S} \to \mathcal{B}(\mathcal{K})$, za katero velja $\psi \succeq \varphi_n$ za vsak $n \in \N$.
        Po konstrukciji je $\psi$ maksimalna na $\mathcal{S} \times \bigcup_{n \in \N} \mathcal{H}_n$.
        Ker je $\overline{\bigcup_{n \in \N} \mathcal{H}_n} = \mathcal{K}$, sledi, da je $\psi$ maksimalna tudi na $\mathcal{S} \times \mathcal{K}$. \qedhere
    \end{enumerate}
\end{dokaz}

Sedaj lahko dokažemo glavni rezultat tega razdelka.

\begin{izrek}\label{izr:2.3}
    Za vsak operatorski sistem $\mathcal{S} \subseteq C^*(\mathcal{S})$ obstaja $C^*$-ogrinjača $C^* _{\env} (\mathcal{S})$.
\end{izrek}

\begin{dokaz}
    V prvem koraku dokažemo, da obstaja neka enotska povsem izometrična preslikava $i: \mathcal{S} \to \bh$,
    ki je maksimalna. Iz GNS izreka vemo, da obstaja $*$-monomorfizem $\pi: C^*(\mathcal{S}) \to \mathcal{B}(\mathcal{H}_0)$,
    torej je $\pi|_{\mathcal{S}} \to \mathcal{B}(\mathcal{H}_0)$ enotska povsem izometrija. Po izreku 
    \ref{izr:2.2} obstaja maksimalna preslikava $i: \mathcal{S} \to \bh$, ki je dilatacija $\pi|_\mathcal{S}$.
    Vemo pa, da je dilatacija povsem izometrične preslikave prav tako povsem izometrija; za vsak $a \in \mathcal{S}$ je namreč 
    $$\| \pi^{(n)} (a)\| \leq \| i^{(n)} (a)\| \leq \| i \|_{\po} \cdot \| a\| = \| a\| = \| \pi^{(n)} (a)\|.$$

    Naj bo $\mathfrak{A} = C^* (i (\mathcal{S})) \subseteq \bh$.
    Dokazujemo, da je $C^*$-pokritje $i: \mathcal{S} \to \mathfrak{A}$ $C^*$-ogrinjača za $\mathcal{S}$.
    Vzemimo poljubno $C^*$-pokritje $j: \mathcal{S} \to \mathfrak{B}$. Ker je $i$ maksimalna preslikava,
    je po trditvi \ref{trd:2.2} tudi $i \circ j^{-1}: j(\mathcal{S}) \to \mathfrak{A} \subseteq \bh$ maksimalna preslikava.
    To pomeni, da ima slednja lastnost enolične razširitve, zato jo lahko razširimo 
    do $*$-homomorfizma $\pi: \mathfrak{B} \to \bh$, za katerega po definiciji velja $\pi \circ j ({\mathcal{S}}) = i(\mathcal{S})$. 
    Ker je $\pi$ $*$-homomorfizem, lahko zaključimo 
    $$\pi (\mathfrak{B}) = \pi (C^* (j(\mathcal{S}))) = C^* (\pi (j(\mathcal{S}))) = C^* (i(\mathcal{S})) = \mathfrak{A}.$$
    Torej je $\pi: \mathfrak{B} \to \mathfrak{A}$ $*$-homomorfizem iz univerzalne lastnosti za $C^*$-ogrinjačo, od koder sledi, da je $\mathfrak{A} = C^* _{\env} (\mathcal{S})$.
\end{dokaz}

\begin{opomba}
    V zgornjem dokazu smo pokazali še več: za poljubno povsem izometrično maksimalno preslikavo 
    $i: \mathcal{S} \to \bh$ je $i: \mathcal{S} \to C^*(i(\mathcal{S}))$ $C^*$-ogrinjača za $\mathcal{S}$.
\end{opomba}

Pomembna lastnost $C^*$-ogrinjače je, da je invariantna za izomorfne operatorske sisteme. 

\begin{trditev}
    Naj bosta $\mathcal{S}_1, \mathcal{S}_2$ operatorska sistema in $\varphi: \mathcal{S}_1 \to \mathcal{S}_2$ 
    njun izomorfizem. Naj bosta 
    $i_1: \mathcal{S}_1 \to C_{\env} ^* (\mathcal{S}_1)$, $i_2: \mathcal{S}_2 \to C_{\env} ^* (\mathcal{S}_2)$ 
    njuni $C^*$-ogrinjači.
    Potem obstaja enoličen $*$-izomorfizem $\pi: C_{\env} ^* (\mathcal{S}_1) \to C_{\env} ^* (\mathcal{S}_2)$, 
    tako da je $\pi \circ i_1 = i_2 \circ \varphi$.
\end{trditev}

\begin{dokaz}
    Ker je $i_1 \circ \varphi^{-1}: \mathcal{S}_2 \to C_{\env} ^* (\mathcal{S}_1)$ $C^*$-pokritje,
    obstaja enolično določen $*$-homomorfizem $\pi: C_{\env} ^* (\mathcal{S}_1) \to C_{\env} ^* (\mathcal{S}_2)$,
    tako da je $\pi \circ (i_1 \circ \varphi^{-1}) = i_2$. Po istem razmisleku obstaja 
    enolično določen $*$-homomorfizem $\rho: C_{\env} ^* (\mathcal{S}_2) \to C_{\env} ^* (\mathcal{S}_1)$, tako da je 
    $\rho \circ (i_2 \circ \varphi) = i_1$. Potem je $\rho \circ \pi: C_{\env} ^* (\mathcal{S}_1) \to C_{\env} ^* (\mathcal{S}_1)$ $*$-homomorfizem,
    za katerega velja 
    \begin{align*}
        (\rho \circ \pi) \circ i_1 = \rho \circ (\pi \circ i_1 \circ \varphi^{-1}) \circ \varphi = \rho \circ i_2 \circ \varphi = i_1.
    \end{align*}
    Ker je $\id_{C_{\env} ^* (\mathcal{S}_1)}: C_{\env} ^* (\mathcal{S}_1) \to C_{\env} ^* (\mathcal{S}_1)$ še en tak $*$-homomorfizem,
    velja $\rho \circ \pi = \id_{C_{\env} ^* (\mathcal{S}_1)}$. 
    Po enakem razmisleku sledi tudi $\pi \circ \rho = \id_{C_{\env} ^* (\mathcal{S}_2)}$, s čimer smo dokazali trditev.
\end{dokaz}

\begin{comment}
\section{Nekomutativen rob Šilova in $C^*$-ogrinjača}



Zato se naravno pojavi vprašanje, 

Arveson (\cite{arveson1969}) je pokazal, da je problem obstoja $C^*$-ogrinjače tesno povezan s 
pojmom ideala Šilova za operatorske sisteme. Če ima $\mathcal{S}$ konkretno realizacijo kot sebiadjungiran enotski podprostor v $C^*$-algebri $\mathfrak{A}$,
ki je generirana s $\mathcal{S}$ (v nadaljevanju pišemo $\mathcal{S} \subseteq C^*(\mathcal{S})$),
potem je ideal Šilova zaprt ideal $J \subseteq C^*(\mathcal{S})$, ki `meri',
kako daleč je $C^* (\mathcal{S})$ od $C^*$-ogrinjače $C^* _{\env} (\mathcal{S})$.

Ideal Šilova je posplošitev pojma robu Šilova za funkcijske sisteme iz konveksne geometrije 
(funkcijski sistemi so operatorski sistemi, vloženi v komutativno $C^*$-algebro $C(X)$).
Zato bomo najprej diskusijo umestili v zgodovinski kontekst in na kratko obravnavali ta komutativen primer, nato pa 
ga posplošili na splošne (ne nujno komutativne) operatorske sisteme in pokazali, kako  
ideal Šilova vodi do $C^*$-ogrinjače. Obstoj $C^*$-ogrinjače in ideala Šilova 
sledi iz Arveson-Dritschell-McColloughovega izreka, ki ga bomo dokazali s pomočjo teorije povsem pozitivnih preslikav iz prejšnjega poglavja.
V zadnjem podrazdelku posplošimo še soroden pojem Choquetovega robu iz funkcijskih sistemov do splošnih operatorskih sistemov 
in dokažemo njegov obstoj s pomočjo Davidson-Kennedyjevega izreka.
\end{comment}

\section{Nekomutativen rob Šilova}\label{sec:11}

Naj bo $\mathcal{S} \subseteq C^*(\mathcal{S})$ operatorski sistem.
V tem podrazdelku vpeljemo ideal Šilova, kot ga je definiral Arveson v \cite{arveson1969}.
To nam bo omogočilo, da $C^*$-ogrinjačo $C^* _{\env}(\mathcal{S})$ za nekatere operatorske sisteme tudi eksplicitno določimo; ideal Šilova namreč `meri',
v kolikšni meri se $C^*(\mathcal{S})$ razlikuje od $C^* _{\env}(\mathcal{S})$.
Da bi motivirali definicijo in utemeljili terminologijo, si najprej oglejmo pojem roba Šilova za operatorske sisteme v $C(X)$ in ga nato prevedimo v jezik $C^*$-algeber.

Naj bo torej $\mathcal{S} \subseteq C(X)$ operatorski sistem, za katerega zahtevamo še, da \emph{loči točke na $X$}:
to pomeni, da za poljubni točki $x \neq y \in X$ obstaja $f \in \mathcal{S}$, tako da je $f(x) \neq f(y)$.
Po Stone-Weierstrassovem izreku (\cite[Theorem V.8.1., str. 145]{conway1990}) je to ekvivalentno $C^*(\mathcal{S}) = C(X)$.

\begin{definicija}
    Naj bo $\mathcal{S} \subseteq C(X)$ operatorski sistem, ki loči točke na $X$.
    Zaprta množica $Y \subseteq X$ je \emph{robna množica} za $\mathcal{S}$,
    če vsaka funkcija $f \in \mathcal{S}$
    na njej doseže svojo normo.
    Če obstaja minimalna robna množica za $\mathcal{S}$, ji pravimo \emph{rob Šilova}.
\end{definicija}

\begin{comment}
\begin{opomba}
    Rob Šilova izhaja iz preučevanja \emph{uniformiranih algeber}:
    to so enotske podalgebre (ne nujno zaprte za $*$) $\mathcal{A} \subseteq C(X)$, ki {ločijo točke na $X$}. 
    Znan primer je \emph{racionalna uniformirana algebra} $R(X) \subseteq C(X)$ za kompaktno množico $X \subseteq \C$,
    ki je definirana kot zaprtje vseh racionalnih funkcij, katerih poli ležijo izven $X$.
    Še en primer je \emph{diskna algebra} $\A (\D)$,
    ki vsebuje vse holomorfne funkcije na $\D$, ki se jih da zvezno razširiti na $\overline{\D}$.
    Teorija uniformiranih algeber je podrobneje motivirana ter obravnavana v \cite{gamelin1969}.
\end{opomba}
\end{comment}

\begin{zgled}\label{zgl:2.1}
    Oglejmo si operatorski sistem
    $$\mathcal{S} = \C[z] + \C[z]^* \subseteq C(\overline{\D}).$$
    Ker $\mathcal{S}$ vsebuje afine funkcije, loči točke na $\overline{\D}$.
    Eksplicitno lahko izračunamo, da je rob Šilova za $\mathcal{S}$
    enak krožnici $\T$.

    Najprej pokažemo, da je $\T \subseteq \overline{\D}$ res robna množica za $\mathcal{S}$.
    Če je $f = p + \overline{q} \in \mathcal{S}$,
    potem je $|f|^2 = (p + \overline{q})(\overline{p} + q)$ subharmonična funkcija na $\D$:
    \begin{align*}
        \Delta |f|^2 &= 4 \frac{\partial}{ \partial \overline{z}} \frac{\partial}{\partial z} (p + \overline{q})(\overline{p} + q)\\
        &= 4 \frac{\partial}{ \partial \overline{z}} \left(p' (\overline{p} + q) + (p + \overline{q}) q'\right)\\
        &= 4 (p' \overline{p'} + \overline{q'} q')\\
        &= 4 (|p'|^2 + |q'|^2) \geq 0.
    \end{align*}
    Po principu maksimuma za subharmonične funkcije (\cite[Theorem 2.3.1., str. 25]{ransford1995}) funkcija $|f|^2$ doseže svoj maksimum na $\T$.

    Sedaj dokažemo še, da je $\T$ minimalna robna množica za $\mathcal{S}$.
    Za vsak $\lambda \in \T$ funkcija $z + \lambda \in \mathcal{S}$ doseže maksimum natanko v $\lambda$,
    torej mora biti $\lambda$ element roba Šilova za $\mathcal{S}$.
\end{zgled}



\begin{comment}

Definiramo preslikavo $$\phi: X \to S(\mathcal{S}),\quad x \mapsto \ev_x,$$
kjer je $\ev_x$ stanje, definirano s predpisom $\ev_x (f) = f(x)$ za vsak $f \in \mathcal{S}$.
Operatorski prostor $\mathcal{S}$ loči točke na $X$, zato je $\phi$ injektivna.
Dokažimo še, da je zvezna. Če je $\{x_\lambda\}_\lambda$ mreža v $X$,
ki konvergira k $x \in X$, potem za vsako funckijo $f \in \mathcal{S}$ velja
$$\phi(x_\lambda) (f) = \ev_{x_\lambda} (f) = f(x_\lambda) \to f(x) = \ev_x (f) = \phi(x) (f),$$
torej $\phi(x_\lambda) \to \phi(x)$ v šibki-$*$ topologiji.
Nazadnje opazimo še, da $\phi$ slika iz kompaktnega v Hausdorffov topološki prostor,
zato je zaprta preslikava in v posebnem homeomorfizem na svojo sliko.

Opazimo, da je $S(\mathcal{S})$ konveksna množica. Velja še več; $S(\mathcal{S})$ zaprta konveksna ogrinjača $\phi(X)$.
Denimo, da to ni res, torej obstaja $\omega_0 \in S(\mathcal{S})$, ki ni vsebovan v $\overline{\co} (\phi(X))$.
Po Hahn-Banachovem separacijskem izreku (DOPOLNI) potem obstaja funkcija $f \in \mathcal{S}$,
tako da je $$\sup \{\real \omega(f)\ |\ \omega \in \overline{\co}(\phi(X))\} < \real \omega_0 (f).$$
Brez škode za splošnost predpostavimo, da je $f$ realna (sicer jo nadomestimo z $\real f$) in $f \geq 0$, sicer ji lahko prištejemo poljubno pozitivno konstanto.
Potem je
\begin{align*}
    \| f\| &= \sup \{\omega (f)\ |\ \omega \in \phi(X) \}\\
    &= \sup \{\omega(f)\ |\ \omega \in \overline{\co}(\phi(X))\}\\
    &< \sup \{\real \omega(f)\ |\ \omega \in S(\mathcal{S})\}\\
    &\leq \sup \{\real \omega(f)\ |\ \omega \in C(X)^*,\ \| \omega\| = 1\}\\
    &= \| f\|,
\end{align*}
kar pa vodi v protislovje. 

Ker je $S(\mathcal{S})$ konveksna in kompaktna množica v šibki-$*$ topologiji,
ima po Krein-Milmanovem izreku (DOPOLNI) neprazno množico ekstremnih točk $\ext S(\mathcal{S})$,
tako da velja $\overline{\co} (\ext S(\mathcal{S})) = S(\mathcal{S})$. 
Ekstremnim točkam $S(\mathcal{S})$ pravimo tudi \emph{čista stanja}.
Ker je $\phi(X)$ kompaktna množica,
za katero velja $\overline{\co}(\phi(X)) = S(\mathcal{S})$, mora po Milmanovem izreku (DOPOLNI) veljati $\ext S(\mathcal{S}) \subseteq \phi(X)$. 



\begin{dokaz}
    Najprej moramo dokazati, da je $\overline{\partial_C}$ robna množica za $\mathcal{S}$.
    Vzemimo poljubno funkcijo $f \in \mathcal{S}$ in dokažemo, da $|f|$ na $\overline{\partial_C}$
    doseže dvoj maksimum. Fiksiramo neko točko $x \in X$, kjer $|f|$ doseže svoj maksimum. Definiramo množico stanj
    $$K := \{\omega \in S(\mathcal{S})\ |\ \omega(f) = f(x)\}.$$
    Očitno je $K$ neprazna, saj vsebuje stanje $\ev_{x}$. 
    Prav tako ni težko preveriti, da je konveksna in zaprta v $S(\mathcal{S})$ v šibki-$*$ topologiji, torej je kompaktna.
    Po Krein-Milmanovem izreku ima $K$ neko ekstremno točko $\omega_0$.
    Ker je $K$ lice $S(\mathcal{S})$, je $\omega_0 \in \ext S(\mathcal{S})$.
    To pa pomeni, da je $\omega_0 = \ev_{x_0}$ za nek $x_0 \in \partial_C $.
    Ker je $|f (x_0)| = |\omega_0 (f)| = |f(x)| = \| f\|$, $|f|$ doseže svoj maksimum na $\partial_C $.

    Preostane nam še dokazati, da za vsako robno množico $Y \subseteq X$ velja $\partial_C \subseteq Y$.
    Dokazujemo s protislovjem, torej predpostavimo, da obstaja točka $x_0 \in \partial_C \setminus Y$.
    Posledično ima okolico $U \ni x_0$, tako da je $U \cap Y = \emptyset$. Dokazali bomo,
    da je $\sup_{X \setminus U} |f| < \sup_{U} |f|$, kar bo vodilo v protislovje.
    Ker je $\phi: X \to \phi(X)$ homeomorfizem, je $\phi(U)$ relativno odprta okolica za $\phi(x_0)$ v $\phi(Y)$ v šibki-$*$ topologiji.
    Po definiciji šibke-$*$ topologije obstajajo funkcije $f_1, \dots, f_n \in \mathcal{S}$ in $\varepsilon > 0$, tako da je 
    $$\cap_{i = 1} ^n \{\omega \in \mathcal{S}^*\ |\ |\real \omega (f_i) - \real f_i(x_0)| < \varepsilon\} \cap \phi(Y) \subseteq \phi(U).$$
    Vpeljimo kompaktne konveksne množice 
    $$K_i := \{\omega \in \mathcal{S}^*\ |\ \real \omega (f_i) \geq \varepsilon\} \cap S(\mathcal{S})$$
    in opazimo, da velja $\phi(X) \setminus \phi(U) \subseteq \bigcup_{i = 1} ^n K_i$.
    Potem je $K := \overline{\co}\left(\bigcup_{i = 1} ^n K_i \right)$ kompaktna konveksna množica, ki ne vsebuje $\phi(x_0)$.
    Ponovno uporabimo Hahn-Banachov separacijski izrek, da dobimo funkcijo $f \in \mathcal{S}$, tako da je $\sup_{\omega \in K} \real \omega (f) < \real f(x_0)$.
    Tudi tokrat lahko predpostavimo, da je $f \geq 0$, od koder sledi 
    \begin{align*}
        \sup_{x \in X \setminus U} f(x) &= \sup_{\omega \in \phi(X) \setminus \phi(U)} \omega (f)\\
        &\leq \sup_{\omega \in \bigcup_{i = 1} ^n K_i} \omega (f)\\
        &= \sup_{\omega \in K} \omega (f)\\
        &< f(x_0),
    \end{align*}
    protislovje.
\end{dokaz}

\begin{posledica}
    Rob Šilova za operatorski sistem $$\mathcal{S} = \C [z] + \C [z]^* \subseteq C(\overline{\D})$$
    iz zgleda \ref{zgl:2.1} je $\partial_S = \T$.
\end{posledica}

\end{comment}

V tem poglavju bomo pojem roba Šilova posplošili na
poljuben (ne nujno komutativen) operatorski sistem $\mathcal{S} \subseteq C^*(\mathcal{S})$.
Ker $C^*(\mathcal{S})$ v splošnem ni nujno oblike $C(X)$, ne moremo več govoriti o robu Šilova kot o topološkem podprostoru v $X$.
Opazimo pa lahko, da ima vsak zaprt podprostor $Y \subseteq X$ pripadajoč zaprt ideal 
$$I(Y) = \{f \in C(X)\ |\ f|_{Y} = 0\} \subseteq C(X).$$
Iz naslednje leme sledi, da je ta korespondenca bijektivna.
\begin{lema}\label{lem:aux}
    Preslikava
    \begin{align*}
        \{\text{zaprte podmnožice $X$}\} &\to \{ \text{zaprti ideali v $C(X)$} \}\\
        Y &\mapsto I(Y) 
\end{align*}
    je bijekcija.
\end{lema}

\begin{dokaz}
    Najprej dokažimo injektivnost. Denimo, da sta $Y_1, Y_2 \subseteq X$
    različni zaprti množici, torej obstaja $x \in Y_2 \setminus Y_1$.
    Po Uryssohnovi lemi obstaja $f \in C(X)$,
    tako da je $f(x) = 1$ in $f|_{Y_1} = 0$. 
    Potem je $f \in I(Y_1)$, toda $f \notin I(Y_2)$.

    Dokažimo še surjektivnost. Vzemimo nek poljuben zaprt ideal $I \lhd C(X)$ in definiramo 
    $Y := \bigcap_{f \in I} f^{-1} (0)$. Dokazujemo, da je $I(Y) = I$.
    Očitno je $I \subseteq I(Y)$, zato dokažimo, da velja $I \supseteq I(Y)$.
    Naj bo $f \in I(Y)$ in $U_0 = f^{-1} (-\varepsilon, \varepsilon) \supseteq Y$ odprta okolica.
    Za poljuben $x \notin Y$ obstaja nek $f_x \in I$, tako da je $f_x (x) \neq 0$.
    Definiramo $U_x := \{f \neq 0\}$, ki je očitno okolica točke $x \in X$.
    Potem je družina $U_0 \cup \{U_x\}_{x \in X}$ odprto pokritje $X$, zato obstaja končno podpokritje
    $$U_0 \cup U_{x_1} \cup \dots \cup U_{x_n}.$$ Ker je $X$ kompakten in Hausdorffov, obstaja 
    razčlenitev enote $\rho_0, \dots, \rho_n$, podrejena zgornjemu pokritju.
    Sedaj definiramo funkcijo $f_{\varepsilon} := f \cdot \sum_{i = 1} ^n \rho_i$, za katero velja $\| f - f_{\varepsilon}\| \leq \varepsilon$ in ki jo lahko zapišemo kot
    $$f_{\varepsilon} = \sum_{i = 1} ^n \left( \frac{f}{f_{x_i}} \rho_i \right) \cdot f_{x_i} \in I.$$
    Ker je ideal $I$ zaprt za supremum normo, sledi $f \in I$.     
\end{dokaz}

\begin{opomba}
    Na tej točki opomnimo, da je zaprt ideal $I \lhd \mathfrak{A}$ v $C^*$-algebri vedno zaprt za involucijo, od koder sledi,
    da je $\quot{\mathfrak{A}}{I}$ prav tako $C^*$-algebra in je kanonična kvocientna preslikava 
    $$q_I : \mathfrak{A} \to \quot{\mathfrak{A}}{I}$$
    $*$-homomorfizem (bralec lahko dokaze najde v \cite[Section VIII.4, str. 245-248]{conway1990}).    
\end{opomba}

Zaprta množica $Y \subseteq X$ je robna za operatorski sistem $\mathcal{S} \subseteq C(X) = C^*(\mathcal{S})$ natanko tedaj, ko je zožitev $*$-homomorfizma 
$$r_Y: C(X) \to C(Y),\quad f \mapsto f|_Y$$
na $\mathcal{S}$ izometrija in po posledici \ref{pos:aux} povsem izometrija.
Iz spodnjega komutativnega diagrama pa sledi, da je $r_Y |_{\mathcal{S}}: \mathcal{S} \to C(Y)$
povsem izometrija natanko tedaj, ko je $q_{I(Y)} |_{\mathcal{S}}: \mathcal{S} \to \quot{C(X)}{I(Y)}$ povsem izometrija.
\begin{equation}\label{eq:aux}
    \begin{tikzpicture}[>=stealth]

        \node (CX) at (0,2) {$C(X)$};
        \node (CY) at (4,2) {$C(Y)$};
        \node (Q)  at (0,0) {$\displaystyle \frac{C(X)}{I(Y)}.$};
        
        \draw[->] (CX) -- node[above] {$r_Y$} (CY);
        
        \draw[->] (CX) -- node[left] {$q_{I(Y)}$} (Q);
        
        \draw[->] (Q) -- node[below right] {$\cong$} (CY);
        
        \end{tikzpicture}    
\end{equation}
S tem smo dokazali naslednje.
\begin{trditev}\label{trd:aux}
    Naj bo $\mathcal{S} \subseteq C(X) = C^*(\mathcal{S})$ operatorski sistem.
    Zaprta množica $Y \subseteq X$ je robna za $\mathcal{S}$
    natanko tedaj, ko je zožitev kvocientne preslikave $q_{I(Y)}: C(X) \to \quot{C(X)}{I(Y)}$ na $\mathcal{S}$ povsem izometrija.
\end{trditev}

To vodi do naslednje definicije.

\begin{definicija}\label{def:2.1}
    Naj bo $\mathcal{S} \subseteq C^*(\mathcal{S})$ operatorski sistem.
    Zaprt ideal $I \subseteq C^*(\mathcal{S})$ je \emph{robni ideal} za $\mathcal{S}$,
    če je zožitev kanonične kvocientne preslikave
    $q_I : C^*(\mathcal{S}) \to \quot{C^*(\mathcal{S})}{I}$
    na $\mathcal{S}$ povsem izometrija.
    Robnemu idealu pravimo \emph{ideal Šilova}, če vsebuje vse ostale robne ideale.
\end{definicija}

\begin{zgled}
    Če je $Y \subseteq X$ rob Šilova za operatorski sistem $\mathcal{S} \subseteq C(X) = C^*(\mathcal{S})$,
    potem je 
    $$I(Y) = \{f \in C(X)\ |\ f|_{Y} = 0\}$$
    njegov ideal Šilova.
\end{zgled}

Vsak robni ideal $I \subseteq C^*(\mathcal{S})$ inducira $C^*$-pokritje
$$q_I |_{\mathcal{S}}: \mathcal{S} \to \quot{C^*(\mathcal{S})}{I},$$
pri čemer je $q_I : \mathcal{C^*(\mathcal{S})} \to \quot{C^*(\mathcal{S})}{I}$ morfizem med $C^*$-pokritjema
$\mathcal{S} \hookrightarrow C^*(\mathcal{S})$ in $q_I |_{\mathcal{S}}$.
\begin{equation*}
    \begin{tikzpicture}[>=stealth, node distance=3cm]

        \node (S) at (0,2) {$\mathcal{S}$};
        \node (A) at (4,2) {$C^*(\mathcal{S})$};
        \node (Q) at (4,0) {$\frac{C^*(\mathcal{S})}{I}$.};
        
        % top arrow (unlabeled)
        \draw[{Hooks[right]}->]  (S) -- (A);
        
        % left diagonal arrow
        \draw[->] (S) -- node[left] {$q_I |_{\mathcal{S}}$} (Q);
        
        % right vertical arrow
        \draw[->] (A) -- node[right] {$q_I$} (Q);
        
    \end{tikzpicture}    
\end{equation*}
Ideal Šilova za $\mathcal{S}$ inducira najmanjše $C^*$-pokritje: $C^*$-ogrinjačo.

\begin{izrek}\label{izr:aux2}
    Za vsak operatorski sistem $\mathcal{S} \subseteq C^*(\mathcal{S})$ obstaja ideal Šilova $I \subseteq C^*(\mathcal{S})$
    in $q_I \big|_{\mathcal{S}}: \mathcal{S} \to \quot{C^*(\mathcal{S})}{I}$ je $C^*$-ogrinjača.
\end{izrek}

\begin{dokaz}
    Naj bo $i: \mathcal{S} \to C^* _{\env} (\mathcal{S})$ $C^*$-ogrinjača za $\mathcal{S}$.
    Po univerzalni lastnosti obstaja $*$-homomorfizem
    $\pi: C^*(\mathcal{S}) \to C_{\env} ^* (\mathcal{S})$, tako da je $\pi|_{\mathcal{S}} = i$.
    Dokazujemo, da je $I := \ker \pi$ ideal Šilova, torej da je zožitev kvocientne projekcije 
    $q_{I} : C^* (\mathcal{S}) \to \quot{C^* (\mathcal{S})}{I}$ na $\mathcal{S}$ povsem izometrija.
    Ker je $\pi: C^*(\mathcal{S}) \to C_{\env} ^* (\mathcal{S})$ surjektiven $*$-homomorfizem, je
    $$\overline{\pi}: \quot{C^*(\mathcal{S})}{I} \to C_{\env} ^* (\mathcal{S}),\quad a + \ker \pi \mapsto \pi (a)$$
    izometrični $*$-izomorfizem. Od tod sledi, da je 
    $$q_I |_{\mathcal{S}} = \overline{\pi}^{-1} \circ i: \mathcal{S} \to \quot{C^*(\mathcal{S})}{I}$$
    povsem izometrična preslikava, zato je $I$ robni ideal.
    
    Naj bo sedaj 
    $J \subseteq C^*(\mathcal{S})$ nek drug robni ideal in 
    $q_{J}: C^*(\mathcal{S}) \to \quot{C^*(\mathcal{S})}{J}$ njegova kanonična kvocientna projekcija,
    tako da je $q_J |_{\mathcal{S}}: \mathcal{S} \to \quot{C^*(\mathcal{S})}{J}$ $C^*$-pokritje za $\mathcal{S}$.
    Tedaj obstaja $*$-homomorfizem $\rho: \quot{C^*(\mathcal{S})}{J} \to C_{\env} ^* (\mathcal{S})$ iz univerzalne lastnosti, tako da velja 
    $\rho \circ q_J |_{\mathcal{S}} = i$. Po enoličnosti iz univerzalne lastnosti je $\rho \circ q_J = \pi$, zato 
    $J = \ker q_J \subseteq \ker \pi = I$.
\end{dokaz}


\begin{comment}
\begin{posledica}
    Za operatorski sistem $\mathcal{S} \subseteq M_n (C(X)) = C^*(\mathcal{S})$ 
    obstaja najmanjši zaprt podprostor $Y \subseteq X$, tako da vsaka matrika $A (x) = [a_{ij} (x)]_{i, j} \in \mathcal{S}$
    doseže svoj maksimum (v normi) na $Y$. 
\end{posledica}

\begin{dokaz}
    Znano dejstvo iz algebre je (DOPOLNI), da je za poljuben kolobar z enico $R$ vsak ideal $M_n (R)$
    oblike $M_n (I)$, kjer je $I$ ideal v $R$. 
    To pomeni, da so vsi zaprti ideali v $M_n (C(X))$ oblike $M_n (I(Y))$, kjer je $Y \subseteq X$ zaprta podmnožica.
    $$\begin{tikzpicture}[>=stealth]

        \node (CX) at (0,2) {$M_n (C(X))$};
        \node (CY) at (4,2) {$M_n (C(Y))$};
        \node (Q)  at (0,0) {$\displaystyle \frac{M_n(C(X))}{M_n (I(Y))}.$};
        
        \draw[->] (CX) -- node[above] {$r_Y$} (CY);
        
        \draw[->] (CX) -- node[left] {$q_{M_n (I(Y))}$} (Q);
        
        \draw[->] (Q) -- node[below right] {$\cong$} (CY);
        
        \end{tikzpicture}$$

    Po prejšnjem izreku za $\mathcal{S}$ obstaja ideal Šilova $I \subseteq M_n (C(X))$.

    To pomeni, da imamo bijektivno korespondenco med zaprtimi množicami v $X$ 
    in zaprtimi ideali v $M_n (C(X))$. V posebnem obstaja zaprta množica $Y \subseteq X$, tako da je 
    $I = M_n (C(Y))$. Od tod pa sledi, da je $Y$ najmanjša zaprta množica v $X$, tako da je 
    $$r_Y : M_n (C(X)) \to M_n (C(Y)),\quad [a_{ij}]_{i, j} \mapsto [a_{ij} |_Y ]_{i, j}$$
    izometrija.
\end{dokaz}
\end{comment}

\begin{posledica}[Šilov]
    Če je $\mathcal{S} \subseteq C(X) = C^*(\mathcal{S})$ operatorski sistem, potem obstaja njegov rob Šilova $Y$
    in $r_Y |_{\mathcal{S}}: \mathcal{S} \to C(Y)$ je njegova $C^*$-ogrinjača.
\end{posledica}

\begin{dokaz}
    Po izreku \ref{izr:2.3} obstaja ideal Šilova za $\mathcal{S}$, ki mora po lemi \ref{lem:aux} biti oblike 
    $$I(Y) = \{f \in C(X)\ |\ f|_Y = 0\}$$ za neko zaprto podmnožico $Y \subseteq X$. 
    Po trditvi \ref{trd:aux} je $Y$ rob Šilova za $\mathcal{S}$.    
    Zadnji stavek potem sledi iz izreka \ref{izr:aux2} in diagrama \eqref{eq:aux}.
\end{dokaz}

\begin{zgled}
    Oglejmo si ponovno operatorski sistem 
    $\mathcal{S} = \C[z] + \C[z]^* \subseteq C(\overline{\D})$
    iz zgleda \ref{zgl:2.1}. 
    Dokazali smo že, da je $\T$ njegov rob Šilova, zato je 
    $$I(\T) = \{f \in C(\overline{\D})\}$$ 
    njegov ideal Šilova ter
    $r_Y |_{\mathcal{S}}: \mathcal{S} \to C(\T)$ njegova $C^*$-ogrinjača.
\end{zgled}

Naslednji zgled je podrobneje obravnavan v \cite{connes2021}.

\begin{zgled}[Connes-van Suijlekom] \label{zgl:2.2}
    Označimo s $\mathcal{T}^{(n)}$ \emph{Toeplitzov operatorski sistem}, ki ga sestavljajo \emph{Toeplitzove matrike}  
    $$\begin{bmatrix}
        a_0 & a_{-1} & a_{-2} & \cdots & a_{-n}\\
        a_{1} & a_0 & a_{-1} & \ddots & \vdots\\
        a_{2} & a_1 & a_0 & \ddots & a_{-2}\\
        \vdots & \ddots & \ddots & \ddots & a_{-1}\\
        a_{n} & \cdots & a_{2} & a_{1} & a_0
    \end{bmatrix} \in M_n (\C)$$
    
    Najprej skicirajmo, zakaj velja $C^*(\mathcal{T}^{(n)}) = M_n (\C)$.
    Označimo matriko $S = \sum_{k = 1} ^n E_{k+1, k} \in M_n (\C)$, ki je Toeplitzova matrika za $a_{1} = 1$ in $a_k = 0, k \neq 1$.
    Potem je $S S^* = I - E_{11}$, zato je $E_{11} \in C^*(\mathcal{T}^{(n)})$.
    Z indukcijo dokažemo, da so vse matrike oblike $E_{kk}$ v $C^*(\mathcal{T}^{(n)})$:
    če je namreč $E_{k - 1, k - 1} \in C^*(\mathcal{T}^{(n)})$, potem enako velja tudi za $E_{k, k - 1} = S E_{k - 1, k - 1}$
    in $E_{k-1, k} = E_{k, k - 1}^*$ ter posledično $E_{kk} = E_{k, k - 1} E_{k - 1, k}$.
    Sedaj, ko imamo vse diagonalne matrične enote, dobimo tudi 
    $E_{k + i, k} = S^i E_{k, k}$ in $E_{k, k + i} = E_{k + i, k}^*$.

    Po izreku \ref{izr:2.3} potem obstaja njegov ideal Šilova $I \subseteq M_n (\C)$.
    Ker pa je $M_n (\C)$ enostaven kolobar, je $I = (0)$ in $C^* _{\env} (\mathcal{T}^{(n)}) = \quot{M_n (\C)}{I} = M_n (\C)$.
\end{zgled}

Toeplitzove matrike iz prejšnjega zgleda bomo v naslednjem zgledu posplošili do neskončno dimenzionalnih matrik nad Hilbertovim prostorom $\ell^2 (\Z_+)$.
Dokaze bomo izpustili in se zgolj sklicali na knjigo \cite{arveson2002}, v kateri je ta material podrobneje obravnavan.

\begin{zgled}\label{zgl:2.3}
    Oglejmo si operatorje $T \in \mathcal{B}(\ell^2 (\Z_+))$, ki imajo glede na standardno bazo $\{e_n\}_{n \in \Z}$ matriko oblike 
    \begin{equation}\label{eq:aux2}
        \begin{bmatrix}
        a_0 & a_{-1} & a_{-2} & \cdots \\
        a_{1} & a_0 & a_{-1} & \cdots \\
        a_{2} & a_1 & a_0 & \cdots \\
        \vdots & \vdots & \vdots & \ddots
    \end{bmatrix}\end{equation}
    Izkaže se, da za takšne operatorje obstaja zelo lepa karakterizacija s t.i.~\emph{Toeplitzovimi operatorji}.

    Začnemo s prostorom $L^2 (\T, m)$, kjer je $m$ normalizirana Lebesgueova mera na $\T$.
    Iz Fouriereve teorije vemo, da funkcije $\{z^n\}_{n \in \Z}$ tvorijo ortonormirano bazo Hilbertovega prostora $L^2 (\T, m)$.
    Sedaj definiramo \emph{Hardyjev prostor} $$H^2 := \overline{\lin\{1, z, z^2, \dots\}} \subseteq L^2 (\T, m).$$
    To je torej prostor $L^2$ funkcij $f$ na $\T$, ki imajo Fourierovo vrsto oblike 
    $$f(e^{i\theta}) \sim \sum_{n = 0} ^\infty a_n e^{i n \theta}.$$
    Ker funkcije $\{z^n\}_{n \in \Z_+}$ tvorijo ortonormirano bazo Hilbertovega prostora $H^2$,
    lahko $H^2$ na očiten način identificiramo s prostorom $\ell^2 (\Z_+)$. Nadalje lahko za vsako funkcijo $f \in L^\infty (\T, m)$
    definiramo operator 
    $$M_f \in \mathcal{B}(L^2 (\T, m)),\quad M_f (g) = fg$$
    in Toeplitzov operator 
    $$T_f \in \mathcal{B} (H^2),\quad T_f = P_{H^2} M_f |_{H^2}.$$
    Po izreku \cite[Theorem 4.2.4, str. 108]{arveson2002} so operatorji $T \in \mathcal{B}(\ell^2 (\Z_+))$ oblike \eqref{eq:aux2}
    natanko operatorji $T_f$ za $f \in L^\infty (\T, m)$, pri čemer je $\| T_f \| = \| f\|_\infty$.
    Z drugimi besedami: formalna Toeplitzova matrika 
    $$\begin{bmatrix}
        a_0 & a_{-1} & a_{-2} & \cdots \\
        a_{1} & a_0 & a_{-1} & \cdots \\
        a_{2} & a_1 & a_0 & \cdots \\
        \vdots & \vdots & \vdots & \ddots
    \end{bmatrix}$$
    je matrika omejenega operatorja natanko tedaj, ko obstaja funkcija $f \in L^\infty (\T, m)$,
    ki ima Fourierevo vrsto $f(e^{i\theta}) \sim \sum_{n = -\infty} ^\infty a_n e^{i n \theta}$.

    Posebno lepa teorija je za Toeplitzove operatorje \emph{z zveznim simbolom}:
    $$\mathcal{T}(C(\T)) := \{T_f\ |\ f \in C(\T)\}.$$
    Označimo $S := T_z \in \mathcal{T} (C(\T))$, ki ustreza operatorju desnega premika na $\ell^2 (\Z_+)$,
    torej $S e_n = e_{n + 1}.$ Operator $S$ generira algebro $C^*(S)$, ki je izjemnega pomena v funkcionalni analizi: rečemo ji \emph{Toeplitzova algebra}. 
    Izrek \cite[Theorem 4.3.2, str. 111]{arveson2002} nam pove, da je Toeplitzova algebra oblike 
    $$C^*(S) = \{T_f + K\ |\ f \in C(\T),\ {K \in \mathcal{K}(H^2)}\},$$
    pri čemer s $\mathcal{K}(H^2)$ označujemo kompaktne operatorje na $H^2$.
    Izrek nam pove še več: vsak element Toeplitzove algebre lahko na enoličen način razpišemo kot $T_f + K$.

    Sedaj lahko določimo ideal Šilova za $\mathcal{T}(C(\T))$. Najprej dokažimo,
    da je to sploh operatorski sistem. Hitro lahko preverimo, da velja 
    $$T_{\alpha f + \beta g} = \alpha T_f + \beta T_g,\quad T_{f} ^* = T_{\overline{f}},$$
    zato je $\mathcal{T}(C(\T))$ vektorski prostor, ki je zaprt za adjungacijo in vsebuje $1 = T_1.$
    Iz prejšnjega odstavka nato takoj sledi, da $\mathcal{T}(C(\T))$ generira Toeplitzovo algebro $C^*(S)$.
    Dokazujemo, da je njegov ideal Šilova kar $\mathcal{K}(H^2) \subseteq C^*(S)$.
    
    Znano je namreč, da kompaktni operatorji tvorijo zaprt ideal v algebri operatorjev na Hilbertovem prostoru.
    Za začetek dokažimo, da je $\mathcal{K}({H^2})$ robni ideal za $\mathcal{T}(C(\T))$.
    Ker ima vsak element v $C^*(S)$ enoličen zapis oblike $T_f + K$, lahko definiramo preslikavo
    $$\varphi: C^*(S) \to C(\T),\quad T_f + K \mapsto f.$$
    Enostavno je preveriti, da je $\varphi$ $*$-homomorfizem in $\ker \varphi = \mathcal{K}(H^2)$, zato imamo komutativni diagram  
    \begin{equation*}
        \begin{tikzpicture}[>=stealth, node distance=3cm]
    
            \node (S) at (0,2) {$C^*(S)$};
            \node (A) at (4,2) {$C(\T)$};
            \node (Q) at (0,0) {$\frac{C^*({S})}{\mathcal{K}(H^2)}$.};
            
            % top arrow (unlabeled)
            \draw[->]  (S) -- node[above] {$\varphi$} (A);
            
            % left diagonal arrow
            \draw[->] (S) -- node[left] {$q_{\mathcal{K}(H^2)}$} (Q);
            
            % right vertical arrow
            \draw[->] (Q) -- node[below right] {$\cong $} (A);
            
        \end{tikzpicture}    
    \end{equation*}
    Za vsak $f \in C(\T)$
    je $$\| \varphi(T_f) \| = \| f\| = \| T_f\|,$$
    zato je $\varphi|_{\mathcal{T}(C(\T))}: \mathcal{T}(C(\T)) \to C(\T)$ enotska izometrija in posledično povsem skrčitev.
    Enako velja tudi za njen inverz, zato je $\varphi|_{\mathcal{T}(C(\T))}: \mathcal{T}(C(\T)) \to C(\T)$ enotska povsem izometrija.
    Po zgornjem diagramu je tudi $q_{\mathcal{K}(H^2)} |_{\mathcal{T}(C(\T))}: \mathcal{T}(C(\T)) \to \quot{C^*({S})}{\mathcal{K}(H^2)}$
    povsem izometrija in $\mathcal{K}(H^2)$ je robni ideal. 
    Dokažimo še, da je tudi ideal Šilova.
    
    Vzemimo nek poljuben ideal $I \supseteq \mathcal{K}(H^2)$ v $C^*(S)$.
    Potem mora $I$ vsebovati nek Toeplitzov operator $T_f$, kjer je $f \in C(\T)$ neničelna funkcija.
    To pa pomeni, da $q_I |_{\mathcal{T}(C(\T))}: \mathcal{T}(C(\T)) \to \quot{C^*(S)}{I}$ uniči $T_f \neq 0$,
    zato ne more biti povsem izometrija.
\end{zgled}

\section{Nekomutativen Choquetov rob}\label{sec:12}

V tem poglavju definiramo Arvesonove robne upodobitve (\cite{arveson1969}), ki predstavljajo nekomutativno posplošitev točk v Choquetovem robu za operatorske sisteme v $C(X)$.
Dokazali bomo, da za vsak operatorski sistem $\mathcal{S}$ obstaja dovolj robnih upodobitev, da {povsem normirajo} $\mathcal{S}$.
Pri tem sledimo članku Davidsona in Kennedyja \cite{davidson2015},
v katerem avtorja posplošita argumente Arveson-Dritschel-McColloughovega izreka za obstoj čistih maksimalnih dilatacij.

Kot v prejšnjem poglavju tudi tokrat najprej začnemo z operatorskim sistemom $\mathcal{S} \subseteq C(X)$,
ki loči točke na $X$.
V tem kontekstu pravimo, da regularna Borelova verjetnostna mera (od sedaj naprej le: verjetnostna mera) $\mu$ na $X$ \emph{zastopa} točko $x \in X$,
če velja $$f(x) = \int_X f\, d\mu,\quad \forall f \in \mathcal{S}.$$
Takoj je očitno, da vsako točko $x \in X$ zastopa njena Diracova delta $\delta_x$.

\begin{definicija}
    Naj bo $\mathcal{S} \subseteq C(X)$ operatorski sistem, ki loči točke na $X$.
    \emph{Choquetov rob} $\mathcal{S}$ je množica vseh točk $x \in X$,
    za katere je Diracova delta $\delta_x$ edina verjetnostna mera, ki jih zastopa.
\end{definicija}


\begin{zgled}
    Vrnimo se k operatorskemu sistem
    $$\mathcal{S} = \C [z] + \C[z]^* \subseteq C(\overline{\D}).$$
    Dokazujemo, da je tudi njegov Choquetov rob enak $\T$.

    Vzemimo poljubno točko $\lambda \in \T$ in naj bo $\mu$ verjetnostna mera, ki jo zastopa.
    Potem je $$2 = \left| \int_{\overline{\D}} (z + \lambda)\, d\mu \right| \leq \int_{\overline{\D}} \left| z + \lambda\right| \, d\mu \leq 2,$$
    torej je $\int_{\overline{\D}} (2 - |z + \lambda|)\, d\mu = 0$ in $2 = | z + \lambda|$ $\mu$-skoraj povsod na $\overline{\D}$.
    Ker velja $2 - |z + \lambda| = 0$ samo za $z = \lambda$, to pomeni, da je $\mu (\{\lambda\}) = 1$.
    Slednje implicira, da je $\mu = \delta_\lambda$.

    Dokažimo še, da poljubna točka $\lambda \in {\D}$ ni v Choquetovem robu.    
    Definiramo Poissonovo jedro za $\overline{\D}$
    $$P(\lambda, z) = \frac{1 - |\lambda|^2}{|z - \lambda|^2},\quad z \in \T.$$
    Znano dejstvo je (glej na primer \cite[Theorem 15, str. 41]{evans2010}), da za vsako realno harmonično funkcijo $g$ na $\D$, ki je zvezna na $\overline{\D}$, 
    velja Poissonova reprezentacijska formula
    $$g(\lambda) = \int_{\T} g(z) P(\lambda, z)\, dm ,$$
    pri čemer je $m$ normalizirana Lebesgueova mera na $\T$. 
    Opazimo, da lahko vsako funkcijo $f \in \mathcal{S}$ zapišemo kot $f(x, y) = p(x, y) + i q(x, y)$,
    kjer sta $p, q$ realna polinoma (in zato harmonični funkciji na $\D$). To pomeni, da je 
    \begin{align*}
        f(\lambda) &= p(\lambda) + i q(\lambda) \\
        &= \int_{\T} p(z) P(\lambda, z)\, dm + i \int_{\T} q(z) P(\lambda, z)\, dm \\ 
        &= \int_{\T} (p(z) + i q(z)) P(\lambda, z)\, dm \\ 
        &= \int_{\T} f(z) P(\lambda, z)\, dm.
    \end{align*}
    Definiramo sedaj mero $\omega_\lambda$ na $\overline{\D}$,
    tako da je za vsako Borelovo množico $E \subseteq \overline{\D}$ 
    $$\omega_\lambda (E) := \int_{E \cap \T} P(\lambda, z)\, dm (z).$$ Potem je $\omega_\lambda$
    verjetnostna mera na $\overline{\D}$, tako da velja 
    $$f(\lambda) = \int_{\overline{\D}} f\, d\omega_\lambda$$
    za vse $f \in \mathcal{S}$. Ker je $\omega_\lambda \neq \delta_\lambda$,
    točka $\lambda$ ni v Choquetovem robu.

    Opazimo lahko, da rob Šilova in Choqueta za operatorski sistem $\mathcal{S}$ sovpadata. 
    Kot bomo videli kasneje, to ni naključje.
\end{zgled}

Če želimo definicijo Choquetovega robu prenesti na splošne operatorske sisteme,
jo moramo najprej prevesti v jezik funkcionalne analize.
Za poljubno točko $x \in X$ lahko definiramo $*$-homomorfizem 
$$\ev_x: C(X) \to \C,\quad f \mapsto f(x).$$
Ker po Rieszovem reprezentacijskem izreku vsakemu enotskemu pozitivnemu funkcionalu na $C(X)$
ustreza enolično določena verjetnostna mera na $X$, 
je vsaka točka $x \in X$ v Choquetovem robu natanko tedaj, ko ima funkcional $\ev_x |_{\mathcal{S}}: \mathcal{S} \to \C$
enolično pozitivno razširitev do $C(X)$.
Od tod takoj sledi naslednja trditev.

\begin{trditev}
    Choquetov rob $\mathcal{S}$ sestavljajo natanko točke $x \in X$, za katere ima $\ev_x |_{\mathcal{S}}$ lastnost enolične preslikave.
\end{trditev}


Preostane še najti primeren nekomutativen analog preslikavam $\ev_x :C(X) \to \C$ za $x \in X$.
Ta analog bodo \emph{nerazcepne upodobitve}. To so upodobitve $C^*$-algeber $\pi: \mathfrak{A} \to \bh$,
ki nimajo nobenega netrivialnega zaprtega invariantnega podprostora v $\mathcal{H}$.
Po standardnem izreku iz funkcionalne analize (\cite[Theorem 5.1.5., str. 143]{murphy1990}) je slednji pogoj ekvivalenten $\pi(\mathfrak{A}) ' = \C \cdot I_{\mathcal{H}}$.
Ni težko videti, da je preslikava $\ev_x: C(X) \to \C$ nerazcepna upodobitev $C^*$-algebre $C(X)$ za poljubno točko $x \in X$.
Velja pa tudi obratno.

\begin{trditev}
    Vsaka nerazcepna upodobitev $C(X)$ je oblike $\ev_x : C(X) \to \C$
    za neko točko $x \in X$.
\end{trditev}

\begin{dokaz}
    Naj bo $\rho: C(X) \to \bh$ nerazcepna upodobitev, torej $\rho(C(X))' = \C \cdot I_{\mathcal{H}}$.
    Ker je $C(X)$ komutativna algebra, je tudi $\rho(C(X))$ komutativna algebra in v posebnem velja 
    $\rho(C(X)) \subseteq \rho(C(X))' = \C \cdot I_{\mathcal{H}}$. To pomeni, da je vsak podprostor v $\mathcal{H}$
    invarianten za $\rho$, zato mora biti $\dim \mathcal{H} = 1$ in lahko vzamemo $\mathcal{H} \equiv \C$.
    To pomeni, da je $\rho: C(X) \to \mathcal{B}(\C) \equiv \C$ multiplikativen funkcional.

    Naprej dokazujemo, da je njegovo jedro $\ker \rho$ maksimalen ideal v $C(X)$. Vzemimo poljuben ideal $\ker \rho \subsetneqq I \subseteq C(X)$,
    tako da obstaja $f \in I \setminus \ker \rho$, za katerega velja $\rho(f) \neq 0$ in $1 - \frac{f}{\rho(f)} \in \ker \rho$.
    Od tod sledi $$1 = \left(1 - \frac{f}{\rho(f)}\right) + \frac{1}{\rho(f)} \cdot f \in I,$$
    zato je $I = C(X)$ in $\ker \rho$ je maksimalen (zaprt) ideal v $C(X)$.
    
    Po korespondenci med zaprtimi ideali v $C(X)$ in zaprtimi množicami v $X$ obstaja neka točka $x \in X$,
    tako da je $\ker \rho = I(\{x\})$. Potem za poljubno funkcijo $f \in C(X)$ velja
    $f - f(x) \cdot 1 \in \ker \rho$, od koder sledi 
    $$\rho(f) - f(x) = 0$$
    in $\rho = \ev_x$.
\end{dokaz}

Zadnji dve trditvi motivirata naslednjo definicijo.

\begin{definicija}
    \emph{Robna upodobitev} operatorskega sistema $\mathcal{S} \subseteq C^*(\mathcal{S})$ je nerazcepna upodobitev $\rho: C^*(\mathcal{S}) \to \bh$,
    tako da ima $\rho|_{\mathcal{S}}$ lastnost enolične razširitve.
\end{definicija}

Dokazali smo že, da so robne upodobitve za operatorski sistem $\mathcal{S} \subseteq C^*(\mathcal{S}) = C(X)$
sovpadajo s preslikavami $\ev_x$ za točke $x$ v Choquetovem robu. Vendar pa za splošen operatorski sistem
iz definicije ni takoj očitno, da robne upodobitve sploh obstajajo. V tem poglavju bomo dokazali, da za poljuben operatorski sistem $\mathcal{S} \subseteq C^*(\mathcal{S})$
obstaja dovolj robnih upodobitev, da \emph{povsem normirajo} $\mathcal{S}$: to pomeni,
da za vsako matriko $A \in M_n (\mathcal{S})$ obstaja robna upodobitev $\varphi$, tako da je $\| \varphi^{(n)} (A)\| = \| A\|$.

Najprej uvedemo nov pojem.

\begin{definicija}
    Preslikava $\varphi \in \PP (\mathcal{S}, \bh)$ je \emph{čista},
    če je vsak $\psi \in \PP (\mathcal{S}, \bh)$, za katerega velja $0 \leq \psi \leq \varphi$,
    oblike $\psi =  t \varphi$ za nek $0 \leq t \leq 1$.
\end{definicija}

\begin{zgled}
    Oglejmo si množico stanj 
    $$S(\mathcal{S}) = \{\omega \in \mathcal{S}^*\ |\ \text{$\omega$ pozitiven}, \| \omega\| = 1 \}$$ 
    na operatorskem sistemu $\mathcal{S}$. Potem je $S(\mathcal{S})$ neprazna, konveksna 
    in po Banach-Alaoglujevem izreku kompaktna v šibki-$*$ topologiji.
    Po Krein-Milmanovem izreku \cite[Theorem V.7.4., str. 142]{conway1990} ima zato neprazno množico ekstremnih točk, ki jo označujemo z $\ext S(\mathcal{S})$.
    Dokazujemo, da je stanje na $\mathcal{S}$ čisto natanko tedaj, ko je element $\ext S(\mathcal{S})$
    (slednje običajno vzamemo za definicijo čistih stanj).

    Denimo najprej, da stanje $\omega$ ni v $\ext S(\mathcal{S})$, torej obstajata $\omega_1, \omega_2 \in S(\mathcal{S})$, ki nista enaka $\omega$ in za katera velja
    $\omega = \frac{1}{2} \omega_1 + \frac{1}{2} \omega_2$. Potem je $0 \leq \frac{1}{2} \omega_1 \leq \omega$,
    vendar pa $\frac{1}{2} \omega_1$ ni skalarni večkratnik $\omega$. 
    
    V obratno smer predpostavimo, da obstaja pozitiven funkcional $\psi$,
    ki ni skalarni večkratnik $\omega$ in velja $0 \leq \psi \leq \omega$. Potem označimo $t = \psi(1)$
    in definiramo stanji $\omega_1 := \frac{1}{t} \psi$ ter $\omega_2 := \frac{1}{1 - t} (\omega - \psi)$, ki nista enaki $\omega$.
    Potem je $\omega = \frac{1}{2} \omega_1 + \frac{1}{2} \omega_2$, torej ne more biti čisto stanje.
\end{zgled}

Razlog za vpeljavo čistih preslikav je naslednja karakterizacija Choquetovega robu (\cite[Proposition 6.2, str. 29]{phelps2001}):
\emph{če je $\mathcal{S} \subseteq C(X)$ operatorski sistem, ki loči točke na $X$,
potem je točka $x \in X$ v Choquetovem robu natanko tedaj, ko je $\ev_x$ čisto stanje}.
V nadaljevanju bomo potrebovali različico te ekvivalence za splošne operatorske sisteme.

\begin{trditev}\label{trd:aux2}
    Naj bo $\mathcal{S} \subseteq C^* (\mathcal{S})$ operatorski sistem in $\rho: C^*(\mathcal{S}) \to \bh$
    upodobitev. Potem je $\rho$ robna upodobitev za $\mathcal{S}$ natanko tedaj, ko je $\varphi|_{\mathcal{S}}: \mathcal{S} \to \bh$
    čista in maksimalna.
%    Če je enotska povsem pozitivna preslikava $\varphi: \mathcal{S} \to \bh$ čista in maksimalna, potem je njena enolična razširitev 
%    $\rho: C^* (\mathcal{S}) \to \bh$ robna upodobitev.
\end{trditev}

Za dokaz potrebujemo naslednjo lemo.

\begin{lema}
    Naj bo $\mathfrak{A}$ $C^*$-algebra.
    Preslikava $\varphi \in \PP (\mathfrak{A}, \bh)$ je {čista} natanko tedaj, ko ima 
    njena minimalna Stinespringova upodobitev $(\pi, V, \mathcal{K})$ nerazcepno upodobitev $\pi$.
\end{lema}

\begin{dokaz}
    Po trditvi \ref{trd:1.4} imamo bijektivno korespondenco med množicama 
    $$\{\psi \in \PP (\mathfrak{A}, \bh)\ |\ 0 \leq \psi \leq \varphi\}$$
    ter $$\{T \in \pi(\mathfrak{A})'\ |\ 0 \leq T \leq I\}.$$
    Slednja množica vsebuje skalarje natanko tedaj, ko je $\pi(\mathfrak{A})' = \C \cdot I$.
\end{dokaz}

\begin{dokaz}[Dokaz trditve \ref{trd:aux2}]
    Najprej predpostavimo, da je $\rho|_{\mathcal{S}}$ čista in maksimalna.
    Dovolj je dokazati, da je $\rho$ nerazcepna upodobitev.
    Denimo nasprotno, torej obstaja nek Hilbertov podprostor $(0) \lneqq \mathcal{K} \lneqq \mathcal{H}$,
    ki reducira $\rho$. Naj bo $P \in \bh$ ortogonalna projekcija na $\mathcal{K}$, torej $P \in \rho(C^* (\mathcal{S}))'$.
    Potem je $\varphi := P \rho$ povsem pozitivna preslikava in $0 \leq \varphi|_{\mathcal{S}} \leq \rho|_{\mathcal{S}}$.
    Ker pa $\varphi (1) = P$ ni skalar, smo prišli v protislovje s tem, da je $\rho|_{\mathcal{S}}$ čista.

    Sedaj dokažimo še nasprotno. Denimo, da je $\rho: C^*(\mathcal{S}) \to \bh$ robna upodobitev.
    Ker je $\rho$ upodobitev, je sama svoja minimalna Stinespringova upodobitev, in ker je nerazcepna,
    je po prejšnji lemi tudi čista preslikava. Zadošča dokazati, da je potem tudi $\rho|_{\mathcal{S}}: \mathcal{S} \to \bh$
    čista preslikava. Denimo nasprotno, torej obstajata $\varphi_1, \varphi_2 \in \PP(\mathcal{S}, \bh)$, ki nista skalarna večkratnika $\rho|_{\mathcal{S}}$
    in za katera velja $\rho|_{\mathcal{S}} = \varphi_1 + \varphi_2$.
    Po Arvesonovem razširitvenem izreku ju lahko razširimo do povsem pozitivnih preslikav $\psi_1, \psi_2 \in \PP(C^*(\mathcal{S}), \bh)$.
    Ker ima $\rho|_{\mathcal{S}}$ lastnost enolične preslikave, je $\rho = \psi_1 + \psi_2$.
    Ker je $0 \leq \psi_1 \leq \rho$, smo v protislovju s tem, da je $\rho$ čista preslikava.
\end{dokaz}

Trditev \ref{trd:aux2} nam pove, da lahko vsako čisto maksimalno preslikavo na operatorskem sistemu $\mathcal{S}$ razširimo do robne upodobitve.
Preostane nam še, da konstruiramo čiste maksimalne preslikave. To bomo naredili na podoben način, kot v poglavju o $C^*$-algebrah,
v katerem smo konstruirali maksimalne dilatacije.
Naslednji izrek je analog Arveson-Dritschel-McColloughovega izreka za čiste preslikave.

\begin{izrek}[Davidson-Kennedy]\label{izr:3.1}
    Vsako čisto enotsko povsem pozitivno preslikavo $\varphi: \mathcal{S} \to \mathcal{B}(\mathcal{H})$ lahko dilatiramo do maksimalne čiste preslikave
    $\psi: \mathcal{S} \to \mathcal{B}(\mathcal{K})$, ki ima razširitev na robno upodobitev.
\end{izrek}

Pri dokazu bomo uporabili skoraj enake argumente kot pri dokazu izreka \ref{izr:2.2},
le da bodo sedaj dilatacije čiste. Spomnimo se leme \ref{lem:2.1}, ki pravi, da za vsako 
usmerjeno množico enotskih povsem pozitivnih preslikav $\{\varphi_{\lambda} : \mathcal{S} \to \mathcal{B} (\mathcal{H}_\lambda)\}_{\lambda \in \Lambda}$ 
z relacijo $\preceq$ obstaja enolično določena enotska povsem pozitivna preslikava $\varphi_{\infty}: \mathcal{S} \to \mathcal{B} (\mathcal{H}_\infty)$,
tako da je $\varphi_\infty \succeq \varphi_{\lambda}$ za vsak $\lambda \in \Lambda$. 

Naslednja lema je njen analog za čiste dilatacije.

\begin{lema}
    Naj bo $\{\varphi_{\lambda} : \mathcal{S} \to \mathcal{B} (\mathcal{H}_\lambda)\}_{\lambda \in \Lambda}$ 
    usmerjena množica čistih enotskih povsem pozitivnih preslikav z relacijo $\preceq$.
    Definirajmo $\mathcal{H}_\infty$ kot napolnitev prostora $\bigcup_{\lambda \in \Lambda} \mathcal{H}_\lambda$.
    Potem obstaja čista enolično določena preslikava $\varphi_{\infty} \in \EPP (\mathcal{S}, \mathcal{B} (\mathcal{H}_\infty))$,
    tako da je $\varphi_\infty \succeq \varphi_{\lambda}$ za vsak $\lambda \in \Lambda$. 
\end{lema}

\begin{dokaz}
    Obstoj in enoličnost dilatacije $\varphi_\infty$ že imamo po lemi \ref{lem:2.1}.
    Dokazati moramo le, da je $\varphi_\infty$ tudi čista preslikava.
    Vzemimo poljubno $\psi \in \PP (\mathcal{S}, \mathcal{B}(\mathcal{H}_\infty))$, 
    za katero velja $0 \leq \psi \leq \varphi_\infty$. Potem za vsak $\lambda \in \Lambda$ velja $0 \leq P_{\mathcal{H}_\lambda} \psi |_{\mathcal{H}_\lambda} \leq \varphi_\lambda$,
    zato je $P_{\mathcal{H}_\lambda} \psi |_{\mathcal{H}_\lambda} = t_\lambda \varphi_\lambda$ za skalarje $t_\lambda \in [0, 1]$.
    Ker je $P_{\mathcal{H}_\lambda} \psi(1) |_{\mathcal{H}_\lambda} = t_\lambda \cdot I_{\mathcal{H}_\lambda}$, so vsi skalarji $t_\lambda$ enaki, torej pišemo $t_\lambda = t$.
    Od tod sledi $P_{\mathcal{H}_\lambda} \psi |_{\mathcal{H}_\lambda} = t \varphi_\lambda$ za vsak $\lambda \in \Lambda$.
    Ker je $t \cdot \varphi_\infty \succeq t \cdot \varphi_\lambda$ in $\psi \succeq t \cdot \varphi_\lambda$ za vse $\lambda \in \Lambda$, od koder po enoličnosti v lemi \ref{lem:2.1}
    dobimo $\psi = t \cdot \varphi_\infty$. 
\end{dokaz}

Naslednja trditev je posplošitev trditve \ref{trd:2.3}.

\begin{trditev}\label{trd:3.1}
    Če je $\mathcal{S}$ operatorski sistem in $\varphi: \mathcal{S} \to \bh$ čista
    enotska povsem pozitivna preslikava, potem obstaja njena čista dilatacija $\psi \succeq \varphi$, ki je maksimalna na $(a_0, x_0) \in \mathcal{S} \times \mathcal{H}$.
\end{trditev}

\begin{dokaz}
    Predpostavimo, da $\varphi$ ni maksimalen v $(a_0, x_0)$, torej 
    $$\| \varphi (a_0) x_0\| < \sup_{\psi \succeq \varphi} \| \psi (a_0) x_0\| =: M.$$
    Definirajmo $C = \sqrt{M^2 - \| \varphi(a_0) x_0 \|^2}$. Po trditvi \ref{trd:2.3}
    obstaja dilatacija $\psi: \mathcal{S} \to \mathcal{B}(\mathcal{H} \oplus \C)$,
    ki je maksimalna na $(a_0, x_0)$. Brez škode za splošnost lahko predpostavimo, da zanjo velja $\psi(a_0)x_0 = (\varphi(a_0)x_0, C)$
    (sicer jo lahko konjugiramo z unitarno matriko $I \oplus \lambda$, tako da rotiramo $C$).
    
    Naj bo $$L = \{\psi: \mathcal{S} \to \mathcal{B}(\mathcal{H} \oplus \C) \text{ e.p.p.}\ |\ \psi \succeq \varphi,\ \psi(a_0)x_0 = (\varphi(a_0)x_0, C) \}$$
    in $$K = \{\psi: \mathcal{S} \to \mathcal{B}(\mathcal{H} \oplus \C) \text{ e.p.p.}\ |\ \psi \succeq \varphi \}.$$
    Po zgornjem razmisleku je $L$ neprazna množica, prav tako pa je konveksna in kompaktna v OŠ topologiji.
    Dokažimo, da je $L$ lice $K$. Predpostavimo, da imamo $\varphi_1, \varphi_2 \in K$
    ter $\psi = \frac{1}{2} (\varphi_1 + \varphi_2) \in L.$
    Naj bo $\varphi_k (a) h = (\varphi(a)h, x_k)$ za $k = 1, 2$,
    pri čemer je $|x_k| \leq C$. Potem je 
    $$C = \frac{x_1 + x_2}{2} = \left| \frac{x_1 + x_2}{2} \right| \leq \frac{1}{2} (|x_1| + |x_2|) \leq C,$$
    zato v trikotniški neenakosti $|x_1 + x_2| \leq |x_1 | + |x_2|$ velja enačaj. Od tod sledi $x_1 = x_2 = C$ in 
    $\varphi_1, \varphi_2 \in L$. 
    
    Po Krein-Milmanovem izreku ima $L$ neko ekstremno točko $\psi_0$.
    Dokazali bomo, da je $\psi_0$ čista preslikava. Naj bo $0 \leq \psi_1 \leq \psi_0$ in $\psi_2 := \psi_0 - \psi_1$.
    Brez škode za splošnost lahko predpostavimo, da je $\psi_1 (1) \geq \varepsilon I$ in $\psi_2 (1) \geq \varepsilon I$ za nek majhen $\varepsilon > 0$.
    V nasprotnem primeru bi lahko vzeli
    $\psi_k' = (1 - 2\varepsilon)\psi_k + \varepsilon \psi_0$ za $k = 1, 2$
    in če dokažemo, da velja 
    $\psi_1 ' = t \psi_0$ za $0 \leq t \leq 1$, potem je $\psi_1 = \frac{t - \varepsilon}{1 - 2\varepsilon} \psi_0$.
    Od tod sledi, da sta $\psi_1 (1), \psi_2 (1)$ obrnljiva. Opazimo tudi, da velja $P_{\mathcal{H}} \psi_k \big|_{\mathcal{H}} \leq \varphi$ za $k = 1, 2$,
    zato po predpostavki obstajata skalarja $0 \leq \lambda_1, \lambda_2 \leq 1$, tako da je 
    $P_{\mathcal{H}} \psi_k \big|_{\mathcal{H}} = \lambda_k \varphi$ za $k = 1, 2$.
    Prav tako velja $\lambda_1 + \lambda_2 = 1$ ter $\lambda_k \geq \varepsilon$ za $k = 1, 2$. 
    Potem obstajata vektorja $u_1, u_2 \in \mathcal{H}$ ter $\alpha_1, \alpha_2 >0$, tako da velja
    $$\psi_k (1) = \begin{bmatrix}
        \lambda_k & \lambda_k ^{\frac{1}{2}} u_k\\
        \lambda_k ^{\frac{1}{2}}  u_k ^* & \alpha_k 
    \end{bmatrix} = \underbrace{\begin{bmatrix}
        \lambda_k ^{\frac{1}{2}} & 0\\
        u_k ^* & \beta_k
    \end{bmatrix}}_{R_k ^*} \underbrace{\begin{bmatrix}
        \lambda_k ^{\frac{1}{2}} & u_k\\
        0 & \beta_k
    \end{bmatrix}}_{R_k},$$
    pri čemer je $\beta_k = \sqrt{\alpha_k - \| \xi \|^2} > 0$, saj sta $\psi_1 (1), \psi_2 (1)$ obrnljiva.
    Ker je $\psi_1 (1) + \psi_2 (1) = I$, mora veljati 
    $$\alpha_1 + \alpha_2 = 1,\quad \lambda_1 ^{\frac{1}{2}} u_1 + \lambda_2 ^{\frac{1}{2}} u_2 = 0.$$
    Prav tako je 
    $$R_k ^{-1} = \begin{bmatrix}
        \lambda_k ^{-\frac{1}{2}} & - \lambda_k ^{-\frac{1}{2}} \beta_k ^{-1} u_k\\
        0 & \beta_k ^{-1}
    \end{bmatrix}.$$
    Sedaj sta 
    $$\tau_k (a) = {R_{k} ^{-1}}^* \psi_k (a) R_k ^{-1} = \begin{bmatrix}
        \varphi_k (a) & \Gamma_k (a)\\
        \Delta_k (a) & \omega (a)
    \end{bmatrix},\quad k = 1, 2$$
    enotski povsem pozitivni preslikavi iz $\mathcal{S}$ v $\mathcal{B}(\mathcal{H} \oplus \C)$, ki dilatirata $\varphi$.

    V posebnem velja $\Gamma_1, \Gamma_2 \in \mathcal{B}(\mathcal{S}, \mathcal{H})$,
    $\Delta_1, \Delta_2 \in \mathcal{B}(\mathcal{S}, \mathcal{H}^*)$ in $\omega_1, \omega_2$ sta stanji na $\mathcal{S}$.
    Ker sta $\tau_1, \tau_2$ dilataciji $\varphi$, mora veljati $|\Gamma_1 (a) x_0|, |\Gamma_2 (a) x_0| \leq C$.
    Opazimo, da velja 
    \begin{align*}
        \psi_k (a) &= R_k ^* \tau_k (a) R_k\\ 
        &= \begin{bmatrix}
            \lambda_k ^{\frac{1}{2}} & 0\\
            u_k ^* & \beta_k
        \end{bmatrix} 
        \begin{bmatrix}
            \varphi_k (a) & \Gamma_k (a)\\
            \Delta_k (a) & \omega (a)
        \end{bmatrix}
        \begin{bmatrix}
            \lambda_k ^{\frac{1}{2}} & u_k\\
            0 & \beta_k
        \end{bmatrix}\\
        &= \begin{bmatrix}
            \lambda_k \varphi (a) & *\\
            \lambda_k ^{\frac{1}{2}} (u_k^* \varphi(a) + \beta_k \Delta_k (a)) & *
        \end{bmatrix}
    \end{align*}
    za $k = 1, 2$.
    Ker je $\psi_1 + \psi_2 = \psi_0$, je $\psi_1 (a_0) x_0 + \psi_2 (a_0) x_0 = (\varphi(a_0)x_0, C)$.
    Od tod sledi 
    \begin{align}
        C &= \lambda_1 ^{\frac{1}{2}} (u_1^* \varphi(a_0) + \beta_1 \Delta_1 (a_0))x_0 + \lambda_2 ^{\frac{1}{2}} (u_2^* \varphi(a_0) + \beta_2 \Delta_2 (a_0)) x_0 \nonumber\\
        &= \left| \lambda_1 ^{\frac{1}{2}} (u_1^* \varphi(a_0) + \beta_1 \Delta_1 (a_0)) x_0 + \lambda_2 ^{\frac{1}{2}} (u_2^* \varphi(a_0) + \beta_2 \Delta_2 (a_0)) x_0 \right| \nonumber\\
        &= \left| \langle \lambda_1 ^{\frac{1}{2}} u_1 + \lambda_2 ^{\frac{1}{2}} u_2, \varphi(a_0) \rangle x_0 + \lambda_1 ^{\frac{1}{2}}\beta_1 \Delta_1 (a_0) x_0 + \lambda_2 ^{\frac{1}{2}}\beta_2 \Delta_2 (a_0) x_0 \right| \nonumber\\
        &= \left| \lambda_1 ^{\frac{1}{2}}\beta_1 \Delta_1 (a_0) x_0 + \lambda_2 ^{\frac{1}{2}}\beta_2 \Delta_2 (a_0) x_0 \right| \nonumber\\
        &\leq \left| \lambda_1 ^{\frac{1}{2}}\beta_1 \Delta_1 (a_0) x_0 \right| + \left| \lambda_2 ^{\frac{1}{2}}\beta_2 \Delta_2 (a_0) x_0 \right| \label{eq:2.1}\\ 
        &\leq \lambda_1 ^{\frac{1}{2}}\beta_1  \left| \Delta_1 (a_0) x_0 \right| + \lambda_2 ^{\frac{1}{2}}\beta_2 \left| \Delta_2 (a_0) x_0 \right| \nonumber\\
        &\leq (\lambda_1 ^{\frac{1}{2}}\beta_1  + \lambda_2 ^{\frac{1}{2}}\beta_2) C \nonumber\\
        &\leq (\lambda_1 ^{\frac{1}{2}}\alpha_1 ^{\frac{1}{2}}  + \lambda_2 ^{\frac{1}{2}}\alpha_2 ^{\frac{1}{2}}) C \label{eq:2.2}  \\
        &\leq (\lambda_1 + \lambda_2)^{\frac{1}{2}} (\alpha_1 + \alpha_2) ^{\frac{1}{2}} C \label{eq:2.3} \\
        &= C. \nonumber
    \end{align}

    Najprej opazimo, da imamo v trikotniški neenakosti v vrstici \eqref{eq:2.1} enačaj in zato $\Delta_1 (a_0) x_0 = \Delta_2 (a_0) x_0 = C$, torej $\tau_1, \tau_2 \in L$.
    Prav tako imamo v Cauchyjevi neenakosti v vrstici \eqref{eq:2.3} enačaj, zato sta vektorja $(\lambda_1^{\frac{1}{2}}, \lambda_2 ^{\frac{1}{2}})$
    in $(\alpha_1^{\frac{1}{2}}, \alpha_2 ^{\frac{1}{2}})$ kolinearna. Ker sta oba tudi enotska, sta enaka in imamo $\lambda_1 = \alpha_1$ ter $\lambda_2 = \alpha_2$.
    Iz vrstice \eqref{eq:2.2} sledi, da je $\beta_1 = \sqrt{\alpha_1 - \| u_1 \|^2} = \alpha_1 ^{\frac{1}{2}}$ in 
    $\beta_2 = \sqrt{\alpha_2 - \| u_2 \|^2} = \alpha_2 ^{\frac{1}{2}}$, torej $\| u_1\| = \| u_2\| = 0$ in posledično $h_1 = h_2 = 0$.
    To pa pomeni, da je $R_k = \lambda_k ^{\frac{1}{2}} I$ za $k = 1, 2$, od koder sledi 
    $\psi_k = \lambda_k \tau_k $ za $k = 1, 2$. 
    S tem smo dokazali, da velja 
    $$\psi_0 = \lambda_1 \tau_1 + \lambda_2 \tau_2, \quad \lambda_1 + \lambda_2 = 1,$$
    pri čemer sta $\tau_1, \tau_2 \in L$.
    Ker je $\psi_0$ ekstremna točka $L$, mora veljati $\tau_1 = \tau_2 = \psi_0$,
    zato je $\psi_1 = \lambda_1 \psi_0$, kot smo hoteli dokazati.
\end{dokaz}

Sedaj lahko dokažemo izrek \ref{izr:3.1}. Pri tem skoraj popolnoma ponovimo dokaz izreka \ref{izr:2.2},
le da bodo tokrat vse dilatacije čiste.

\begin{dokaz}[Dokaz izreka \ref{izr:3.1}]
    Po izreku o dobri urejenosti na množici $\mathcal{S} \times \mathcal{H} = \{(a_\lambda, x_\lambda)\}_{\lambda \in \Lambda}$ obstaja dobra urejenost.
    \begin{enumerate}
        \item V prvem koraku za vsak $\lambda \in \Lambda$ konstruiramo čisto enotsko povsem pozitivno preslikavo $\varphi_\lambda: \mathcal{S} \to \mathcal{B} (\mathcal{H}_\lambda)$,
        ki je maksimalna na $\{(a_\mu, x_\mu) \ |\ \mu \preceq \lambda\}$
        in za katero velja $\varphi_\lambda \succeq \varphi_\mu$ za vse $\mu \preceq \lambda$.
        Dokaz tega koraka ponovno poteka s transfinitno indukcijo.

        Označimo z $(a_0, x_0)$ minimum dobre urejenosti. Po trditvi \ref{trd:3.1} obstaja čista dilatacija $\varphi_0: \mathcal{S} \to \mathcal{B}(\mathcal{H}_0)$,
        tako da je $\varphi_0 \succeq \varphi$ maksimalna na $(a_0, x_0)$.    

        Denimo sedaj, da je $(a_\lambda, x_\lambda)$ naslednik nekega $x_\nu \in \mathcal{S} \times \mathcal{H}$.
        Po indukcijski predpostavki obstaja čista enotska povsem pozitivna preslikava $\varphi_\nu: \mathcal{S} \to \mathcal{B}(\mathcal{H}_\nu)$,
        tako da je $\varphi_\nu$ maksimalna na množici $\{(a_\mu, x_\mu) \ |\ \mu \preceq \nu\}$ in $\lambda_\nu \succeq \lambda_\mu$
        za vsak $\nu \preceq \mu$.
        Tedaj po trditvi \ref{trd:3.1} obstaja čista dilatacija $\varphi_\lambda \succeq \varphi_\nu$,
        tako da je $\varphi_\lambda$ maksimalna na $(a_\lambda, x_\lambda)$. Od tod sledi, da je 
        $\varphi_\lambda$ čista in maksimalna na množici $\{(a_\mu, x_\mu) \ |\ \mu \preceq \lambda\}$ in $\varphi_\mu \preceq \varphi_\lambda$
        za vse $\mu \preceq \lambda$.

        Nazadnje obravnavamo primer, ko je $(a_\lambda, x_\lambda)$ limitna točka v dobri urejenosti.
        Potem je $\{\varphi_\mu\ |\ \mu \preceq \lambda\}$ usmerjena množica čistih enotskih povsem pozitivnih preslikav.
        Označimo s $\mathcal{H}_\lambda$ napolnitev prostora $\bigcup_{\mu \preceq \lambda} \mathcal{H}_\mu$.
        Potem obstaja čista dilatacija $\varphi_\lambda: \mathcal{S} \to \mathcal{B} (\mathcal{H}_\infty)$, tako da je $\varphi_\lambda \succeq \varphi_\mu$
        za vsak $\mu \preceq \lambda$. Od tod sledi tudi, da je $\varphi_\lambda$ maksimalna na množici $\{(a_\mu, x_\mu)\ |\ \mu \preceq \lambda\}$,
        s čimer smo zaključili transfinitno indukcijo. 

        \item V drugem koraku dokazujemo, da lahko $\varphi$ dilatiramo do čiste preslikave $\varphi_1 : \mathcal{S} \to \mathcal{B} (\mathcal{H}_1)$,
        ki je maksimalna na $\mathcal{S} \times \mathcal{H}$. 
        Vzemimo čiste dilatacije $\varphi_\lambda: \mathcal{S} \to \mathcal{B}(\mathcal{H}_\lambda)$, ki smo jih konstruirali v prvem koraku.
        Naj bo $\mathcal{H}_1$ napolnitev prostora $\bigcup_{\lambda} \mathcal{H}_\lambda$. 
        Potem obstaja čista enotska povsem pozitivna preslikava $\varphi_1: \mathcal{S} \to \mathcal{B}(\mathcal{H}_1)$,
        za katero velja $\varphi_1 \succeq \varphi_\lambda$ za vsak $\lambda$. S tem smo konstruirali čisto dilatacijo $\varphi_1 \succeq \varphi$,
        ki je maksimalna na vseh točkah $(a_\lambda, x_\lambda) \in \mathcal{S} \times \mathcal{H}$.

        \item V tretjem koraku uporabimo argument iz drugega koraka, da konstruiramo zaporedje čistih dilatacij
        $\varphi_n: \mathcal{S} \to \mathcal{B}(\mathcal{H}_n)$,
        tako da je $\varphi_n \succeq \varphi_{n - 1}$ in $\varphi_n$ maksimalna na $\mathcal{S} \times \mathcal{H}_{n - 1}$.
        Nazadnje na napolnitvi $\mathcal{K}$ prostora $\bigcup_{n \in \N} \mathcal{H}_n$
        dobimo čisto dilatacijo $\psi: \mathcal{S} \to \mathcal{B}(\mathcal{K})$, za katero velja $\psi \succeq \varphi_n$ za vsak $n \in \N$.
        Po konstrukciji je $\psi$ maksimalna na $\mathcal{S} \times \bigcup_{n \in \N} \mathcal{H}_n$.
        Ker je $\overline{\bigcup_{n \in \N} \mathcal{H}_n} = \mathcal{K}$, sledi, da je $\psi$ maksimalna tudi na $\mathcal{S} \times \mathcal{K}$. \qedhere
    \end{enumerate}
\end{dokaz}

S pomočjo Davidson-Kennedyjevega izreka se lahko lotimo glavnega izreka tega poglavja.

\begin{izrek}\label{izr:2.4}
    Za vsako matriko $A \in M_n (\mathcal{S})$ obstaja robna upodobitev $\rho: C^*(\mathcal{S}) \to \mathcal{B}(\mathcal{K})$,
    tako da je $\| A\| = \|\rho^{(n)} (A) \|$.
\end{izrek}

Za dokaz bomo potrebovali naslednji dve lemi.

\begin{lema}
    Če so $A, B, C, D \in \bh$, potem
    $$\left\| \begin{bmatrix}
        A & B\\
        C & D
    \end{bmatrix} \right\| \leq \left\| \begin{bmatrix}
        \| A\| & \| B\| \\
         \|C\| & \|D\|
    \end{bmatrix} \right\|.$$
\end{lema}

\begin{dokaz}
    Vzemimo enotski vektor $[x, y]^\top \in \mathcal{H}^2$. Potem je 
        \begin{align*}
            \left\|
            \begin{bmatrix}
                A & B\\
                C & D
            \end{bmatrix} \begin{bmatrix}
                x \\ y
            \end{bmatrix} \right\|^2 &= \left\| \begin{bmatrix}
                A x + B y\\  Cx + Dy
            \end{bmatrix} \right\|^2 \\
            &= \| A x + B y \|^2 + \| Cx + Dy \|^2 \\
            &\leq (\| A\| \| x\| + \| B\| \| y\|)^2 + (\| C\| \| x\| + \| D\| \| y\|)^2 \\
            &= \left\| \begin{bmatrix}
                \| A\| \| x\| + \| B\| \| y\| \\ \| C\| \| x\| + \| D\| \| y\|
            \end{bmatrix} \right\|^2\\
            &= \left\| \begin{bmatrix}
                \| A\| & \| B\|\\
                \| C\| & \| D\|
            \end{bmatrix} \begin{bmatrix}
                \| x\| \\ \| y\|
            \end{bmatrix} \right\|^2\\
            &\leq \left\| \begin{bmatrix}
                \| A\| & \| B\|\\
                \| C\| & \| D\|
            \end{bmatrix} \right\|^2. \qedhere
        \end{align*}
\end{dokaz}

\begin{lema}\label{lem:3.1}
    Naj bo $\mathfrak{A}$ $C^*$-algebra in $\pi: M_n (\mathfrak{A}) \to \bk$ upodobitev.
    Potem obstaja upodobitev $\rho: \mathfrak{A} \to \bh$, tako da je  
    $\pi$ unitarno ekvivalentna $\rho^{(n)}: M_n (\mathfrak{A}) \to \mathcal{B}(\mathcal{H}^n)$.
\end{lema} 

\begin{proof}
    Za poljuben par indeksov $i, j$ definirajmo $P_{ij} := \pi (E_{ij}) \in \bk$.
    Očitno velja $P_{ij} ^* = P_{ji}$ in 
    \begin{equation}\label{eq:1.7}
       P_{ij} P_{kl} = \delta_{jk} P_{il}. 
    \end{equation}
    V posebnem je 
    $$P_{ii} ^2 = P_{ii} = P_{ii}^*,$$
    torej je $P_{ii}$ ortogonalna projekcija na zaprt (po izreku o zaprtem grafu)
    podprostor $\mathcal{H}_i \subseteq \mathcal{K}$.
    Ker je $$P_{11} + \cdots + P_{nn} = I,\quad P_{ii} P_{jj} = 0,$$
    mora veljati $$\mathcal{H}_1 + \cdots + \mathcal{H}_n = \mathcal{K},\quad \mathcal{H}_i \perp \mathcal{H}_j$$
    in posledično $$\mathcal{H}_1 \oplus \cdots \oplus \mathcal{H}_n = \mathcal{K}.$$
    Za vsak $a \in \mathfrak{A}$ ter $P_{ii}$ velja $P_{ii} \pi(aI) = \pi(aI) P_{ii}$, torej je vsak $\mathcal{H}_i$
    reducirajoč prostor za $\pi(aI)$ in je 
    $$\pi_{ii}: \mathfrak{A} \to \mathcal{B}(\mathcal{H}_i),\quad \pi_{ii} (a) = \pi(aI) \big|_{\mathcal{H}_i}$$
    upodobitev.

    Po enačbi \eqref{eq:1.7} imamo $P_{ij} (\mathcal{H}_k) = \delta_{jk} \mathcal{H}_{i}$, 
    kar pomeni, da lahko $P_{ij}$ zapišemo v matrični obliki na $\mathcal{H}_1 \oplus \cdots \oplus \mathcal{H}_n$ kot
    $$P_{ij} = \begin{bmatrix}
    0 & 0 & 0 \\
    0  & 0 & \underbrace{P_{ij}\big|_{\mathcal{H}_j}}_{(i, j)} \\
     0 & 0 & 0 
 \end{bmatrix}.$$
    Iz \eqref{eq:1.7} prav tako sledi $$P_{ij} ^* P_{ij} = P_{ji} P_{ij} = P_{jj},$$
    od koder dobimo
    $$(P_{ij}\big|_{\mathcal{H}_j})^* (P_{ij}\big|_{\mathcal{H}_j}) = (P_{ji}\big|_{\mathcal{H}_i}) (P_{ij}\big|_{\mathcal{H}_j}) = P_{jj} \big|_{\mathcal{H}_j} = \id_{\mathcal{H}_j}.$$
    To implicira, da je $P_{ij}\big|_{\mathcal{H}_j}: \mathcal{H}_j \to \mathcal{H}_i$ unitaren operator
    in $$P_{ij} \pi_{jj}(a) P_{ji} = \pi_{ii} (a),\quad \forall a \in \mathfrak{A}.$$

    Sedaj definiramo $\mathcal{H} := \mathcal{H}_1$, $\rho := \pi_{11}$ in $U_i := P_{i1} \big|_{\mathcal{H}_1} : \mathcal{H}_1 \to\mathcal{H}_i$, zato je 
    $$U = \oplus_i ^n U_i : \mathcal{H}^n \to \oplus_{i = 1} ^n \mathcal{H}_i = \mathcal{K}$$
    unitaren operator. Ker je
    \begin{align*}
        U^* \pi(a E_{ij}) U &= \begin{bmatrix}
            U_1 ^* & & \\
            & \ddots & \\
            & & U_n ^* 
        \end{bmatrix} \pi_{ii} (a) P_{ij} \begin{bmatrix}
            U_1  & & \\
            & \ddots & \\
            & & U_n  
        \end{bmatrix}\\
        &= \begin{bmatrix}
            P_{11}\big|_{\mathcal{H}_1}  & & \\
            & \ddots & \\
            & & P_{1n} \big|_{\mathcal{H}_n} 
        \end{bmatrix} \begin{bmatrix}
           0 & 0 & 0 \\
           0  & 0 & \pi_{ii} (a) {P_{ij}\big|_{\mathcal{H}_j}} \\
            0 & 0 & 0 
        \end{bmatrix}
        \begin{bmatrix}
            P_{11}\big|_{\mathcal{H}_1}  & & \\
            & \ddots & \\
            & & P_{n1} \big|_{\mathcal{H}_1} 
        \end{bmatrix}\\
        &= \begin{bmatrix}
            0 & 0 & 0 \\
            0  & 0 & (P_{1i} \big|_{\mathcal{H}_i}) \pi_{ii} (a){(P_{ij}\big|_{\mathcal{H}_j})}(P_{j1} \big|_{\mathcal{H}_1}) \\
             0 & 0 & 0 
         \end{bmatrix}\\
         &= \begin{bmatrix}
            0 & 0 & 0 \\
            0  & 0 & (P_{1i} \big|_{\mathcal{H}_i}) \pi_{ii} (a)(P_{i1} \big|_{\mathcal{H}_1}) \\
             0 & 0 & 0 
         \end{bmatrix}\\
         &= \begin{bmatrix}
            0 & 0 & 0 \\
            0  & 0 & \pi(a)  \\
             0 & 0 & 0 
         \end{bmatrix},
    \end{align*}
    sledi $U^* \pi([a_{ij}]_{i, j}) U = [\pi(a_{ij})]_{i, j}$.
\end{proof}

\begin{dokaz}[Dokaz izreka \ref{izr:2.4}]
    Brez škode za splošnost naj bo $\| A\| = 1$. Oglejmo si matriko 
    $$B := \begin{bmatrix}
        I & A\\
        A^* & I
    \end{bmatrix} \in M_{2n} (\mathcal{S}).$$
    Ta matrika je pozitivna in dokazali bomo, da je $\| B\| = 2$.

    Za poljubna enotska vektorja $x, y \in \mathcal{K}^n$ velja 
    $$\left\langle B \begin{bmatrix}
        x/\sqrt{2}\\ y/\sqrt{2}
    \end{bmatrix}, \begin{bmatrix}
        x/\sqrt{2}\\ y/\sqrt{2}
    \end{bmatrix} \right\rangle = \frac{1}{2} \| x\|^2 + \real \langle A x, y \rangle + \frac{1}{2} \| y\|^2 = 1 + \real \langle A x, y \rangle.$$
    Če vzamemo supremum po vseh enotskih $x, y \in \mathcal{K}^n$, dobimo na desni strani enačbe $1 + \|A\| = 2$,
    torej $\| B\| \geq 2$. Po drugi strani pa je 
    $$\| B\| \leq \left\| \begin{bmatrix}
        1 & 1\\
        1 & 1
    \end{bmatrix} \right\| = \left\| \begin{bmatrix}
        1 \\ 1
    \end{bmatrix} \begin{bmatrix}
        1 & 1
    \end{bmatrix} \right\| = \left\| \begin{bmatrix}
        1 & 1
    \end{bmatrix} \right\|^2 = 2.$$
    
    Ker je $B$ pozitivna, obstaja neko stanje $\omega$ na $M_{2n} (C^* (\mathcal{S}))$.
    Označimo s $$K = \{\omega \in S(M_{2n} (\mathcal{S}))\ |\ \omega (B) = 2\}$$
    množico vseh takšnih stanj na $M_{2n} (\mathcal{S})$.
    Ta množica je konveksna in kompaktna v OŠ topologiji, zato ima po Krein-Milmanovem izreku 
    ekstremno točko $\omega_0$. Slednja je čisto stanje, torej ga lahko po prejšnjem izreku 
    dilatiramo do maksimalne čiste enotske povsem pozitivne preslikave $\psi: M_{2n} (\mathcal{S}) \to \mathcal{B}(\mathcal{K})$,
    ki ima razširitev na robno upodobitev $\psi': M_{2n}(C^* (\mathcal{S})) \to \mathcal{B}(\mathcal{K})$.
    Po lemi \ref{lem:3.1} obstaja upodobitev $\rho : C^* (\mathcal{S}) \to \mathcal{B}(\mathcal{K}_0)$, tako da je $\psi'$ unitarno ekvivalentna $\rho ^{(2n)}$,
    torej obstaja unitarna preslikava $U: \mathcal{K}_0 ^{2n} \to \mathcal{K}$, tako da je $\psi' = U \rho^{(2n)} U^*$.
    Ker je $\psi'$ nerazcepna, velja enako tudi za $\rho^{(2n)}$ in posledično $\rho$.
    Dokazujemo, da je $\rho$ želena robna preslikava.

    Najprej dokažimo, da ima $\rho \big|_{\mathcal{S}}$ lastnost enolične preslikave.
    Naj bo ${\rho'}: C^*(\mathcal{S}) \to \mathcal{B}(\mathcal{K}_0)$ neka njena enotska povsem pozitivna razširitev.
    Potem je $$U {\rho'} ^{(2n)} U^*: M_{2n} (C^*(\mathcal{S})) \to \mathcal{B}(\mathcal{K})$$ enotska povsem pozitivna razširitev $\psi$, zato je po 
    lastnosti enolične preslikave $$U {\rho'}^{(2n)} U^* = \psi' = U {\rho}^{(2n)} U^*.$$
    Od tod sledi ${\rho'} ^{(2n)} = \rho ^{(2n)}$ in posledično $\rho' = \rho$.
    
    Nazadnje dokažimo še, da je $\| \rho^{(n)} (A) \|= 1$.
    Ker je $\rho$ $*$-homomorfizem, velja $\| \rho^{(n)}\| \leq 1$.
    Ker je $B \in M_{2n} (\mathcal{S})_+$ in $\psi$ enotska povsem pozitivna preslikava,
    je $\| B \| = \| \psi(B) \|$. Od tod pa sledi 
    \begin{align*}
        2 =\| B\| &= \| \psi(B)\| = \| \rho^{(2n)} (B) \|\\
        &= \left\| \begin{bmatrix}
            I & \rho^{(n)} (A)\\
            \rho^{(n)} (A)^* & I
        \end{bmatrix} \right\|\\
        &\leq \left\| \begin{bmatrix}
            1 & \| \rho^{(n)} (A) \|\\
            \| \rho^{(n)} (A) \| & 1
        \end{bmatrix} \right\|\\
        &= 1 + \| \rho^{(n)} A\| \leq 2.
    \end{align*}
    Tako dobimo $\| \rho^{(n)} (A)\| = 1$.
\end{dokaz}

\begin{posledica}
    Naj bo $\mathcal{S} \subseteq C^*(\mathcal{S})$  
    in $\{\rho_\lambda\}_{\lambda \in \Lambda}$ množica predstavnikov vseh unitarnih ekvivalenčnih razredov robnih upodobitev $\mathcal{S}$.
    Če definiramo direktno vsoto $$\pi := \bigoplus_{\lambda \in \Lambda} \rho_\lambda: C^*(\mathcal{S}) \to \bh,$$
    potem je $\pi|_{\mathcal{S}}: \mathcal{S} \to \pi(C^*(\mathcal{S}))$
    $C^*$-ogrinjača za $\mathcal{S}$ in 
    $$\ker \pi = \bigcap_{\lambda \in \Lambda} \ker \rho_\lambda \subseteq C^*(\mathcal{S})$$
    ideal Šilova za $\mathcal{S}$.
\end{posledica}

\begin{dokaz}
    Ker je $\pi$ $*$-homomorfizem, je povsem skrčitev. Po izreku \ref{izr:2.4}
    sledi, da je tudi povsem izometrična preslikava. Nadalje je $\pi\big|_{\mathcal{S}}$
    direktna vsota maksimalnih preslikav, zato je tudi sama maksimalna.
    Iz izreka \ref{izr:2.3} nato sledi, da je $\pi\big|_{\mathcal{S}}: \mathcal{S} \to \pi(C^*(\mathcal{S}))$ 
    $C^*$-ogrinjača za $\mathcal{S}$ in po izreku \ref{izr:aux2} je $\ker \pi$ ideal Šilova za $\mathcal{S}$.
\end{dokaz}

\begin{posledica}[Choquet]
    Če je $\mathcal{S} \subseteq C(X) = C^*(\mathcal{S})$ operatorski sistem, potem je 
    njegov rob Šilova enak zaprtju Choquetovega robu.
\end{posledica}

\begin{dokaz}
    Naj bo $Y \subseteq X$ rob Šilova in $Z \subseteq X$ Choquetov rob za $\mathcal{S}$. Ker so vse robne upodobitve $\mathcal{S}$ 
    oblike $\ev_x : C(X) \to \C$ za $x \in Z$, je po prejšnji posledici 
    \begin{align*}
        I(Y) = \bigcap_{x \in Z} \ker \ev_x = \bigcap_{x \in Z} I(\{x\}) = I(\overline{Z}).
    \end{align*}
    Preostane le še utemeljiti zadnjo enakost. Očitno je $I(\overline{Z}) \subseteq \bigcap_{x \in Z} I(\{x\})$.
    Obratno mora vsaka funkcija $f \in \bigcap_{x \in Z} I(\{x\})$, ki je torej ničelna na $Z$, po zveznosti biti ničelna tudi na zaprtju $\overline{Z}$,
    od koder dobimo $f \in I(\overline{Z})$.
\end{dokaz}

\begin{zgled}
    Vrnimo se k zgledu \ref{zgl:2.2} in poiščimo vse robne upodobitve za operatorski sistem 
    $\mathcal{T}^{(n)} \subseteq M_n (\C)$ do unitarne ekvivalence. Po lemi \ref{lem:3.1} so vse nerazcepne upodobitve
    $\rho: M_n (\C) \to \bh$ unitarno ekvivalentne identiteti $\id: M_n (\C) \to M_n (\C)$.
    Od tod sledi, da je identiteta do unitarne ekvivalence natanko edina robna upodobitev $\mathcal{T}^{(n)}$.
\end{zgled}

\begin{zgled}
    Nadaljujemo zgled \ref{zgl:2.3} in določimo še vse robne upodobitve za operatorski sistem 
    $\mathcal{T}(C(\T)) \subseteq C^*(S)$. Vemo že, da je $\mathcal{K}(H^2)$ njegov ideal Šilova,
    zato za vsako robno upodobitev $\rho: C^*(\mathcal{S}) \to \bh$ velja $\mathcal{K}(H^2) \subseteq \ker \rho$.
    To pomeni, da komutira diagram 
    \begin{equation*}
        \begin{tikzpicture}[>=stealth, node distance=3cm]
    
            \node (S) at (0,2) {$C^*(S)$};
            \node (A) at (4,2) {$\bh$};
            \node (Q) at (0,0) {$\frac{C^*({S})}{\mathcal{K}(H^2)}$,};
            
            % top arrow (unlabeled)
            \draw[->]  (S) -- node[above] {$\rho$} (A);
            
            % left diagonal arrow
            \draw[->] (S) -- node[left] {$q_{\mathcal{K}(H^2)}$} (Q);
            
            % right vertical arrow
            \draw[->] (Q) -- node[below right] {$\overline{\rho} $} (A);
            
        \end{tikzpicture}    
    \end{equation*}
    pri čemer je $\overline{\rho}: \quot{C^*({S})}{\mathcal{K}(H^2)} \to \bh$
    nerazcepna upodobitev za $\quot{C^*({S})}{\mathcal{K}(H^2)}$.
    
    V zgledu \ref{zgl:2.3} smo dokazali, da velja $\quot{C^*({S})}{\mathcal{K}(H^2)} \cong C(\T)$,
    zato je $\mathcal{H} = \C$ in $\overline{\rho} = \ev_\lambda: C(\T) \to \C$ za nek $\lambda \in \T$.
    S tem smo pokazali, da so vse robne upodobitve oblike 
    $$\rho_\lambda: C^*(S) \to \C,\quad \rho_\lambda(T_f + K) = f(\lambda).$$ 
    Dokazati moramo še, da so vse takšne upodobitve $\rho_\lambda$ res robne za $\mathcal{T}(C(\T))$.
    
    Vzemimo poljuben $\lambda \in \T$. Po izreku \ref{izr:2.4} mora obstajati neka robna upodobitev $\rho_\xi$ za $\xi \in \T$, tako da je
    $$\| \rho_\xi (T_{z + \lambda})\| = \| z + \lambda\| = 2.$$
    Ker pa funkcija $z + \lambda \in C(\T)$ zavzame svoj maksimum natanko v $z = \lambda$,
    mora veljati $\xi = \lambda$ in zato je $\rho_\lambda = \rho_\xi$ robna upodobitev.
\end{zgled}







\end{document}
